\chapter{Disturbance growth in thermally excited hypersonic
boundary layers}\label{chapter:disturbance_growth}

In hypersonic boundary layers, transition is compounded by real-gas effects in the post-shock thermal layer. High-enthalpy conditions drive \emph{thermal} nonequilibrium where vibrational modes depart from the translational--rotational temperature in the post shock region and relax on a finite time-scale, altering modal stability characteristics~\cite{Candler2019}. Capturing this behaviour is essential for reliable transition prediction and control. Replicating such conditions while maintaining a quiet freestream is uncommon because of high power demands and stringent acoustic control~\cite{DiRenzo2021}. Nonetheless, modern numerical studies have yielded useful insight into transition behaviour under these conditions. In low-noise flight conditions, transition onset is highly sensitive to small environmental and thermochemical changes, which can move it upstream or downstream. Such low-disturbance environments promote the most natural form of transition; see Chapter~\ref{chapter:introduction} (\textit{Path A} in \fref{figure:transitionmarkovin}). Receptivity first converts external acoustic, entropy, vorticity disturbances and mild surface or thermal non-uniformities into boundary-layer fluctuations. These are organised into discrete linear instability modes, as established in Mack’s seminal analyses~\cite{Mack1984}, which then undergo exponential amplification. Slow streamwise variations in the mean state can introduce \emph{parametric} modulation, altering growth rates and generating sidebands. This pathway ultimately culminates in breakdown to turbulence; eigenmode growth substantially increases aerothermodynamic loads where heat transfer and skin friction rise sharply, with thermal loads reported up to \(30\times\) laminar values~\cite{BitterShepherd2015}. Furthermore, these spectra and growth rates are governed by compressibility, the wall thermal state \(T_w/T_\infty\), and the vibrational relaxation time \(\tau_v\). Thermal nonequilibrium modifies bulk-viscous damping and phase speed, shifts synchronisation, and reshapes the overall stability envelope. Accurately representing the thermal state and operating conditions is of utmost importance for predicting the transition location.





% In hypersonic boundary layers, transition is compounded by real-gas effects in the post-shock thermal layer. High-enthalpy conditions drive \emph{thermal} nonequilibrium: vibrational modes depart from the translational/rotational temperature and relax on a finite time-scale \(\tau_v\), altering modal stability characteristics~\cite{Candler2019}. Replicating such conditions while maintaining a quiet freestream is uncommon because of high power demands and stringent acoustic control~\cite{DiRenzo2021}. In low-disturbance, flight-like environments, vanishingly small unstable perturbations grow exponentially (Mack first/second modes). Spectra and growth rates are governed by compressibility, wall thermal state \(T_w/T_\infty\), and vibrational relaxation \(\tau_v\). Thermal nonequilibrium modifies bulk-viscous damping and phase speed, shifts synchronisation, and reshapes the second-mode envelope. Receptivity to acoustic, entropy, and vorticity forcing sets initial amplitudes; downstream, synchronisation and thermo--viscous coupling control amplification, precipitating secondary instability (fundamental or subharmonic oblique) and breakdown. The aerothermodynamic penalties are significant: skin friction rises and wall heat flux can increase by an order of magnitude, with thermal loads reported near \(30\times\)~\cite{Bitter2015}. 

% In hypersonic boundary layers, transition is compounded by real-gas effects in the post-shock thermal layer. High-enthalpy conditions drive \emph{thermal} nonequilibrium where vibrational modes depart from the translational/rotational temperature in the post shock region and relax on a finite time-scale \(\tau_v\), altering modal stability characteristics~\cite{Candler2019}. Capturing this behaviour is essential for reliable transition prediction and control. Replicating such conditions while maintaining a quiet freestream is uncommon because of high power demands and stringent acoustic control~\cite{DiRenzo2021}. Nonetheless, modern numerical studies have yielded useful insight into transition behaviour under these conditions. In low-noise flight conditions, transition onset is highly sensitive to small environmental and thermochemical changes, which can move it upstream or downstream. 
% This eigenmode growth disrupts the flow and leads to turbulence, substantially increasing aerothermodynamic loads: heat transfer and skin friction rise sharply, with thermal loads reported up to \(30\times\) laminar values~\cite{Bitter2015}. Furthermore, spectra and growth rates are governed by compressibility, the wall thermal state \(T_w/T_\infty\), and the vibrational relaxation time \(\tau_v\). Thermal nonequilibrium modifies bulk-viscous damping and phase speed, shifts synchronisation, and reshapes the overall stability envelope. Accurately representing the thermal state and operating conditions is of utmost importance for predicting the transition location.


\begin{figure}[hbt!]
    \centering
    \fontsize{8}{9}\selectfont\includesvg[width=0.95\textwidth]{resources/figures/disturbance/schematic_wedgeflow.svg}
    \captionof{figure}{Schematic of hypersonic flow over a two-dimensional wedge, highlighting vibrational nonequilibrium phenomena. Figure adapted from~\cite{williams2021}. }\label{figure:wedgeflow}
\end{figure}

% In what follows, we focus on the role of thermal nonequilibrium in hypersonic boundary layers. Figure~\ref{figure:wedgeflow}, adapted from \citet{williams2021}, schematically delineates regions of thermal equilibrium (green) and thermally frozen flow (blue). The upper panel depicts flow over an extended wedge, where the boundary layer (bounded by the dashed line) gradually relaxes downstream towards equilibrium. The lower panel magnifies the near-wall region and includes the immediate post-shock flow, which is treated as equilibrated at the onset of boundary-layer development near the leading edge. Owing to the nose geometry, a favourable pressure gradient arises, and additional pressure-gradient effects may be induced by thermal nonequilibrium through the thermal response to local shock curvature. Elevated temperatures also develop within the boundary layer due to viscous work and its dissipation into heat. In the present study, we demonstrate that vibrational energy relaxes towards equilibrium in the very region where boundary-layer instabilities are amplifying, establishing a direct spatial overlap between nonequilibrium relaxation and disturbance growth.

% Consider a \gls{HGV} shaped as a wedge flying at hypersonic speed at low altitude. As sketched in Figure~\ref{figure:wedgeflow} (adapted from \citet{williams2021}), the flow field partitions into zones of \gls{TEQ} (green) and \gls{TNE} (blue). Immediately behind the oblique shock near the leading edge, the post-shock gas is close to \gls{TEQ}, while within the boundary layer the vibrational temperature lags the translational–rotational temperature, producing \gls{TNE} that relaxes downstream over a finite time scale. The extent of each zone depends on ambient density and heating: at low altitude (higher density), relaxation is faster and \gls{TEQ} occupies a larger fraction of the near-field; farther downstream, viscous dissipation raises temperatures and \gls{TNE} can persist until vibrational modes re-equilibrate. Local shock curvature and the favourable pressure gradient induced by the nose further modulate these \gls{TEQ}/\gls{TNE} regions along the surface.

Consider a \gls{HGV} shaped as a wedge flying at hypersonic speed at low altitude. As sketched in Figure~\ref{figure:wedgeflow} (adapted from \citet{williams2021}), the flow field partitions into zones of \gls{TEQ} (green) and \gls{TNE} (blue). Immediately behind the oblique shock near the leading edge, the post-shock gas is close to \gls{TEQ}, while within the boundary layer the vibrational temperature lags the translational–rotational temperature, producing \gls{TNE} that relaxes downstream over a finite time scale. The extent of each zone depends on ambient density and heating: at low altitude (higher density), relaxation is faster and \gls{TEQ} occupies a larger fraction of the near-field; farther downstream, viscous dissipation raises temperatures and \gls{TNE} can persist until vibrational modes re-equilibrate. Local shock curvature and the favourable pressure gradient induced by the nose further modulate these \gls{TEQ}/\gls{TNE} regions along the surface. Instability modes respond differently across these zones: growth rates, phase speeds, and synchronisation are modified in \gls{TNE} owing to vibrational relaxation and altered bulk-viscous damping, whereas \gls{TEQ} exhibits the classical compressible behaviour. For clarity in what follows, we therefore consider a nominally zero streamwise pressure gradient to isolate thermally driven effects on the spectra and growth of instability modes along the domain.


% This chapter quantifies linear-mode growth in thermally excited hypersonic boundary layers across conditions parameterised by \(M_\infty\), \(T_w/T_\infty\), and \(\tau_v\), by co-analysing \gls{DNS} and \gls{LST} on matched self-similar laminar bases with harmonised scalings, spectral definitions, and normalisations, yielding a reproducible protocol for cross-comparison and interpretation.

% This chapter provides a comprehensive quantification of linear-mode growth in thermally excited hypersonic boundary layers across flight-relevant conditions parameterised by \(M_\infty\), \(T_w/T_\infty\), and \(\tau_v\). We co-analyse \gls{DNS} and \gls{LST} on matched self-similar laminar bases with identical gas models, transport closures, and wall thermal boundary conditions, and enforce harmonised scalings (semilocal units), spectral definitions (frequency–spanwise wavenumber pairs), and energy-based normalisations to ensure like-for-like comparison. The protocol standardises base-flow construction and verification, receptivity-to-forcing conventions, and the extraction of modal diagnostics—streamwise growth rates \(\alpha_i(x)\), phase speed \(c_p\), synchronisation locations, and cumulative \(N\)-factors—together with uncertainty controls (DNS grid/time-step independence and LST eigenvalue tolerances). This framework yields a reproducible stability envelope and a consistent interpretation of how \(M_\infty\), \(T_w/T_\infty\), and \(\tau_v\) regulate amplification, thereby linking numerical predictions to flight-condition trends.

This chapter develops a unified assessment of linear-mode growth in thermally excited hypersonic boundary layers over flight-relevant ranges of \(M_\infty\), wall-to-freestream temperature \(T_w/T_\infty\), and vibrational relaxation time \(\tau_v\). Emphasis is placed on the mechanisms by which thermal nonequilibrium modifies receptivity to external disturbances, phase speed and synchronisation, and the balance between bulk-viscous damping and acoustic trapping, culminating in changes to streamwise growth rates \(\alpha_i(x)\) and cumulative \(N\)-factors. The resulting stability envelope, mapped consistently across \gls{TEQ} and \gls{TNE} bases, establishes how these parameters regulate amplification and thereby influence the predicted transition location. This foundation is used in later chapters to interpret early nonlinear evolution and to relate numerical predictions to trends expected under representative flight conditions.










% In hypersonic boundary layers, transition is compounded by real-gas effects in the post-shock thermal layer. High-enthalpy conditions activate thermochemical nonequilibrium—most notably vibrational excitation and (where relevant) finite-rate chemistry—which introduces additional relaxation time-scales and modifies stability characteristics~\cite{Candler2019}. Accurately reproducing this physics remains difficult: only a subset of facilities can sustain the required enthalpy while preserving “quiet” freestreams for transition studies, owing to high power demands~\cite{DiRenzo2021} and stringent acoustic/vibrational controls. As a result, the literature is disproportionately weighted towards initial amplification, and a unified picture of hypersonic boundary-layer transition remains incomplete. The present study combines \gls{DNS} with \gls{LST} based on the presented self-similar laminar base flows to isolate nonequilibrium effects and clarify the mechanisms precipitating turbulence.

% In hypersonic boundary layers, transition is compounded by real-gas effects in the post-shock thermal layer. High-enthalpy conditions drive \emph{thermal} nonequilibrium where vibrational modes depart from the translational/rotational temperature in the post-shock region and relax on a finite time-scale \(\tau_v\), altering modal stability characteristics~\cite{Candler2019}. Capturing this behaviour is essential for reliable transition prediction and control. Replicating such conditions while maintaining a quiet freestream is uncommon because of high power demands and stringent acoustic control~\cite{DiRenzo2021}. Nonetheless, modern numerical studies have yielded useful insight into transition behaviour under these conditions. In low-noise flight conditions, transition onset is highly sensitive to small environmental and thermochemical changes, which can move it upstream or downstream. Such low-disturbance environments promote the most natural transition pathway; further discussion is provided in \ref{chapter:introduction} (\textbf{Path A} in \fref{figure:transitionmarkovin}). This eigenmode growth disrupts the flow and leads to turbulence, substantially increasing aerothermodynamic loads: heat transfer and skin friction rise sharply, with thermal loads reported up to \(30\times\) laminar values~\cite{Bitter2015}. Furthermore, spectra and growth rates are governed by compressibility, the wall thermal state \(T_w/T_\infty\), and the vibrational relaxation time \(\tau_v\). Thermal nonequilibrium modifies bulk-viscous damping and phase speed, shifts synchronisation, and reshapes the overall stability envelope. Accurately representing the thermal state and operating conditions is of utmost importance for predicting the transition location.

% Low-disturbance, flight-like environments permit extremely small unstable perturbations to undergo exponential modal growth (e.g. Mack first/second modes), with eigenvalue spectra and growth rates strongly influenced by compressibility and relaxation processes. Receptivity sets the initial amplitudes; downstream, synchronisation and thermo-viscous/thermochemical coupling govern amplification, culminating in secondary instabilities and breakdown. The aerothermodynamic penalties are substantial: skin-friction drag rises and wall heat flux can increase by an order of magnitude or more, with reports of thermal loads approaching a 30× amplification~\cite{Bitter2015}. Quantifying the initiation, growth, and nonlinear breakdown of these disturbances is therefore central to reliable design. A detailed account of the breakdown pathway is provided in \ref{chapter:introduction} (\textbf{Path A} in \fref{figure:transitionmarkovin}).

% Low-disturbance, flight-like environments allow vanishingly small unstable perturbations to grow exponentially (Mack first/second modes), with spectra and growth rates governed by compressibility, wall thermal state \(T_w/T_\infty\), and finite-rate vibrational relaxation \(\tau_v\). Thermal nonequilibrium modifies effective bulk-viscous damping and phase speed, shifting the synchronisation location and reshaping the second-mode envelope. Receptivity to acoustic, entropy, and vorticity forcing sets initial amplitudes; downstream, branch coalescence (synchronisation) and thermo--viscous coupling—mediated by temperature-dependent transport and vibrational energy exchange—control amplification, leading to secondary instability (fundamental or subharmonic oblique routes) and breakdown. Aerothermodynamic penalties are severe: skin friction increases and wall heat flux can rise by an order of magnitude, with thermal loads reported near \(30\times\)~\cite{Bitter2015}. Precise quantification of receptivity, linear amplification, secondary instability, and nonlinear breakdown under \emph{thermal} nonequilibrium is therefore essential for reliable design. A detailed breakdown pathway is given in \ref{chapter:introduction} (\textbf{Path A} in \fref{figure:transitionmarkovin}).


% The aim of this chapter is to quantify the growth of linear instability modes in thermally excited hypersonic boundary layers across a range of flight conditions (characterised by \(M_\infty\), \(T_w/T_\infty\), and \(\tau_v\)), by systematically comparing \gls{DNS} with \gls{LST} on self-similar laminar bases and establishing clear pathways for consistent cross-comparison and interpretation across modelling frameworks and datasets.


% \section{Laminar turbulent transitions}
% \section{Transitional flows}
% The transition from a laminar to a turbulent boundary layer is a critical aerodynamic phenomenon with significant implications for skin friction, heat transfer, and aerodynamic drag. The destabilisation of an otherwise laminar flow into a chaotic, turbulent regime had attracted interest since the late 1800s when Osborne Reynolds first demonstrated that laminar flow could become unstable under certain conditions. Over the decades, extensive research has sought to understand, predict, and control this process through theoretical and experimental efforts. Within this broader context, \citet{Morkovin1994} laid out a conceptual foundation that synthesised these developments into the modern view of the multiple paths to transition that can occur in boundary layer flows. \fref{figure:transitionmarkovin}, reproduced from \citet{Fedorov2011}, presents a schematic overview of these paths to transition, which applies to both low-speed and high-speed boundary layers, although the specific instability modes involved differ between the two regimes.


% \begin{figure}[t!]
%     \centering
%     \fontsize{9}{10}\selectfont\includesvg[width=0.50\textwidth]{resources/figures/introduction/schematic_fedorov.svg}
%     \caption{\cite{Morkovin1994} roadmap to transition, diagram of pathways by which a boundary layer might transition to turbulence \citep{Fedorov2011}.}
%     \label{figure:transitionmarkovin}
% \end{figure}


% Transition begins with the appearance and amplification of disturbances from environmental sources such as freestream turbulence, acoustic noise, surface roughness, or shock interactions. These interact with the boundary layer, generating oscillatory modes whose amplitude, frequency, and phase define the initial conditions for instability waves that grow and ultimately breakdown. This \textit{receptivity process} governs how external disturbances are converted into instability waves, strongly influencing when and where transition occurs \citep{Morkovin1994}. In hypersonic boundary layers, key destabilising factors include freestream perturbations (acoustic, entropy, and vorticity fluctuations), surface roughness that introduces localised bypass routes, shock interactions that create unsteady pressure and velocity fields, and high-temperature effects like vibrational excitation and dissociation, all of which can significantly enhance instability growth \citep{vandereynde2015,Bitter2015}. Because these environmental disturbances are inherently complex and broadband, the receptivity problem remains largely unresolved for compressible boundary layers, despite extensive research \citep{dCerminara2017} Ma Zhong.

% In the most general case, low-amplitude oscillatory modes are triggered within the boundary layer and amplify through linear modal growth, followed by parametric instabilities and mode interactions that transfer energy between modes and lead to breakdown to turbulence; this transition pathway is depicted in \fref{figure:transitionmarkovin} as \textit{path A}. This path has been widely studied since the pioneering work of Tollmien and Schlichting, which established the well-known two-dimensional Tollmien--Schlichting instability waves. For higher disturbance levels, the superposition of perturbations may first exhibit algebraic transient growth before feeding into modal amplification (\textit{path B}). As the disturbance level increases further, transient growth can lead directly to parametric instabilities and mode interactions or, if sufficiently strong, bypass these stages entirely and transition straight to breakdown (\textit{paths C} and \textit{D}). Only \textit{paths A, B,} and \textit{C} are relevant for external hypersonic boundary layers, where disturbances grow through receptivity, transient growth, and modal amplification. \textit{Path D} mainly applies to internal flows with elevated turbulence levels, while \textit{path E} represents extreme forcing such as severe surface roughness or intense freestream turbulence where linear growth mechanisms are completely bypassed.


% Most of the groundwork done for  

% Ok, hear me out, most of the groundwork done for transition prediction was on incompressible flows which is not completely true but we say that anyway. The first studies about boundary layer stability was performed my Lees and Lin, who were quantifying the inflection point phenomenon, then came the Tollmein-Schliting waves, confirmed by further experiments by Schuberger or whatever their names are. It was not until Mack studies the stability of supersonic boudnary layer that we have a rigiorous definition of boundary layers and what they actually are. Mack performed stability analysis on supersonic base flows and characterised instability modes into numbers (first, second and third mode, so on). A detailed description of these instabilities is discussed in chapter 4, but the general gist of it is, first modes are promiment in subsonic incompressible flows, often associated with lower frequencies and the second mode is an acoustic trapped wave that is a feature of flows past supersonic regime. 


% The transition begins with the appearance and subsequent amplification of disturbances, which can arise from environmental sources such as freestream turbulence, acoustic noise, surface roughness, or shock interactions, depending on the flow regime. These environmental disturbances interact with the boundary layer, appearing as fluctuations in the base flow that generate oscillatory modes whose amplitude, frequency, and phase set the initial conditions for the instability waves that develop and ultimately lead to breakdown. This process, known as the receptivity mechanism, converts external disturbances into instability waves; their initial amplitude plays a key role in determining how and when the transition occurs. Because environmental disturbances are often complex and contain a wide range of information, the receptivity problem remains largely unsolved for compressible boundary layers, although several studies have tried to tackle the issue at hand. 

% Transition begins with the appearance and subsequent amplification of disturbances originating from environmental sources such as freestream turbulence, acoustic noise, surface roughness, or shock interactions, depending on the flow regime. These external disturbances interact with the boundary layer, manifesting as fluctuations in the base flow that generate oscillatory modes whose amplitude, frequency, and phase define the initial conditions for instability waves that develop and eventually breakdown. This process, known as receptivity, converts external disturbances into instability waves whose initial amplitude plays a critical role in determining the onset and location of transition. In hypersonic boundary layers, key destabilising factors influencing transition in hypersonic flows include: \textbf{Freestream disturbances}: Acoustic waves, entropy fluctuations, and vorticity perturbations in the freestream can penetrate the boundary layer and trigger instability growth. \textbf{Surface roughness}: Imperfections on the surface can introduce localised perturbations that bypass classical instability mechanisms, leading to premature transition. \textbf{Shock interactions}: Shock-induced perturbations can introduce unsteady pressure and velocity fluctuations, modifying the growth of boundary layer instabilities. \textbf{High-temperature effects}: Vibrational excitation, dissociation, and recombination reactions in high-enthalpy flows alter the stability properties of the boundary layer, often enhancing instability growth. Because environmental disturbances are typically complex and broadband, the receptivity problem remains largely unsolved for compressible boundary layers, despite numerous efforts to address it.



% Transition begins with the appearance and amplification of disturbances from environmental sources such as freestream turbulence, acoustic noise, surface roughness, or shock interactions. These interact with the boundary layer, generating oscillatory modes whose amplitude, frequency, and phase define the initial conditions for instability waves that grow and ultimately breakdown. This \textit{receptivity process} governs how external disturbances are converted into instability waves, strongly influencing when and where transition occurs \citep{Morkovin1994}. In hypersonic boundary layers, key destabilising factors include freestream perturbations (acoustic, entropy, and vorticity fluctuations), surface roughness that introduces localised bypass routes, shock interactions that create unsteady pressure and velocity fields, and high-temperature effects like vibrational excitation and dissociation, all of which can significantly enhance instability growth \citep{vandereynde2015,Bitter2015}. Because these environmental disturbances are inherently complex and broadband, the receptivity problem remains largely unresolved for compressible boundary layers, despite extensive research \citep{dCerminara2017} Ma Zhong.

% In the most general case, low-amplitude oscillatory modes are triggered within the boundary layer and amplify through linear modal growth, followed by parametric instabilities and mode interactions that transfer energy between modes and lead to breakdown to turbulence; this transition pathway is depicted in \fref{figure:transitionmarkovin} as \textit{path A}. This path has been widely studied since the pioneering work of Tollmien and Schlichting, which established the well-known two-dimensional Tollmien--Schlichting instability waves. For higher disturbance levels, the superposition of perturbations may first exhibit algebraic transient growth before feeding into modal amplification (\textit{path B}). As the disturbance level increases further, transient growth can lead directly to parametric instabilities and mode interactions or, if sufficiently strong, bypass these stages entirely and transition straight to breakdown (\textit{paths C} and \textit{D}). Only \textit{paths A, B,} and \textit{C} are relevant for external hypersonic boundary layers, where disturbances grow through receptivity, transient growth, and modal amplification. \textit{Path D} mainly applies to internal flows with elevated turbulence levels, while \textit{path E} represents extreme forcing such as severe surface roughness or intense freestream turbulence where linear growth mechanisms are completely bypassed.


% Most of the groundwork done for  

% Ok, hear me out, most of the groundwork done for transition prediction was on incompressible flows which is not completely true but we say that anyway. The first studies about boundary layer stability was performed my Lees and Lin, who were quantifying the inflection point phenomenon, then came the Tollmein-Schliting waves, confirmed by further experiments by Schuberger or whatever their names are. It was not until Mack studies the stability of supersonic boudnary layer that we have a rigiorous definition of boundary layers and what they actually are. Mack performed stability analysis on supersonic base flows and characterised instability modes into numbers (first, second and third mode, so on). A detailed description of these instabilities is discussed in chapter 4, but the general gist of it is, first modes are promiment in subsonic incompressible flows, often associated with lower frequencies and the second mode is an acoustic trapped wave that is a feature of flows past supersonic regime. 



% The primary focus of the current thesis is to understand the modal decomposition of \textit{Path A}, 

% Transition to turbulence in hypersonic boundary layers has been studied in the past, both in \textit{`cold flow'} and high-enthalpy regime at suborbital enthalpies, despite the challenges imposed by the later environment. The studies concern stages of transition through perturbations in the laminar base flow through linear and non-linear stability analyses. \cite{Cerminara2019, cerminarasandham2020} studied the receptivity mechanisms at the leading edge of a blunt body in cold hypersonic flow without considering finite-rate chemical models. Studies of transition and control for chemically perfect gases, and flows in chemical equilibrium were first performed by  \citep{Malik1989, Malik1990b, Malik1990a, Malik1990c}. This was extended further by \cite{Stuckert1994} and \cite{Hudson1997}, who introduced chemical and vibrational non-equilibrium extensions with Park's two-temperature model. The studies conclusively show that high enthalpy effects bring non-negligible contribution to boundary layer stability; the endothermic nature of the reactions remove energy from the thermal boundary layer, leaving the flow more subjected to second-mode (Mack modes) disturbances, almost mimicking wall-cooling effects. \cite{Malik2003} and \cite{Johnson2005}, conducted parabolised stability equations (PSE) to the subject, with the former study concluding that some dissociation reactions actually reduce the growth of perturbations caused by the chemical reactions on the base flow. \cite{Johnson2005} has also performed simulations using multi-species models to better capture the stability during transition. \cite{Franko2010} concluded that the chemical and viscosity model used in the analysis has a substantial effect on the stability characteristic. Consequently, \cite{Klentzman2013} showed that viscosity stabilises the base flow while investigating viscous effects on the receptivity of boundary layers and \cite{miromiro2018} extended the work on an adiabatic flat plate, observing the flow to most destabilise when mass diffusion is hundreds of times faster than viscous diffusion and further investigated the effects of choice of thermochemical and thermophysical framework on the transition prediction \citep{MiroMiro2019}. Furthermore, work of \cite{WartemannEtal2019} show that thermochemical non-equilibrium approach for the mean flow and the stability code are almost identical to the approach of considering just the base flow in thermochemical non-equilibrium. Most of the work in the literature focused on the modal growth of the transition process predominantly for sub-orbital enthalpies. The stability of high-speed flows at low altitudes, where the Reynolds numbers are significantly higher and a comprehensive study of the evolution of the boundary layer to complete turbulence transition in such flows has been largely unexplored areas.  


\section{Second mode propagation}
The transition from laminar to turbulent flow on a vehicle surface occurs due to the emergence and subsequent amplification of perturbations in the flow. These perturbations can originate from various sources, including free-stream disturbances, acoustic waves, shock interactions, structural vibrations, surface roughness, surface catalysis or internal fluctuations. Typically there is a receptivity phase, during which the flow field responds to the perturbations. For small amplitude perturbations, this sets the conditions for an initially linear growth followed by nonlinear breakdown to turbulence, as described in~\cite{reshotko2008}. The initial modal growth of the unstable boundary-layer waves can be obtained by solving the compressible Orr-Sommerfeld equation. The most unstable eigenmodes computed by solving the eigenvalue problem are Tollmien–Schlichting waves in low-speed flow, whereas for high-speed boundary layer flows multiple unstable modes exist. Early studies of~\cite{Mack1969} provided a wider understanding of these modes, classifying them into two distinct types of linear instabilities known as Mack's \textit{first} and \textit{second modes} which is also the lowest of the acoustic modes. The first mode is primarily vortical in nature. In supersonic boundary layers, the first-mode instability amplifies most readily at oblique orientations, nonzero spanwise wavenumber. Although its local growth rate can be small, the long streamwise development region allows substantial energy accumulation. Several mechanisms damp the first mode in high-speed flows notably wall cooling which strongly stabilises first-mode waves, favourable pressure gradients, and enhanced thermal-viscous or bulk-viscous dissipation~\citep{Mack1984}.
The second mode is characterised by trapped acoustic waves between the wall and the sonic line relative to the relative perturbation propagation speed~\cite{Fedorov2001} as depicted in ~\fref{figure:schematicmackmode}. The second-mode instability exhibits growth rates typically much larger than those of the first mode, but over a shorter streamwise interval. It is predominantly two-dimensional, spanwise-uniform, high-frequency and acoustic in character, with phase speeds near the local sound speed and energy trapped between the wall and the sonic line~\citep{DeTullio2013}. Wall cooling destabilises this mode, making it the principal instability mechanism in many hypersonic boundary layers. Its amplification is sensitive to the mean thermal state, through sound-speed and viscous or bulk-viscous effects, and to slight changes in synchronisation conditions, which confine its growth to a relatively narrow streamwise band.


\begin{figure}[hbt!]
    \centering
    \fontsize{9}{10}\selectfont\includesvg[width=0.9\textwidth]{resources/figures/disturbance/schematic-mackmode.svg}
    \captionof{figure}{Schematic of acoustic second mode instability in a high-speed boundary layer, where U(y) is the mean streamwise flow profile. Adapted from \citep{Fedorov2011}}.
    \label{figure:schematicmackmode}
\end{figure}

% The early stages of transition can be accurately captured by \gls{LST}, provided the disturbance environment has low amplitude and the Reynolds number is high enough for the boundary layer to be only weakly non-parallel. Early results for \gls{LST} with real-gas effects began with studies by \cite{Malik1989, Malik1990a} for local thermochemical equilibrium effects compared to a \gls{CPG}, followed by ~\cite{stuckertreed1994}, who studied chemical nonequilibrium.~\cite{Hudson1997} advanced this work by incorporating both chemical and vibrational nonequilibrium using Park's two-temperature model. They found that endothermic reactions, which extract energy from the boundary layer, suppress first-mode instabilities while enhancing second-mode instability, similar to wall-cooling effects. It was found that flows in thermochemical nonequilibrium exhibit higher growth rates than \gls{CPG} flows, but lower growth rates than thermal and chemical equilibrium flows. This suggests that a combination of thermochemical nonequilibrium has a stabilising effect on the flow. These results were corroborated by \gls{LST} and experimental studies from the T5 shock tunnel by ~\cite{JohnsonEtal1998}. Their study of high-enthalpy flows past a cone revealed that endothermic reactions stabilise the flow, while exothermic reactions have the opposite effect. ~\cite{JohnsonCandler2005} demonstrated that the predicted transition location can shift forward by nearly 50\% when finite-rate reactions are modelled.~\cite{MortensenZhong2013} incorporated multi-species chemical effects and extensively compared \gls{LST} to \gls{DNS} for various flow conditions past a cone, focusing on surface ablation and roughness and noticing that change in transport properties significantly affects the growth rate predictions. Recently,~\cite{Bitter2015} investigated the influence of vibrational nonequilibrium on stability, concluding that this can play an important role.

The early stages of transition can be captured accurately by \gls{LST}, provided the disturbance environment has low amplitude and the Reynolds number is high enough that the boundary layer is only weakly non-parallel. Early \gls{LST} results including real-gas effects began with \cite{Malik1989, Malik1990a}, who examined local thermochemical equilibrium relative to a \gls{CPG}, followed by \cite{stuckertreed1994}, who studied chemical nonequilibrium. \cite{Hudson1997} advanced this line of work by incorporating both chemical and vibrational nonequilibrium using Park’s two-temperature model. They found that endothermic reactions, which extract energy from the boundary layer, suppress first-mode instabilities while enhancing the second mode analogous to wall cooling. Overall, thermochemical nonequilibrium produced higher growth rates than \gls{CPG} flows but lower than thermal and chemical equilibrium, indicating an overall stabilising influence. These trends were corroborated by \gls{LST} and experiments in the T5 shock tunnel by \cite{JohnsonEtal1998}, who showed that endothermic reactions stabilise the flow, whereas exothermic reactions destabilise it. \cite{JohnsonCandler2005} further demonstrated that the predicted transition location can shift upstream by nearly 50\% when finite-rate reactions are modelled. \cite{MortensenZhong2013} incorporated multi-species chemical effects and compared \gls{LST} with \gls{DNS} for various cone flows, including surface ablation and roughness, noting that changes in transport properties significantly affect growth-rate predictions. Most recently, \cite{BitterShepherd2015} investigated vibrational nonequilibrium and concluded that it can play an important role in stability. 

% Previous \gls{LST} studies, such as \cite{MiroMiro2019}, have shown that local thermochemical equilibrium can significantly amplify the $N$-factors, potentially altering the transition breakdown process. However, it remains unclear whether this shift arises primarily from chemical equilibrium effects or from thermal equilibrium. Chemical nonequilibrium effects on modal growth are analysed in \cite{FrankoMacCormackLele2010}, with \cite{LyttleReed2005} offering the only partial consideration of transport properties. While both works acknowledge that discrepancies stem chiefly from the thermal boundary layer, they stop short of isolating transport-model influences on the thermal state or quantifying their magnitude. The goal of the current chapter is to assess the effects of Mach number and wall temperature on thermally excited boundary layers, and to examine the progression of $N$-factors through a combination of \gls{LST} and \gls{DNS} in typical flight scenarios.

Previous \gls{LST} studies have demonstrated that local thermochemical equilibrium can significantly amplify the $N$-factors, potentially altering the transition breakdown process \citep{MiroMiro2019}. However, it remains unclear whether this shift arises primarily from chemical equilibrium effects or from thermal equilibrium. Chemical nonequilibrium effects on modal growth are examined in the literature \citep{FrankoMacCormackLele2010}, with only partial consideration given to transport properties and their influence on the thermal state \citep{LyttleReed2005}. Whilst these studies acknowledge that discrepancies stem chiefly from the thermal boundary layer, the magnitude of transport-model influences remains inadequately quantified. The goal of the current chapter is to assess the effects of Mach number and wall temperature on thermally excited boundary layers, and to examine the progression of $N$-factors through a combination of \gls{LST} and \gls{DNS} in typical flight scenarios.


% \section{Linear stability framework}


\section{Linear stability framework}
% The propagation of small-amplitude disturbances can be accurately captured by a normal-mode ansatz that reduces the dynamics to growth rates and phases across the boundary layer. The linear stability framework examines the 
% In what follows, small-amplitude disturbances are analysed about a known base state to determine where and how they amplify. The local, parallel viewpoint provides growth rates, phase speeds, and wall-normal structures that delineate neutral curves and most-amplified frequency bands, from which $N$-factors are accumulated to estimate transition onset. The approach is especially useful in compressible boundary layers, where Mach number, wall temperature, and transport-property variation shape amplification; it also guides DNS/LES set-ups by informing frequency selection, domain placement, and sensor/forcing locations. While restricted to weak non-parallelism and the early, linear stages of evolution, it offers a practical bridge from a computed base flow to quantitative predictions of initial amplification and the streamwise window over which breakdown is expected~\citep{SchmidHenningson2001}.



The linear propagation and subsequent decay of small, superposed disturbances in a given base state can be accurately captured by stability analysis. A suite of frameworks, ranging from local normal mode to parabolised and fully global analyses, rests on different treatments of parallelism and nonparallel effects and, taken together, yields workable estimates of the transition stage, including neutral boundaries, peak amplification bands, \textit{N}-factor levels, and probable breakdown windows; a detailed review of these methods can be found in \citep{Theofilis2011}. Focusing on local disturbances in a nearly parallel boundary layer, the problem is cast as a one-dimensional eigenvalue problem in the wall-normal coordinate. The resulting spectrum provides amplification rates, phase speeds, and wall-normal eigenfunctions, from which neutral curves and bands of peak gain are assembled and \textit{N}-factors accumulated to estimate the onset of transition. Two modelling choices make this tractable. Linearity assumes disturbances are sufficiently small that quadratic interactions can be neglected; although breakdown is nonlinear, the linear growth stage is typically the slowest and thus sets the effective transition Reynolds number, and the most-amplified frequency in this regime tends to persist into the nonlinear range \citep{Reshotko2001,Mack1984}. Local parallelism assumes that, over a short streamwise interval, the base flow varies only with the wall-normal coordinate; this has proved adequate for charting initial amplification and is supported by classic experiments that identified instability waves and their quantitative trends.

% Within these bounds, the framework connects a computed mean flow to usable predictions—guiding DNS/LES set-ups and clarifying how Mach number, wall temperature, and transport-property variations influence disturbance growth \citep{SchmidHenningson2001}.
% \begin{equation}
%   \boldsymbol{q}(x_j,t) \;=\; \bar{\boldsymbol{q}}(x_j,t) \;+\; \tilde{\boldsymbol{q}}(x_j,t),
% \end{equation} 

The first step is to introduce the disturbance field explicitly and linearise about the base state to obtain the eigenproblem for growth rates and phases. The instantaneous state is written as a mean plus a small fluctuation,
\begin{equation}
  \boldsymbol{q} \;=\; \bar{\boldsymbol{q}}\;+\; \tilde{\boldsymbol{q}}(x_j,t),
\end{equation} 
where \(\bar{\boldsymbol{q}}=(\bar{\rho},\,\bar{u},\,\bar{v},\,\bar{w},\,\bar{T})^{\mathsf{T}}\) denotes the base flow and \(\tilde{\boldsymbol{q}}=(\tilde{\rho},\,\tilde{u},\,\tilde{v},\tilde{w},\,\tilde{T})^{\mathsf{T}}\) the small disturbance. The boundary layer is deemed unstable where the disturbance amplitude increases downstream (or in time), and stable where it decays. The governing equations are linearised about \(\bar{\boldsymbol{q}}\); under a locally parallel assumption the coefficients depend only on \(y\).
Accordingly, a normal–mode ansatz is adopted for the disturbance field,
\begin{equation}
  \tilde{\boldsymbol{q}}(x,y,z,t)
  = \hat{\boldsymbol{q}}(y)\,\exp\!\big(i(\alpha x+\beta z-\omega t)\big),
\end{equation}
where \(\hat{\boldsymbol{q}}(y)\) is the complex eigenfunction, \(\alpha,\beta\) are the streamwise/spanwise wavenumbers, and \(\omega\) the frequency.
Substitution of this ansatz into the linearised system yields a linear eigenvalue problem for \(\omega\) (temporal) or \(\alpha\) (spatial).



The governing equations are linearised about \(\bar{\boldsymbol{q}}\); under a locally parallel assumption the coefficients depend only on \(y\).
Accordingly, a normal–mode ansatz is adopted for the disturbance field,
\begin{equation}
  \tilde{\boldsymbol{q}}(x,y,z,t)
  = \hat{\boldsymbol{q}}(y)\,\exp\!\big(i(\alpha x+\beta z-\omega t)\big),
\end{equation}
where \(\hat{\boldsymbol{q}}(y)\) is the complex eigenfunction, \(\alpha,\beta\) are the streamwise/spanwise wavenumbers, and \(\omega\) the frequency.
Substitution of this ansatz into the linearised system yields a linear eigenvalue problem for \(\omega\) (temporal) or \(\alpha\) (spatial).

Substitution of the normal-mode ansatz into the linearised Navier--Stokes equations yields a linear eigenvalue problem,
\begin{equation*}
  L\,\hat{\boldsymbol{q}} = \omega\,\hat{\boldsymbol{q}}.
\end{equation*}
The interpretation hinges on the spectral parameters \((\alpha,\beta,\omega)\). In what follows, a spatial analysis is adopted: disturbances grow in space with \(\alpha \in \mathbb{C}\) and \(\omega \in \mathbb{R}\), and \(\alpha\) is iterated until \(\omega\) is real. In this setting, the spatial growth rate is \(-\operatorname{Im}(\alpha)\), which quantifies the exponential amplification or decay of the disturbance along the streamwise direction. The complex streamwise wavenumber \(\alpha = \alpha_r + i\,\alpha_i\) obtained from the spatial eigenvalue problem provides two physically important quantities: the spatial growth rate and the phase speed of the disturbance. The spatial growth rate is defined as \(-\alpha_i\), quantifying the exponential amplification (\(\alpha_i < 0\)) or decay (\(\alpha_i > 0\)) of the disturbance envelope in the streamwise direction. The real part \(\alpha_r\) determines the streamwise wavelength, \(\lambda_x = 2\pi / \alpha_r\), and, together with the real frequency \(\omega\), defines the phase speed of the disturbance as
\[
  c_r = \frac{\omega}{\alpha_r}.
\]
This quantity represents the velocity at which a constant phase propagates along the surface. When normalised by the freestream velocity \(u_e\), the nondimensional phase speed \(c_r/u_e\) provides insight into the acoustic nature of the disturbance. For high-speed boundary layers, the second-mode instability typically propagates at speeds close to the local acoustic velocities relative to the wall. These can be approximated as
\[
  \frac{c_r}{u_e} \approx 1 \pm \frac{1}{M_e},
\]
where the \(+\) branch corresponds to the fast (upper) acoustic mode and the \(-\) branch corresponds to the slow (lower) acoustic mode. The former is associated with trapped acoustic waves reflecting between the wall and the relative sonic line, while the latter represents waves propagating in the opposite direction. The observed phase-speed values in spatial stability analyses and DNS therefore provide a direct indication of whether the dominant disturbance is associated with the upper or lower acoustic branch. Throughout this chapter, disturbance propagation is analysed using a base flow \(\bar{\boldsymbol{q}}\) supplied by the similarity solutions introduced in ~\cref{section:locallysimilar}, while the perturbations \(\tilde{\boldsymbol{q}}\) are modelled as a single species \gls{CPG}, as detailed in this section.

\subsection{Governing equations and viscosity effects}
% Governing equations (dimensionless compressible Navier–Stokes for a Newtonian fluid)
A number of previous studies have considered this issue. ~\cite{Wartemann2019} investigated the use of thermal excitation within the mean flow and stability code. Their results demonstrated that thermal excitation in both base flows and stability solver yielded results practically identical to those obtained when only the base flow included thermal excitation. Similar conclusions were drawn in \cite{MortensenZhong2013} amd \cite{miromiroThesis} and in 
\citep{BitterShepherd2015} for an air-mixture, where base flow thermal excitation had a larger effect on the final growth rates compared to thermal excitation in the perturbed flow. A conclusion was that a compressible stability solver with a constant specific heat $\gamma$ produced comparable outcomes to a solver that solved Linearised Navier Stokes Equations (LNSE) for thermally excited flows \cite{Wartemann2019}. Taking this into account, the current study employs a stability solver based on the \gls{CPG} assumption, while the base flow takes thermal excitation into account. The governing equations mentioned in this section form the basis for the linearised Navier--Stokes equations. Conservation of mass, momentum, and energy in the three spatial directions \(x_i\) (\(i = 0, 1, 2\)) yields a system of five partial differential equations written in terms of density \(\rho\), pressure \(p\), temperature \(T\), total energy \(E\), and velocity components \(u_k\):
\begin{align}
\frac{\partial \rho}{\partial t}
+ \frac{\partial}{\partial x_k}\!\left(\rho u_k\right) = 0, \\[0.5em]
\frac{\partial}{\partial t}\!\left(\rho u_i\right)
+ \frac{\partial}{\partial x_k}\!\left(\rho u_i u_k + p\,\delta_{ik} - \tau_{ik}\right) = 0, \\[0.5em]
\frac{\partial}{\partial t}\!\left(\rho E\right)
+ \frac{\partial}{\partial x_k}\!\left(
  \rho u_k\!\left(E + \frac{p}{\rho}\right) + q_k - u_i \tau_{ik}
\right) = 0,
\end{align}
where the heat flux \(q_k\) is given by Fourier’s law,
\begin{align}
q_k = -\,\kappa\,\frac{\partial T}{\partial x_k},
\end{align}
and \(\kappa\) denotes the thermal conductivity.

The equations are non-dimensionalised using the freestream reference quantities for velocity, density, and temperature \((U_\infty^*,\,\rho_\infty^*,\,T_\infty^*)\), together with a characteristic length based on the boundary-layer displacement thickness \(\delta^{*}\) prescribed at the inlet. The temperature-dependent dynamic viscosity \(\mu(T)\) is evaluated using the translational--rotational temperature for \gls{CPG} flows, following either Sutherland’s law, the Blottner--Eucken--Wilke rule, or the Musawi--Sandham model. A constant Prandtl number of \(0.72\) is employed in all cases. 





% Stability analysis investigates the fate of small, superposed disturbances and asks whether they amplify or decay as they propagate over a given base state. I

% In the local, parallel viewpoint this reduces to a wall-normal eigenproblem that yields growth rates, phase behaviour, and eigenfunctions—quantities that map neutral curves and most-amplified frequency bands and underpin \(N\)-factor estimates of transition onset. Two simplifying assumptions make the problem tractable. First, linearity: disturbances are taken to be sufficiently small that quadratic interaction terms can be neglected, so the governing equations about the base state are linear. Although transition is ultimately nonlinear, linear amplification is typically the slowest stage and therefore sets the effective transition Reynolds number, and the most-amplified linear frequency often persists into the nonlinear regime \citep{Reshotko2001,Mack1975}. Second, local parallelism: over a short streamwise extent the base flow varies only in the wall-normal direction, which is adequate for charting initial amplification and is well supported by classic experiments that confirmed instability waves and their quantitative trends \citep{SchubauerSkramstad1947}. Within these bounds, the framework provides a practical bridge from a computed mean flow to quantitative predictions used to guide DNS/LES set-ups, interpret measurements, and assess the influence of Mach number, wall temperature, and transport properties on disturbance growth \citep{SchmidHenningson2001}.

% The analysis proceeds by decomposing the instantaneous state into a base component and a small fluctuation,
% \begin{equation}
%   \boldsymbol{q}(x_j,t) \;=\; \bar{\boldsymbol{q}}(x_j,t) \;+\; \tilde{\boldsymbol{q}}(x_j,t),
% \end{equation}
% where \(\bar{\boldsymbol{q}}=(\bar{\rho},\,\bar{u},\,\bar{v},\,\bar{w},\,\bar{T})^{\mathsf{T}}\) denotes the base flow and \(\tilde{\boldsymbol{q}}=(\tilde{\rho},\,\tilde{u},\,\tilde{v},\,\tilde{w},\,\tilde{T})^{\mathsf{T}}\) the small disturbance. The boundary layer is deemed unstable where the disturbance amplitude increases downstream (or in time), and stable where it decays.




% The propagation of small amplitude disturbances in the flow field can be accurately captured by the linear stability framework. 


% Linear stability theory examines the evolution of small-amplitude disturbances superposed on a steady base flow. In the local viewpoint relevant to parallel shear layers, the base state is assumed to vary only in the wall-normal direction over a short streamwise extent, rendering the streamwise and spanwise directions effectively homogeneous. With this simplification, the linearised equations reduce to a one-dimensional eigenvalue problem for the wall-normal structure of the perturbations, yielding growth rates and phase information that indicate where instability first sets in. The approach is accurate so long as disturbance amplitudes remain small and departures from parallelism are modest, and it provides a practical bridge between a computed base flow and predictions of initial amplification prior to nonlinear breakdown. \citep{SchmidHenningson2001}


% LST assumes a parallel base flow given by
% \begin{equation}
%   \bar{\rho},\,\bar{u},\,\bar{T} = f(y) \qquad \bar{v} = \bar{w} = 0,
% \end{equation}
% where $\bar{\rho}, \bar{u}, \bar{T}$ are the base-flow density, streamwise velocity, and temperature, with $\bar{v}$ and $\bar{w}$ the wall-normal and spanwise velocities, respectively. Disturbances are assumed to have the form of travelling waves,
% \begin{equation}
%   \tilde{\boldsymbol{q}}(x,y,z,t) = \hat{\boldsymbol{q}}(y)\, \exp\!\big(i(\alpha x + \beta z - \omega t)\big),
% \end{equation}
% where $\hat{\boldsymbol{q}}(y)$ is the complex eigenfunction, $\alpha$ and $\beta$ are the streamwise ($x$) and spanwise ($z$) wavenumbers, and $\omega$ is the disturbance frequency. Substitution of the normal-mode ansatz into the linearised Navier--Stokes equations yields a linear eigenvalue problem,
% \begin{equation*}
%   L\,\hat{\boldsymbol{q}} = \omega\,\hat{\boldsymbol{q}}.
% \end{equation*}
% The interpretation hinges on the spectral parameters $(\alpha,\beta,\omega)$. In what follows, a spatial analysis is adopted: disturbances grow in space with $\alpha \in \mathbb{C}$ and $\omega \in \mathbb{R}$, and $\alpha$ is iterated until $\omega$ is real. In this setting, the spatial growth rate is $-\operatorname{Im}(\alpha)$, which quantifies the exponential amplification or decay of the disturbance along the streamwise direction.


% LST assumes a parallel base flow given by
% \begin{equation}
%     \Bar{\rho}, \Bar{u}, \Bar{T} = f(y) \quad \quad \Bar{v} = \Bar{w} = 0
% \end{equation} 
% where $\Bar{\rho}, \Bar{u}, \Bar{T}$  are the base flow density, streamwise velocity and temperature with $\Bar{v}$ and $\Bar{w}$ respectively the wall-normal velocity and spanwise velocity. Disturbances are assumed to have the form of travelling waves according to
% \begin{equation}
%     \Tilde{\boldsymbol{q}} = \hat{q}(y)  e^{(i(\alpha x + \beta z - \omega t))}
% \end{equation} 
% where $\hat{q}(y)$ is the complex eigenfunction, $\alpha$ and $\beta$ are the streamwise ($x$) and spanwise ($z$) wavenumbers, and $\omega$ is the disturbance frequency. When the normal mode solution is substituted into the LNSE, a linear system is obtained, which can be represented as \begin{equation*}
%     L \boldsymbol{\hat{q}} = \omega \boldsymbol{\hat{q}}.
% \end{equation*}
% The form of the wavenumbers and frequency depends on the type of the problem. Here we solve the spatial stability problem, i.e.~with disturbances growing in space, $\alpha \to \mathbb{C} \ \text{and} \ \omega \to \mathbb{R}$, by iterating on $\alpha$ until $\omega$ is real. 


% The similarity solutions of the previous section provide a basis for linear stability analysis. The question then arises of the influence of thermal effects on small perturbations. A number of previous studies have considered this issue. ~\cite{Wartemann2019} investigated the use of thermal excitation within the mean flow and stability code. Their results demonstrated that thermal excitation in both base flows and stability solver yielded results practically identical to those obtained when only the base flow included thermal excitation. Similar conclusions were drawn in \cite{Bitter2015,Mortensen2013,miromiroThesis}, where base flow thermal excitation had a larger effect on the final growth rates compared to thermal excitation in the perturbed flow. A conclusion was that a compressible stability solver with a constant specific heat $\gamma$ produced comparable outcomes to a solver that solved Linearised Navier Stokes Equations (LNSE) for thermally excited flows \cite{Wartemann2019}. Taking this into account, the current study employs a stability solver based on the CPG assumption, while the base flow takes into account thermal excitation. 



% \subsection{Viscosity effects}

\subsection{$e^N$ method}

% For a given frequency and Reynolds number, the locally parallel analysis yields only the \emph{local} spatial growth rate, $\operatorname{Im}(\alpha)$. Transition prediction requires the \emph{net} amplification accumulated as the disturbance convects downstream. In the $e^N$ framework \citep{SmithGamberoni1956,vanIngen1956}, the disturbance amplitude $A(x)$ relative to its value $A(x_0)$ at a reference location $x_0$ is
% \begin{equation}
%   N(x) \;=\; \ln\!\left(\frac{A(x)}{A(x_0)}\right)
%           \;=\; - \int_{x_0}^{x} \operatorname{Im}\!\left(\alpha(\xi)\right)\, \mathrm{d}\xi,
% \end{equation}


% so that $A(x)=A(x_0)\,e^{N(x)}$. In practice, $N$ is maximised over frequency to form an envelope, and a ``critical'' value is calibrated against experiments; in conventional (noisy) tunnels, transition is often observed for $N\approx5$--$10$ \citep{Schneider2015}. Because initial disturbance levels are not included, the calibrated threshold depends on facility, geometry, and disturbance environment; a fully predictive approach would couple this amplification with receptivity-based initial amplitudes, but in their absence the calibrated $e^N$ method provides a compact correlation once tuned to the configuration.


For a given frequency and Reynolds number, the locally parallel stability analysis delivers only the \emph{local} spatial growth rate, $\operatorname{Im}(\alpha)$, whereas transition prediction requires the \emph{net} amplification accumulated as a disturbance convects downstream. The $e^N$ method \citep{SmithGamberoni1956,vanIngen1956} provides this link by integrating the spatial growth rate to obtain the amplitude ratio relative to a reference location $x_0$:
\begin{equation}
    N \ = \ln{ \left( \cfrac{A(x)}{A(x_0)} \right )} \ = - \int_{x_0}^{x} \alpha_i \ \dd x
\end{equation}
so that $A(x)=A(x_0)\,e^{N(x)}$. In practice $N(x)$ is maximised over frequency to form an envelope used for correlation with experiments. A “critical” value of $N$ is then calibrated: in conventional (noisy) wind tunnels, transition is commonly observed for $N \approx 5$–$10$ \citep{Schneider2015}, corresponding to an amplification of $e^{5}$–$e^{10}\!\approx\!150{,}000$–$22{,}000$. 

It should be stressed that transition occurs when disturbance amplitudes become large enough for nonlinear interactions to matter; thus the absolute amplitude $A$—not merely the ratio $A/A_0$ or its logarithm $N$—is the physically relevant quantity. An ideal predictive framework would therefore couple the linear amplification computed here with receptivity-based initial amplitudes derived from identified disturbance sources. Because such initial levels are not included explicitly, the calibrated “critical $N$” depends on facility, geometry, and disturbance environment. Nonetheless, once tuned to a given configuration, the $e^N$ approach has demonstrated practical predictive value.


% \subsection{N-factor analysis}
% For a particular frequency and Reynolds number, the locally-parallel stability analysis presented above yields only the local growth rate, $\alpha$i. However, in predicting boundary layer transition one is less interested in the local growth rate and more interested in the net amplification that a disturbance experiences as it travels down- stream. An intuitive method of predicting this net growth is to simply integrate the spatial growth rate as a function of downstream distance. The amplitude A of the disturbance relative to its initial amplitude Ao at some point xo is then


% \begin{equation}
%     N \ = \ln{ \left( \cfrac{A(x)}{A(x_0)} \right )} \ = - \int_{x_0}^{x} \alpha_i \ \dd x
% \end{equation} 

% The maximum N factor (maximized over all frequencies) is often used to correlate experimental measurements of the transition distance with stability calculations, an approach which was pioneered by Smith and Gamberoni (1956). For example, boundary layer transition in conventional wind tunnels is usually observed at the Reynolds number for which N = 5-10, which corresponds to a disturbance amplification of 150-22,000. Therefore it is often assumed that transition will occur at some “critical N factor of transition” in the range of 5-10, where the precise value of the critical N factor is calibrated using experimental data. This method of estimating transition is known in the literature as the e\textsuperscript{th} method. Of course, the reality is that transition occurs when the disturbances reach a
% large enough amplitude that nonlinear interactions become important, so the actual amplitude A is more relevant to transition than the amplification ratio A/Ao or N. An ideal transition prediction method would therefore couple the amplifications predicted here with initial amplitudes obtained from receptivity calculations based on known disturbance sources in the flow. Unfortunately, theories of receptivity have not yet been developed to the point where this is feasible, and the idea of a “critical N factor” for transition is used instead. But because this method does not account for the initial disturbance amplitudes, the precise value of N at which transition occurs is different for each flow geometry and wind tunnel facility. However, there is some evidence that the eN method described above has predictive value once calibrated to a particular flow configuration (Schneider, 2015).

% The eigenvalue problem is solved using a Chebyshev spectral method with geometric mapping. \cref{table:lstcomp} shows a validation of both the similarity solution and the LST eigenvalues (using $n_{\mathrm{cheb}}$=200) against case IV in ~\cite{Groot2018}, showing agreement up to the fifth decimal place in the real and imaginary parts of the wavenumber $\alpha$. Finally, the growth in amplitude is obtained as an $N$-factor by integrating the LST growth rate according to
% \begin{equation}\label{eq:lstnfactor}
%     N \ = \ln{ \left( \cfrac{A(x)}{A(x_0)} \right )} \ = \int_{x_0}^{x} - \alpha_i \ \dd x
% \end{equation} 



\subsection{Validation}

The \gls{LST} eigen value problem was solved using Chebyshev spectral differentiation method, with geometric wall-normal clustering introduced via the Malik mapping \citep{Malik1990a}, which maps $[-1,1]\to[0,\eta_{\max}]$. In the spatial formulation adopted here, the problem is posed for complex $\alpha$ at prescribed $(\omega,\beta)$, yielding a generalised eigenvalue problem whose spectrum provides $\alpha_r$ and $\alpha_i$ (with spatial growth rate $-\alpha_i$). The framework is compared against \citet{Groot2018}, and the case-specific conditions in Table~\ref{table:lstcomp} are chosen to isolate Mack’s second mode.

\begin{table}[ht]
  \centering
  \caption{Case-specific parameters (subset of \citet{Groot2018}) and numerical settings used here.}
  \label{table:lstcomp}
  \vspace{0.25em}
  \begin{tabular}{@{}lccccccccc@{}}
    \toprule
    Case & $M$ & type & $\mathrm{Re}_\ell$ &
    $\overline{T}_e$ (K) & $\overline{T}_w$ (K) & $\mathrm{Re}$ &
    $\omega$ & $\beta$ & $n_{\mathrm{cheb}}$ \\
    \midrule
    IV & 2.5  & M1 & $5.0\cdot10^{6}$ & $600/4.05$ & adiab. & 3000 & 0.04  & 0.1 & 200 \\
    V  & 10.0 & M2 & $9.8425\cdot10^{6}$ & 278       & adiab. & 2000 & 0.075 & 0   & 500 \\
    \bottomrule
  \end{tabular}
\end{table}


Throughout the current validation study, the viscosity law follows Sutherland’s law with $T_{\mathrm{ref}}=273.15\ \mathrm{K}$, $S=110.6\ \mathrm{K}$, and $\mu_{\mathrm{ref}}=1.716\times10^{-5}\ \mathrm{kg/(m\,s)}$, and we use $Pr=0.70$ and $\gamma=1.4$, consistent with the benchmark setup of \citet{Groot2018}. \gls{LST} is extremely sensitive to the viscosity model, and slight changes in the viscosity and thermal conductivity can yield contrasting results. DEKAF is a spectral basic-state solver for boundary-layer stability that computes high-accuracy self-similar base flows across regimes using Chebyshev pseudo-spectral discretisation and Newton–Raphson iteration. It provides benchmark base states and integral metrics for validating \gls{LST} and assessing how basic-state accuracy propagates into stability predictions~\citep{MiroMiro2019}.


% The eigenvalue problem is solved using a Chebyshev spectral discretisation on CGL nodes with a geometric clustering of points in the wall-normal direction via the Malik mapping \citep{Malik1990a}, which transforms the computational interval $[-1,1]$ to $[0,\eta_{\max}]$. The resulting LST operator is assembled on this staggered grid and solved for complex $\alpha$ at prescribed $(\omega,\beta)$. \cref{table:lstcomp} validates both the similarity solution and the \gls{LST} eigenvalues (with $n_{\mathrm{cheb}}=200$ unless otherwise noted) against cases~IV and~V of \citet{Groot2018}, showing agreement to five decimal places in both $\Re(\alpha)$ and $\Im(\alpha)$. Throughout, viscosity follows Sutherland’s law with $T_{\mathrm{ref}}=273.15\ \mathrm{K}$, $S=110.6\ \mathrm{K}$, and $\mu_{\mathrm{ref}}=1.716\times10^{-5}\ \mathrm{kg/(m\,s)}$, and we use $Pr=0.70$ and $\gamma=1.4$, consistent with the benchmark setup of \citet{Groot2018}. The case-specific conditions in Table~\ref{table:lstcomp} are chosen to isolate Mack’s second mode.

% \begin{table}[ht]
%   \centering
%   \caption{Integral parameters for cases IV and V: DEKAF values (from \citet{Groot2018}) versus Present (this work). Present values are placeholders.}
%   \label{tab:integral_params_compare_iv_v}
%   \setlength{\tabcolsep}{5pt}
%   \small
%   \begin{tabular}{@{}l
%                   *{2}{c}
%                   *{2}{c}
%                   *{2}{c}
%                   *{2}{c}
%                   *{2}{c}@{}}
%     \toprule
%     & \multicolumn{2}{c}{$\delta^*/\ell$}
%     & \multicolumn{2}{c}{$\theta^*/\ell$}
%     & \multicolumn{2}{c}{$H^*=\delta^*/\theta^*$}
%     & \multicolumn{2}{c}{$\delta_e^*/\ell$}
%     & \multicolumn{2}{c}{$\delta_h^*/\ell$} \\
%     \cmidrule(lr){2-3}\cmidrule(lr){4-5}\cmidrule(lr){6-7}\cmidrule(lr){8-9}\cmidrule(lr){10-11}
%     Case & DEKAF & Present & DEKAF & Present & DEKAF & Present & DEKAF & Present & DEKAF & Present \\
%     \midrule
%     IV  & 4.2571  & 0 & 0.6391  & 0 & 6.6615  & 0 & 1.0086 & 0 & 1.2099 & 0 \\
%     V   & 27.0430 & 0 & 0.4188  & 0 & 64.5748 & 0 & 0.6745 & 0 & 0.8245 & 0 \\
%     \bottomrule
%   \end{tabular}
% \end{table}

% \begin{table}[ht]
%   \centering
%   \caption{Integral parameters for cases IV and V: DEKAF values (from \citet{Groot2018}) versus Present (this work).}
%   \label{tab:integral_params_compare_iv_v}
%   \setlength{\tabcolsep}{5pt}
%   \small
%   \begin{tabular}{@{}l*{2}{c}*{2}{c}*{2}{c}*{2}{c}*{2}{c}@{}}
%     \toprule
%     & \multicolumn{2}{c}{$\delta^*/\ell$}
%     & \multicolumn{2}{c}{$\theta^*/\ell$}
%     & \multicolumn{2}{c}{$H^*=\delta^*/\theta^*$}
%     & \multicolumn{2}{c}{$\delta_e^*/\ell$}
%     & \multicolumn{2}{c}{$\delta_h^*/\ell$} \\
%     \cmidrule(lr){2-3}\cmidrule(lr){4-5}\cmidrule(lr){6-7}\cmidrule(lr){8-9}\cmidrule(lr){10-11}
%     Case & DEKAF & Present & DEKAF & Present & DEKAF & Present & DEKAF & Present & DEKAF & Present \\
%     \midrule
%     IV  & 4.2571  & 4.2571  & 0.6391  & 0.6391  & 6.6615  & 6.6615  & 1.0086 & 1.0086 & 1.2099 & 1.2099 \\
%     V   & 27.0430 & 27.0430 & 0.4188  & 0.4188  & 64.5748 & 64.5748 & 0.6745 & 0.6745 & 0.8245 & 0.8245 \\
%     \bottomrule
%   \end{tabular}
% \end{table}

\begin{table}[ht]
  \centering
  \caption{Integral parameters for cases IV and V: DEKAF vs Present.}
  \label{tab:integral_params_compare_iv_v}
  \setlength{\tabcolsep}{6pt}
  \small
  \begin{tabular}{@{}lcccc@{}}
    \toprule
    & \multicolumn{2}{c}{Case IV} & \multicolumn{2}{c}{Case V} \\
    \cmidrule(lr){2-3}\cmidrule(lr){4-5}
    Metric & DEKAF & Present & DEKAF & Present \\
    \midrule
    $\delta^*/\ell$        & 4.2571  & 4.2571  & 27.0430 & 27.0430 \\
    $\theta^*/\ell$        & 0.6391  & 0.6391  & 0.4188  & 0.4188  \\
    $H^*=\delta^*/\theta^*$& 6.6615  & 6.6615  & 64.5748 & 64.5748 \\
    $\delta_e^*/\ell$      & 1.0086  & 1.0086  & 0.6745  & 0.6745  \\
    $\delta_h^*/\ell$      & 1.2099  & 1.2099  & 0.8245  & 0.8245  \\
    \bottomrule
  \end{tabular}
\end{table}




% To indicate the accuracy of the solutions in a different way, the boundary-layer integral parameters are evaluated. In particular, the displacement thickness $\delta^*$, the momentum thickness $\theta^*$, the shape factor $H^*$, the energy thickness $\delta_e^*$, and the enthalpy thickness $\delta_h^*$ are tabulated. The resulting values of the integral parameters are reported in Table~5, whereas the velocity-gradient values and mapping parameters, $y_i$ and $y_{\max}$, are reported in Table~\ref{tab:integral_params_compare_iv_v}. As can be seen, the similarity solution presented in the previous chapter yields results that match the reference to five decimal places, thereby increasing confidence in the method. Note that the similarity solution is computed using a Chebyshev differentiation approach and required $N=200$ points to converge the shape factor $H^*$. The similarity framework used for comparison is the five-species frozen \gls{CPG} formulation.
As a first assessment of accuracy, boundary-layer integral parameters are evaluated: the displacement thickness $\delta^*$, momentum thickness $\theta^*$, shape factor $H^*$, energy thickness $\delta_e^*$, and enthalpy thickness $\delta_h^*$. The resulting values are reported in Table~5, while the velocity-gradient metrics and mapping parameters, $y_i$ and $y_{\max}$, are given in Table~\ref{tab:integral_params_compare_iv_v}. The similarity solution presented in the previous chapter matches the reference to five decimal places, increasing confidence in the method. It is computed using a Chebyshev differentiation approach and required $N=200$ collocation points to converge $H^*$. The comparison uses the five-species frozen \gls{CPG} formulation.



% \begin{table}[ht]
% \renewcommand{\arraystretch}{1.4}
% \centering
% \caption{Validation of the stability solver for M=2.5, adiabatic wall, comparing to \cite{Groot2018}} \label{table:lstcomp}
% \begin{tabular}{|l|c|c|c|}
% \hline
% %\toprule
%  & $\Delta^{*}$ & $\alpha_r$ & $\alpha_i$  \\
%  \hline
% %\midrule
% Current Solver & 4.257109 & 107.402617 & -1.039767  \\ 
% DEKAF ~\cite{Groot2018} & 4.257110 & 107.402614 &-1.039777  \\ 
% %\bottomrule
% \hline
% \end{tabular}
% \end{table}




\begin{table}[ht]
  \centering
  \caption{Eigenvalues obtained using Chebyshev spectral collocation.}
  \label{tab:eigs_iv_v}
  \setlength{\tabcolsep}{7pt}
  \begin{tabular}{@{}lcccc@{}}
    \toprule
    \multicolumn{5}{c}{\textbf{IV}}\\
    \midrule
    solver & $\alpha_r\;(\mathrm{rad/m})$ & $\epsilon(\alpha_r)$
           & $\alpha_i\;(\mathrm{rad/m})$ & $\epsilon(\alpha_i)$ \\
    \midrule
    VESTA   & $1.07402614379007\cdot 10^{2}$ & $9.4\cdot 10^{-13}$ &
               $-1.03977656691188\cdot 10^{0}$ & $5.1\cdot 10^{-12}$ \\
    TUD     & $1.07402614379108\cdot 10^{2}$ & --- &
               $-1.03977656636347\cdot 10^{0}$ & --- \\
    Current & $1.07402619911674\cdot 10^{2}$  & $5.2\cdot 10^{-8}$ &
               $-1.03975787207327\cdot 10^{0}$ & $1.8\cdot 10^{-5}$ \\
    \midrule
    \multicolumn{5}{c}{\textbf{V}}\\
    \midrule
    solver & $\alpha_r\;(\mathrm{rad/m})$ & $\epsilon(\alpha_r)$
           & $\alpha_i\;(\mathrm{rad/m})$ & $\epsilon(\alpha_i)$ \\
    \midrule
    VESTA   & $3.86915377502036\cdot 10^{2}$ & $4.7\cdot 10^{-12}$ &
               $-7.98862348202598\cdot 10^{0}$ & $3.5\cdot 10^{-12}$ \\
    TUD     & $3.86915377503853\cdot 10^{2}$ & --- &
               $-7.98862348069057\cdot 10^{0}$ & --- \\
    Current & $3.86914179040329\cdot 10^{2}$  & $3.1\cdot 10^{-6}$ &
               $-7.98994070166233\cdot 10^{0}$ & $1.6\cdot 10^{-4}$ \\
    \bottomrule
  \end{tabular}
\end{table}




% Finally, Table~\ref{tab:eigs_iv_v} compares the complex wavenumber $\alpha=\alpha_r+i\alpha_i$ from the present spatial LST solver against the DEKAF benchmarks, which themselves were verified against the VESTA stability code~\citep{Pinna2012} and the TU~Delft solver~\citep{Groot2013}. For both case~IV ($M=2.5$) and case~V ($M=10$), the present $\alpha_r$ and $\alpha_i$ agree with DEKAF to at least five decimal places, confirming consistency of the growth rate (via $-\alpha_i$) and phase speed (via $\alpha_r$). Figure~\ref{figure:lst_groot_comp_eigenfunc} further compares the corresponding eigenfunctions at $M=2.5$ and $M=10$ (cases~IV–V), normalised by $\max|\tilde{u}|$. The wall-normal amplitude distributions from the present solver (solid lines) closely overlay the DEKAF profiles (markers), capturing the characteristic second-mode structure in both cases. Case~IV achieves convergence with $N=200$ collocation points, whereas case~V requires $N=500$, consistent with the higher frequency and sharper acoustic trapping of the Mack second mode at $M=10$.
 
Finally, Table~\ref{tab:eigs_iv_v} compares the complex wavenumber $\alpha=\alpha_r+i\alpha_i$ from the present spatial \gls{LST} solver with the DEKAF benchmarks, themselves verified against the VESTA code~\citep{Pinna2012} and the TU~Delft solver~\citep{Groot2013}. For both case~IV ($M=2.5$) and case~V ($M=10$), the present $\alpha_r$ and $\alpha_i$ agree with DEKAF to at least five decimal places, confirming consistency in growth rate (via $-\alpha_i$) and phase speed (via $\alpha_r$). Figure~\ref{figure:lst_groot_comp_eigenfunc} further compares the mode shapes for $M=2.5$ and $M=10$, normalised by $\max|\tilde{u}|$. The wall-normal amplitude distributions from the present solver closely overlay the DEKAF profiles, capturing the perturbations of characteristic second-mode structure in both cases. Case~IV converged with $N=200$ collocation points, whereas case~V requires $N=500$, consistent with the higher frequency and stronger acoustic trapping of the Mack second mode at $M=10$.

\begin{figure}[ht]
    \centering
    \fontsize{9}{10}\selectfont
    \includegraphics[width=0.70\textwidth]{resources/figures/disturbance/result_lst_groot_comparison_DEKAF.pdf}
    \caption[Eigenfunction profiles for cases IV–V (DEKAF \citep{Groot2018} vs Current).]%
    {Amplitude profiles of the spatial LST eigenfunctions plotted versus the wall-normal coordinate $y/\delta^{*}$. Curves are normalised by $\max|\tilde{u}|$. Solid lines denote the \emph{Current} solver and circular markers denote \emph{DEKAF}~\citep{Groot2018}. \textbf{(a)} Case~IV; \textbf{(b)} Case~V.}
    \label{figure:lst_groot_comp_eigenfunc}
\end{figure}



% \section{DNS domain setup and configuration}
\section{Computational domain and configurations}

% The configuration in the current study is shown in Figure~\ref{figure:schematic_flow}, a zero-pressure gradient flat plate boundary layer.  The wedge is in a Mach 10 flow at 36 km altitude, chosen to be the limit of a hypersonic glide vehicle trajectory. A standard model of the atmosphere adopted by the International Civil Aviation Organisation (ICAO) is employed to get the atmospheric conditions at 36 km, giving rise to the post-shock freestream conditions shown in Table \ref{table:flight_conditions}.


% The configuration in the current study is shown in ~\cref{figure:schematicflow}. A wedge is placed in a Mach 10 flow at 36 km altitude, chosen to be the heat rate limit of a hypersonic glide vehicle trajectory \cite{Rizvi2015}. A standard model of the atmosphere adopted by the International Civil Aviation Organisation (ICAO) is employed to get the atmospheric conditions at 36 km, i.e.~$\mathrm{p}_{\mathrm{atm}} = 498.521 \ \mathrm{Pa}$ and $\mathrm{T}_{\mathrm{atm}} = 239.282 \ \mathrm{K}$, with the molar concentrations given as $X_{\ce{N2}} = 0.79$ and $X_{\ce{O2}}=0.21$  giving rise to an atmospheric stagnation enthalpy $H_0 = h_{atm} + u^2_{atm}/2 = 5.07 MJ/kg$. 

The configuration considered here is shown in~\cref{figure:schematicflow}. A wedge is placed in a Mach~10 freestream at an altitude of 36~km, selected to represent the heat-rate limit along a hypersonic glide-vehicle trajectory~\cite{Rizvi2015}. Atmospheric properties at 36~km are obtained from the International Civil Aviation Organisation (ICAO) standard atmosphere and are used without modification, i.e.\ $\mathrm{p}_{\mathrm{atm}} = 498.521~\mathrm{Pa}$ and $\mathrm{T}_{\mathrm{atm}} = 239.282~\mathrm{K}$. The composition is taken as a binary mixture with molar concentrations $X_{\ce{N2}}=0.79$ and $X_{\ce{O2}}=0.21$, which, together with the stated thermodynamic state, define the reference conditions for the flow. The corresponding atmospheric stagnation enthalpy is evaluated as $H_{0}=h_{\mathrm{atm}}+u_{\mathrm{atm}}^{2}/2=5.07~\mathrm{MJ/kg}$, where $h_{\mathrm{atm}}$ and $u_{\mathrm{atm}}$ are the static enthalpy and flow speed consistent with the ICAO values.


\begin{figure}[ht]
    \centering
    \fontsize{8}{10}\selectfont\includesvg[width=0.95\textwidth]{figures/disturbance/schematic-flightconditions_aiaa2.svg}
    \caption{ Schematic of the geometry considered in the current study, showcasing a domain behind the shock of a wedge. The post-shock freestream conditions of the flow over the wedge are provided in \cref{table:flight_conditions}}
    \label{figure:schematicflow}
\end{figure}

% The rectangular computational domain highlighted in the figure is placed downstream of the attached oblique shock generated by the wedge. The corresponding post-shock free-stream conditions are calculated using the shock jump relations, and the mixture composition of air at equilibrium conditions is checked using the Park et al.~\cite{park2001} model. A full breakdown of flow conditions is mentioned in \cref{table:flight_conditions} where $\vartheta$ is the wedge angle altered to get two flow conditions (at Mach 6 and 8 respectively). A subscript $e$ denotes free-stream conditions at the inflow to the computational domain. The two rows, one for each wedge angle, show only small differences in the post-shock velocity through the assumption of \gls{CPG} or \gls{TEQ} conditions. 

The rectangular computational domain highlighted in the figure is placed immediately downstream of the attached oblique shock generated by the wedge, i.e.\ in the post-shock region where the flow has already been turned and compressed by the shock induced at the wedge surface. The corresponding post-shock free-stream conditions are calculated using the standard shock jump relations, and the mixture composition of air at equilibrium conditions is checked for consistency using the ~\cite{Park1990} model, ensuring that thermodynamic state and species partitioning are aligned. A full breakdown of the flow conditions is provided in \cref{table:flight_conditions}, where $\vartheta$ denotes the wedge angle and is altered to realise two nominal operating points (Mach~6 and Mach~8, respectively), with all other quantities reported in a consistent format. Throughout, a subscript $e$ designates free-stream conditions evaluated at the inflow to the computational domain, making clear which values are prescribed as boundary data. The two rows of the table one corresponding to each wedge angle show only small, practically negligible differences in the post-shock streamwise velocity when adopting either \gls{CPG} or \gls{TEQ} conditions, indicating that the inflow kinematics are only weakly sensitive to the chosen thermodynamic closure for these cases.

% The profiles of the conservative variables prescribed at the inflow are obtained from the similarity solution procedure given in \cref{subsection:similarity}. 

% \begin{table}[H]
% \renewcommand{\arraystretch}{1.4}
% \centering
% \caption{ 2D wedge post-shock conditions for h=36 km,  $M_{\infty}=10$ }\label{table:flight_conditions}
% \begin{tabular}{|l|c|c|c|c|c|c|}
% \hline
% %\toprule
% $\vartheta$ & $M_e$  & $T_e$ (K) & $p_e$ (kPa) &  $\rho_e$ (kg/m\textsuperscript{3})  & $u_e$ (m/s) & type  \\
% \hline
% %\midrule
% \multirow{2}{*}{12.41} &\multirow{2}{*}{6} & \multirow{2}{*}{620.0} &\multirow{2}{*}{4.91} &\multirow{2}{*}{2.748$\times$10\textsuperscript{-3} } & 3001.09 & \gls{CPG}  \\ 
%  & &  &   & &  2972.59 & \gls{TEQ}
%  \\
%  \hline
% % \midrule
% \multirow{2}{*}{5.88} &\multirow{2}{*}{8} & \multirow{2}{*}{365.0} &\multirow{2}{*}{1.78} & \multirow{2}{*}{1.6920$\times$10\textsuperscript{-3}}& 3070.21 &   \gls{CPG}   \\ 
% &  &  &   & & 3066.27   & \gls{TEQ}  \\
% \hline
% %\bottomrule
% \end{tabular}
% \end{table} 

\begin{table}[H]
\renewcommand{\arraystretch}{1.4}
\centering
\caption{2D wedge post-shock conditions for $h=36~\mathrm{km}$, $M_{\infty}=10$.}
\label{table:flight_conditions}
\setlength{\tabcolsep}{6pt}
\small
\begin{tabular}{@{}lcccccc@{}}
\toprule
$\vartheta$ & $M_e$  & $T_e$ (K) & $p_e$ (kPa) & $\rho_e$ (kg/m\textsuperscript{3}) & $u_e$ (m/s) & type \\
\midrule
\multirow{2}{*}{12.41} & \multirow{2}{*}{6} & \multirow{2}{*}{620.0} & \multirow{2}{*}{4.91} & \multirow{2}{*}{2.748$\times$10\textsuperscript{-3}} & 3001.09 & \gls{CPG} \\
                       &                         &                         &                         &                                             & 2972.59 & \gls{TEQ} \\
\addlinespace[0.5ex]
\midrule
\multirow{2}{*}{5.88}  & \multirow{2}{*}{8} & \multirow{2}{*}{365.0} & \multirow{2}{*}{1.78} & \multirow{2}{*}{1.6920$\times$10\textsuperscript{-3}} & 3070.21 & \gls{CPG} \\
                       &                     &                        &                        &                                                & 3066.27 & \gls{TEQ} \\
\bottomrule
\end{tabular}
\end{table}


% For two-dimensional calculations, the domain is set as $L_x \times L_y$ = $750\delta^{*}_{0} \times 100\delta^{*}_{0}$ in $x$ and $y$ direction. The streamwise extent equates to $\approx 2.5 \ m$ and $\approx 2.7 \ m$ for the Mach 6 and 8 cases respectively. The computational grid is $N_x \times N_y \ = \  2500 \times 600$ points in the streamwise and wall-normal direction respectively, giving 1.5 million grid points altogether. The points are equally spaced in the streamwise direction, and a grid stretching is applied in the wall-normal direction $y$ to ensure the $\Delta y^{+} = 1$ with points clustered hyperbolically towards the wall as \begin{equation}
%     y = L_y \cfrac{\sinh{s_1 \xi}}{\sinh{s_1}}
% \end{equation} where $s_1$ is the stretching factor, $s_1 = 5.0$ and $\xi = (0,1)$ is an auxiliary coordinate giving approximately $180$ points in the boundary layer at the exit of the domain. Finally, the non-dimensional time step for the simulation was $\Delta t = 0.01$, achieving an acoustic CFL number $CFL = (u+a_f)\Delta t/\Delta x$ of 0.083, comparable to literature. The Reynolds number based on inflow displacement thickness is chosen to be $Re_{\delta^{*}_0} \ = 10,000$ translating to respectively around 0.3 m and 0.5 m behind the leading edge for both cases.
For the two-dimensional calculations, the post-shock computational domain is prescribed as $L_x\times L_y=800\,\delta^{*}_{0}\times 100\,\delta^{*}_{0}$ in the streamwise and wall-normal directions, respectively, which corresponds to a physical streamwise length of $\approx 2.9~\mathrm{m}$ for the Mach~6 cases and $\approx 3.2~\mathrm{m}$ for the Mach~8 cases. The mesh comprises $N_x\times N_y=2800\times 600$ points i.e.\ $\sim 1.68\times 10^{6}$  points in total, which are uniformly spaced in $x$ while a stretched distribution incorporated in $y$ obtained by hyperbolic stretching to guarantee near-wall resolution of $\Delta y^{+}\approx 1$. The wall-normal mapping employed is
\begin{equation}
y=L_y\,\frac{\sinh\!\left(s_1\,\xi\right)}{\sinh s_1},\qquad s_1=5.0,\quad \xi\in(0,1),
\end{equation} which clusters nodes towards the wall and yields approximately $180$ points inside the boundary layer by the downstream end of the domain. Time advancement uses a non-dimensional time step \(\Delta t=0.01\), giving an acoustic Courant number $\mathrm{CFL}=(u+a_f)\,\Delta t/\Delta x=0.083$, consistent with values commonly reported in the literature for comparable resolutions and Mach numbers. The inflow Reynolds number based on the displacement thickness is fixed at $Re_{\delta^{*}_0}=10{,}000$, which places the inflow section at roughly $\approx 0.5~\mathrm{m}$ downstream of the leading edge for the Mach~6 cases configuration and $\approx 0.8~\mathrm{m}$ for the Mach~8 configuration, thereby standardising the viscous scaling at entry. It is important to note that the initial distance from the leading edge is heavily dependent on the thermal excitation, Mach number and wall temperature, therefore the values reported are slightly different for individual cases. 

Discretisation and marching follow a high-order strategy: unless otherwise stated, the spatial derivatives are evaluated with a fourth-order central-difference scheme augmented by Feiereisen splitting to stabilise convective transport in the presence of strong aliasing errors, while temporal evolution is advanced with a low-storage, three-stage Runge–Kutta (RK3) method. Boundary conditions are imposed to minimise spurious reflections and preserve the prescribed inflow state: an inlet pressure transfer condition is applied at $x=x_d$, simple extrapolation is used at the outflow plane, and the farfield boundary is likewise treated by extrapolation consistent with the outer state. The wall is modelled as non-catalytic and isothermal in thermal equilibrium, with temperatures set as $T=T_v\in\{1200~\mathrm{K},\,1800~\mathrm{K}\}$ depending on the case, selected to lie within the melting-range envelope of typical materials used in thermal-protection systems (see \cite{CedillosBarraza2016} for a detailed survey). The thermal-equilibrium assumption at the wall is supported by experiments showing that molecules leaving a heated surface equilibrate vibrationally at the wall \cite{Park1990,CandlerNompelis2021}. Consistent with this, prior stability analyses \cite{Hudson1997,MiroMiro2019,BitterShepherd2015} impose a thermally equilibrated wall, and DNS studies of hypersonic boundary layers \cite{Passiatoreetal2021,DiRenzo2021} have adopted comparable wall treatments under both equilibrium and nonequilibrium bulk-flow conditions.


% Together, these choices ensure a uniformly resolved wall layer, controlled numerical dissipation in regions of steep gradients, and a boundary treatment that neither contaminates the interior solution nor constrains the downstream development. 


% As mentioned previously, the conservative variables are initialised by solving the similarity solution which is transferred into the domain through an inlet pressure extrapolation boundary condition. The outflow and far-field boundary conditions are treated using simple extrapolation and zero-gradient boundary conditions respectively. The wall is considered to be non-catalytic and isothermal in thermal equilibrium, with temperatures set as $T = T_v = 1200\mathrm{K}, 1800\mathrm{K}$ depending on the case, chosen to be within the range of melting points of various materials used in the thermal protection systems (a detailed breakdown of materials can be found in Cedillos-Barrazaeta et al. \cite{cedillosBarrazaetal2016}). The assumption of thermal equilibrium at the wall is justified as experiments measuring the vibrational temperature of molecules leaving a heated surface have shown that they equilibrate at the wall \cite{Park1990,candler2021von}. Furthermore, various stability studies by Hudson et al.~\citep{hudson1997}, Miro-Miro \cite{miromiroThesis} and Bitter~\cite{Bitter2015} have considered a wall in thermal equilibrium, and DNS studies previously mentioned of Passiatore~\cite{Passiatore2021} and Di Renzo and Urzay~\cite{direnzo_urzay_2021} have also employed similar approaches for flows in thermal equilibrium or nonequilibrium

% \subsubsection{Boundary and initial conditions}
\subsection{Wall treatment}
To introduce controlled, small-amplitude perturbations, we employ a wall suction--blowing strip following \cite{AlamSandham2000}, adapted to the two-dimensional (spanwise-uniform) domain. A Gaussian-localised disturbance centred at $x_f$ (see \cref{figure:schematicflow}) is imposed as a wall mass-flux perturbation
\begin{equation}
  \rho v \;=\; a_0\, g(x)\, \sum_{i=1}^{8} \sin\!\big(2\pi f_i t\big),
\end{equation}
with $a_0 = 1\times10^{-5}\,u_e$ the prescribed forcing amplitude, chosen to keep the response in the linear regime. The streamwise shape function is
\begin{equation}
  g(x) \;=\; \exp\!\left[-\,\frac{(x-x_f)^2}{2\sigma^2}\right],
\end{equation}
where \(2\sigma^2 = 20(\delta_0^{*})^2\), we take \(\sigma=\sqrt{10}\,\delta_0^{*}\), and centre the forcing at \(x_f = 25\,\delta_0^{*}\). The Gaussian’s effective width is \(2\sqrt{2\ln 2}\,\sigma \approx 2.35\,\sigma\), so the core spans several grid points. Outside the Gaussian support the wall is impermeable \((\rho v=0)\), and the imposed signal has zero temporal mean, so the net mass flux is zero. Outside the Gaussian support the wall remains impermeable ($\rho v=0$), and the temporal mean of the imposed signal is zero, preserving global mass balance.

The non-dimensional frequency is defined as $\tilde f = f\,\delta_0/u_e$. Discrete tones are selected so that the resulting wavepackets amplify within $x/\delta_0^{*}\in[0,\,800]$. For Mach~6 we use eight tones with $0.10 \le \tilde f \le 0.18$; for Mach~8, $0.08 \le \tilde f \le 0.16$, yielding $f^{*}\approx 80$–$152~\mathrm{kHz}$ in both cases. These bands target Mack-mode receptivity under elevated post-shock pressure and temperature and are consistent with accompanying \gls{LST} predictions. Numerically, suction–blowing is imposed as a Dirichlet condition on the wall-normal mass flux $(\rho v)$ over the Gaussian support; elsewhere a no-penetration condition $(\rho v=0)$ is enforced. This construction provides a compact, well-resolved, reproducible, inflow-independent disturbance source that excites the intended instability content while maintaining small perturbation levels.



% where $2\sigma^2 = 20\,(\delta_0^{*})^2$ and $x_f = 25\,\delta_0^{*}$. The corresponding width in full-width-at-half-maximum is $\mathrm{FWHM}=2\sqrt{2\ln 2}\,\sigma$, ensuring multiple grid points across the core of the forcing envelope

% The non-dimensional frequency is \(\tilde f = f\,\delta_0/u_e\). Discrete tones are chosen so that the resulting wavepackets amplify within \(x/\delta_0^{*}\!\in[0,\,750]\). For Mach~6 we use eight tones with \(0.10 \le \tilde f \le 0.18\); for Mach~8, \(0.09 \le \tilde f \le 0.17\), yielding \(f^{*}\approx 80\)–\(145~\mathrm{kHz}\) in both cases. These bands target Mack-mode receptivity under the elevated post-shock pressure and temperature and are consistent with the accompanying \gls{LST} predictions. Numerically, suction–blowing is applied as a Dirichlet condition on the wall-normal mass flux \((\rho v)\) at boundary nodes within the Gaussian support; elsewhere the no-penetration condition \((\rho v=0)\) is enforced. This construction provides a compact, well-resolved, reproducible, inflow-independent disturbance source that excites the intended instability content while maintaining small perturbation levels.




% The non-dimensional frequency is defined as $\tilde f = f\,\delta_0/u_e$, and the discrete tones are selected so that the resulting wavepackets amplify within $x/\delta_0^{*}\in[0,\,750]$. For Mach~6 we employ eight tones spanning $0.10 \le \tilde f \le 0.18$ (corresponding to $f^{*}\approx 80$–$145~\mathrm{kHz}$); for Mach~8, $0.09 \le \tilde f \le 0.17$ (again $f^{*}\approx 80$–$145~\mathrm{kHz}$). These bands target Mack-mode receptivity under the elevated post-shock pressure and temperature and are consistent with the accompanying \gls{LST} predictions. Numerically, the suction–blowing is enforced as a Dirichlet condition on the wall-normal mass flux at boundary nodes within the Gaussian support; elsewhere the no-penetration condition is applied. The construction provides a compact, well-resolved, reproducible, and inflow-independent disturbance source that excites the intended instability content while maintaining small perturbation levels.

\subsection{Pressure response handling}
The perturbation response is monitored via $500$ wall probes, equally spaced over $x_d \in [110\,\delta^{*}_{0},\,800\,\delta^{*}_{0}]$, recording wall-pressure time series at a sampling frequency $f_s$ chosen to resolve the highest forced tone (Nyquist criterion: $f_s > 2\max_i f_i$). The simulation time step provides $2000$ samples per period of the fundamental frequency. A sample of the wall-pressure evolution over a representative interval is shown in \fref{figure:fftplot}(a), where pressure traces from $10$ monitor points are plotted against time.

\begin{figure}[H]
    \centering
    \includegraphics[width=14cm]{resources/figures/disturbance/result_fft_plot_mach6_tw18.pdf}
    \caption{Wall pressure fluctuations $p_\mathrm{w}^\prime$ (Pa) over a period of time and Fourier transform of the fluctuations with a Hanning filter to give $p_\mathrm{r}^\prime$.}
    \label{figure:fftplot}
\end{figure}



After discarding the initial transient portion, the wall-pressure fluctuations are analysed using a discrete Fourier transform (DFT) to recover the imposed frequencies and their streamwise amplification in wall pressure. To reduce spectral leakage, a Hann (Hanning) window is applied to each record prior to the transform. For each probe location and forced frequency $f_i$, the complex DFT coefficient $\widehat{p}$ is extracted and its amplitude $\hat{a}_{\mathrm{r}}$ computed. With $N_s$ samples retained per probe, the record length is $T = N_s/f_s$ and the frequency resolution is $\Delta f = f_s/N_s$, ensuring that each forced frequency falls on (or within one bin of) the discrete spectrum.

The DNS-based $N$-factors are evaluated from the amplitude ratios as
\begin{equation}\label{eq:dnsnfactor_dense}
    N_{\mathrm{DNS}} = \ln\!\left(\frac{\hat{a}_{r}}{\hat{a}_{r,0}}\right) + c,
\end{equation}
where $\hat{a}_{r,0}$ is the neutral-point amplitude taken from the early-growth baseline and $c$ is a constant offset used to align the DNS-based $N$-factors with those from \gls{LST}. In continuous form, this represents the integrated amplification
\begin{equation}
    N_{\mathrm{DNS}} = \int_{x_0}^{x_d} \frac{1}{a_{\mathrm{rms}}}\,\frac{\mathrm{d}a_{\mathrm{rms}}}{\mathrm{d}x}\,\mathrm{d}x + c,
\end{equation}
with $x_0$ the neutral location. Using wall pressure as the observable provides a robust, directly accessible measure of the acoustic-type Mack-mode content and enables consistent comparison with the growth predicted by \gls{LST}.


% \subsection{Pressure response handling}
% \begin{figure}[H]
%     \centering
%     \includegraphics[width=14cm]{resources/figures/disturbance/result_fft_plot_mach6_tw18.pdf}
%     \caption{ Wall pressure fluctuations $p_\mathrm{w}^\prime \ \mathrm{(Pa)}$ over a period of time and Fourier transform of the fluctuations with a Hanning filter to give $p_\mathrm{r}^\prime$.}
%     \label{figure:fftplot}
% \end{figure}

% The perturbation response is monitored via $500$ wall probes, equally spaced over $x_d \in [\,110\,\delta^{*}_{0},\,800\,\delta^{*}_{0}\,]$, recording wall-pressure time series at a sampling frequency $f_s$ chosen to resolve the highest forced tone (Nyquist: $f_s > 2\max_i f_i$). The simulation time step provides $2000$ samples per period of the fundamental frequency. A sample of the wall-pressure evolution over a representative interval is shown in \fref{figure:fftplot}(a), where pressure traces from $10$ monitor points are plotted against time.

% After discarding the initial transient portion, the wall-pressure fluctuations are analysed using a discrete Fourier transform (DFT) to recover the imposed frequencies and their streamwise amplification in wall pressure, i.e.\ the linear response $\tilde{\boldsymbol{q}}$. To reduce spectral leakage, a Hann (Hanning) window is applied to each record prior to the transform. For each probe at $x_d$ and each forced frequency $f_i$, the complex DFT coefficient $\widehat{p}(x,f_i)$ is extracted; the corresponding amplitude is defined as $\hat a_{\mathrm{r}}(x,f_i)=\lvert \widehat{p}(x,f_i)\rvert$. Let $N_s$ denote the number of samples retained per probe; the record length and frequency resolution are then
% \begin{equation}
%   T \;=\; \frac{N_s}{f_s}, 
%   \qquad 
%   \Delta f \;=\; \frac{1}{T} \;=\; \frac{f_s}{N_s},
% \end{equation}
% which ensures that each $f_i$ falls on (or within one bin of) the discrete spectrum. The DNS-based $N$-factors are evaluated from the amplitude ratios as
% \begin{equation}\label{eq:dnsnfactor_dense}
%     N_{\mathrm{DNS}}(x_d,f_i) \;=\; \ln\!\left(\frac{\hat a_{r}(x,f_i)}{\hat a_{r,0}(f_i)}\right) \;+\; c,
% \end{equation}
% where $a_{\mathrm{rms},0}(f_i)$ is the neutral-point amplitude taken from the early-growth baseline and $c$ is a constant offset used to align the DNS-based $N$-factors with those from \gls{LST}, which does not include receptivity. For reference, the continuous form reads
% \begin{equation}
%     N_{\mathrm{DNS}}(x,f_i) \;=\; \int_{x_0}^{x_d} \frac{1}{a_{\mathrm{rms}}(x,f_i)}\,\frac{\mathrm{d}a_{\mathrm{rms}}(x,f_i)}{\mathrm{d}x}\,\mathrm{d}x \;+\; c,
% \end{equation}
% with $x_0$ the neutral location associated with $a_{\mathrm{rms},0}(f_i)$. Using wall pressure as the observable provides a robust, directly accessible measure of the acoustic-type Mack-mode content and enables consistent comparison with the growth predicted by \gls{LST}.


\subsection{Convergence studies}
To verify the linearity of the disturbance amplitudes, a systematic amplitude-scaling test was performed for two representative configurations: a Mach~6 cold-wall case in \gls{TNE} and a Mach~8 hot-wall case with \gls{TEQ}. These cases bracket the range of wall thermal conditions and vibrational models used in the study and therefore provide a stringent test of linearity. Figures~\ref{figure:linearmach6} and~\ref{figure:linearmach8} compare responses to four initial amplitudes, $\hat{a}_0 \in \{10^{-6}, 10^{-5}, 10^{-4}, 10^{-3}\}$, with each trace rescaled to the $\hat{a}_0 = 10^{-3}$ reference for visual comparison.

For both configurations, the curves for $\hat{a}_0 = 10^{-6}, 10^{-5}, 10^{-4}$ collapse to within visual tolerance, confirming linear behaviour where the wavepacket amplitude $\hat{a}_r(x)$ scales proportionally with the input whilst the peak location and local growth and decay rates remain invariant across the $\tilde{f}$ range. In contrast, the $\hat{a}_0 = 10^{-3}$ responses deviate markedly in both cases: at Mach~6, the peak fundamental amplitude is reduced and the envelope broadens, whilst at Mach~8, higher harmonics become more pronounced. These deviations indicate the onset of weak nonlinearity, consistent with quadratic self-interaction of the fundamental in strictly two-dimensional \gls{DNS}, which generates both a steady mean-flow correction ($k_x = 0$) and higher streamwise harmonics. Both channels drain energy from the fundamental, yielding a lower $\hat{a}_r$ despite the larger forcing. The nonlinear saturation effects observed at large amplitudes underscore the importance of maintaining disturbances in the linear regime for rigorous modal analysis. On the basis of these tests, $\hat{a}_0 = 10^{-5}$ was adopted for all subsequent cases: it lay comfortably above the numerical noise floor whilst remaining strictly within the linear-response range for both test configurations. The remaining cases exhibited the same qualitative trend, so $\hat{a}_0 = 10^{-5}$ was applied uniformly throughout, ensuring consistency across the entire parametric study.

\begin{figure}[ht]
\centering
\includegraphics[width=0.9\textwidth]{resources/figures/disturbance/result_mach6_linear_comp.pdf}
\caption{Amplitude-scaling at $M_e = 6$, $T_w = 1200~\mathrm{K}$ (TNE): streamwise evolution of the fundamental amplitude $\hat{a}_r(x)$ for multiple $\hat{a}_0$ across selected $\tilde{f}$.}
\label{figure:linearmach6}
\end{figure}

\begin{figure}[ht]
\centering
\includegraphics[width=0.9\textwidth]{resources/figures/disturbance/result_mach8_linear_comp.pdf}
\caption{Amplitude-scaling at $M_e = 8$, $T_w = 1800~\mathrm{K}$ (TEQ): streamwise evolution of the fundamental amplitude $\hat{a}_r(x)$ for multiple $\hat{a}_0$ across selected $\tilde{f}$.}
\label{figure:linearmach8}
\end{figure}

Furthermore, a grid study was performed on two additional representative configurations: a Mach~6 hot-wall \gls{TNE} case and a Mach~8 cold-wall \gls{TEQ} case. Two grids were compared: $2800 \times 600$ and $4000 \times 1000$ points in $(x,y)$. The solutions on both grids collapsed for all reported diagnostics, with the streamwise evolution of $\hat{a}_r(x)$ serving as the primary metric of interest. Boundary-layer integral quantities exhibited the same trend, confirming grid independence at $2800 \times 600$. Figure~\ref{figure:convergencemach6} illustrates the grid convergence for the Mach~6 hot-wall \gls{TNE} case, showing the streamwise evolution of wall pressure $\bar{p}_w$ and skin friction coefficient $c_f$. The observed convergence validates the spatial discretisation strategy employed throughout this investigation. Accordingly, $2800 \times 600$ was adopted for all subsequent simulations. The time step was adjusted with grid refinement to maintain a fixed CFL target, with $\Delta t$ reduced in proportion to the local spacing based on $\min\{\Delta x, \Delta y\}$. The same CFL-based policy was applied uniformly to all cases, ensuring temporal resolution consistency across varying mesh densities and physical regimes examined in the study.

\begin{figure}[t]
\centering
\includegraphics[width=0.9\textwidth]{resources/figures/disturbance/result_mach6_convergence_grid.pdf}
\caption{Grid convergence at $M_e = 6$, $T_w = 1800~\mathrm{K}$ (TNE): streamwise evolution of wall pressure $\hat{p}_w$ and skin friction coefficient $c_f$.}
\label{figure:convergencemach6}
\end{figure}

% \subsection{Convergence studies}
% To verify the linearity of the disturbance amplitudes, a systematic amplitude-scaling test was performed for two representative configurations: a Mach~6 cold-wall case in \gls{TNE} and a Mach~8 hot-wall case with \gls{TEQ}. These cases bracket the range of wall thermal conditions and vibrational models used in the study and therefore provide a stringent test of linearity. Figures~\ref{figure:linearmach6} and~\ref{figure:linearmach8} compare responses to four initial amplitudes, $\hat{a}_0 \in \{10^{-6}, 10^{-5}, 10^{-4}, 10^{-3}\}$, with each trace rescaled to the $\hat{a}_0 = 10^{-3}$ reference for visual comparison.

% For both configurations, the curves for $\hat{a}_0 = 10^{-6}, 10^{-5}, 10^{-4}$ collapse to within visual tolerance, confirming linear behaviour where the wavepacket amplitude $\hat{a}_r(x)$ scales proportionally with the input whilst the peak location and local growth and decay rates remain invariant across the $\tilde{f}$ range. In contrast, the $\hat{a}_0 = 10^{-3}$ responses deviate markedly in both cases: at Mach~6, the peak fundamental amplitude is reduced and the envelope broadens, whilst at Mach~8, higher harmonics become more pronounced. These deviations indicate the onset of weak nonlinearity, consistent with quadratic self-interaction of the fundamental in strictly two-dimensional \gls{DNS}, which generates both a steady mean-flow correction ($k_x = 0$) and higher streamwise harmonics. Both channels drain energy from the fundamental, yielding a lower $\hat{a}_r$ despite the larger forcing. On the basis of these tests, $\hat{a}_0 = 10^{-5}$ was adopted for all subsequent cases: it lay comfortably above the numerical noise floor whilst remaining strictly within the linear-response range for both test configurations. The remaining cases exhibited the same qualitative trend, so $\hat{a}_0 = 10^{-5}$ was applied uniformly throughout.

% \begin{figure}[ht]
% \centering
% \includegraphics[width=0.9\textwidth]{resources/figures/disturbance/result_mach6_linear_comp.pdf}
% \caption{Amplitude-scaling at $M_e = 6$, $T_w = 1200~\mathrm{K}$ (TNE): streamwise evolution of the fundamental amplitude $\hat{a}_r(x)$ for multiple $\hat{a}_0$ across selected $\tilde{f}$.}
% \label{figure:linearmach6}
% \end{figure}

% \begin{figure}[ht]
% \centering
% \includegraphics[width=0.9\textwidth]{resources/figures/disturbance/result_mach8_linear_comp.pdf}
% \caption{Amplitude-scaling at $M_e = 8$, $T_w = 1800~\mathrm{K}$ (TEQ): streamwise evolution of the fundamental amplitude $\hat{a}_r(x)$ for multiple $\hat{a}_0$ across selected $\tilde{f}$.}
% \label{figure:linearmach8}
% \end{figure}

% Furthermore, a grid study was performed on two additional representative configurations: a Mach~6 hot-wall \gls{TNE} case and a Mach~8 cold-wall \gls{TEQ} case. Two grids were compared: $2800 \times 600$ and $4000 \times 1000$ points in $(x,y)$. The solutions on both grids collapsed for all reported diagnostics, with the streamwise evolution of $\hat{a}_r(x)$ serving as the primary metric of interest. Boundary-layer integral quantities exhibited the same trend, confirming grid independence at $2800 \times 600$. Figure~\ref{figure:convergencemach6} illustrates the grid convergence for the Mach~6 hot-wall \gls{TNE} case, showing the streamwise evolution of wall pressure $\bar{p}_w$ and skin friction coefficient $c_f$. Accordingly, $2800 \times 600$ was adopted for all subsequent simulations. The time step was adjusted with grid refinement to maintain a fixed CFL target, with $\Delta t$ reduced in proportion to the local spacing based on $\min\{\Delta x, \Delta y\}$. The same CFL-based policy was applied uniformly to all cases.

% \begin{figure}[ht]
% \centering
% \includegraphics[width=0.9\textwidth]{resources/figures/disturbance/result_mach6_convergence_grid.pdf}
% \caption{Grid convergence at $M_e = 6$, $T_w = 1800~\mathrm{K}$ (TNE): streamwise evolution of wall pressure $\hat{p}_w$ and skin friction coefficient $c_f$.}
% \label{figure:convergencemach6}
% \end{figure}


% To verify the linearity of the disturbance amplitudes, a systematic amplitude–scaling test was performed for two representative configurations: a Mach~6 \emph{cold-wall} case in \gls{TNE} and a Mach~8 \emph{hot-wall} case with \gls{TEQ}. These bracket the range of wall thermal conditions and vibrational models used elsewhere in the study and therefore provide a stringent test of linearity. Figures~\ref{figure:linearmach6},~\ref{figure:linearmach8} compare responses to four initial amplitudes, $\hat a_0\in\{10^{-6},10^{-5},10^{-4},10^{-3}\}$. For visual comparison, each trace is rescaled to the $\hat a_0=10^{-3}$ reference. For both configurations, the curves for $\hat a_0=10^{-6},10^{-5},10^{-4}$ collapse to within visual tolerance, indicating linear behaviour, where the wavepacket amplitude $\hat a_r(x)$ scales proportionally with the input, while the peak location and the local growth and decay rates remain invariant across $\tilde f$ range. In contrast, the $\hat a_0=10^{-3}$ responses deviate markedly in both cases: the peak fundamental amplitude is reduced and the envelope broadens, indicating the onset of weak nonlinearity. In a strictly 2D DNS this is consistent with quadratic self–interaction of the fundamental, which generates (i) a steady mean–flow correction ($k_x=0$) and (ii) higher streamwise harmonics. Both channels drain energy from the fundamental, yielding a lower $\hat a_r$ despite the larger forcing. On the basis of these tests, $\hat a_0=10^{-5}$ was adopted for all production runs: it lay comfortably above the numerical noise floor while remaining strictly within the linear–response range for the Mach~6 cold–wall \gls{TNE} and Mach~8 hot–wall \gls{TEQ} cases. The remaining configurations (other wall temperatures and the \gls{CPG}/\gls{TEQ}/\gls{TNE} variants at Mach~6 and Mach~8) exhibited the same qualitative trend, so the choice $\hat a_0=10^{-5}$ was applied uniformly throughout.

% \begin{figure}[ht]
% \centering
% \includegraphics[width=0.9\textwidth]{resources/figures/disturbance/result_mach6_linear_comp.pdf}
% \caption{Amplitude–scaling at $M_e=6$, $T_w=1200~\mathrm{K}$ (TNE): streamwise evolution of the fundamental amplitude $\hat a_r(x)$ for multiple $\hat a_0$ across selected $\tilde f$.}\label{figure:linearmach6}
% \end{figure}

% \begin{figure}[ht]
% \centering
% \includegraphics[width=0.9\textwidth]{resources/figures/disturbance/result_mach8_linear_comp.pdf}
% \caption{Amplitude–scaling at $M_e=8$, $T_w=1800~\mathrm{K}$ (TEQ): streamwise evolution of the fundamental amplitude $\hat a_r(x)$ for multiple $\hat a_0$ across selected $\tilde f$.}\label{figure:linearmach8}
% \end{figure}

% Furthermore, a grid study was performed on two additional representative configurations: a Mach~6 \emph{hot–wall} \gls{TNE} case and a Mach~8 \emph{cold–wall} \gls{TEQ} case. Three grids were considered, $2000\times500$, $2800\times600$, and $4000\times1000$ points in $(x,y)$. The solutions on $2800\times600$ and $4000\times1000$ collapsed for all reported diagnostics; the streamwise evolution of $\hat a_r(x)$ was the primary metric of interest. Although not shown, boundary–layer integral quantities exhibited the same trend, confirming grid independence at $2800\times600$. Accordingly, $2800\times600$ was adopted for subsequent simulations. The time step was adjusted with grid refinement to maintain a fixed CFL target, i.e.\ $\Delta t$ was reduced in proportion to the local spacing (based on $\min\{\Delta x,\Delta y\}$), and the same CFL-based policy was applied uniformly to the remaining cases.


% \begin{figure}[ht]
% \centering
% \includegraphics[width=0.9\textwidth]{resources/figures/disturbance/result_mach6_convergence_grid.pdf}
% \caption{Amplitude–scaling at $M_e=6$, $T_w=1800~\mathrm{K}$ (TNE): streamwise evolution of the wall pressure, $\hat p_w$ and skin friction $c_f$.}\label{figure:convergencemach6}
% \end{figure}


% \subsection{Wall treatment}
% To introduce controlled, small-amplitude perturbations, we employ a wall suction–blowing forcing strip in the spirit of \cite{AlamSandham2000}, adapted to our two-dimensional (spanwise-uniform) domain. A Gaussian-localised disturbance is centred at $x_f$ (see \cref{figure:schematicflow}) and imposed as a wall mass-flux perturbation,
% \begin{equation}
%   \rho v \;=\; a_0\, g(x)\, \sum_{i=1}^{8} \sin\!\big(2\pi f_i t\big),
% \end{equation}
% with $a_0 = 1\times10^{-5}\,u_e$ the prescribed forcing amplitude, chosen to keep the response in the linear regime. The streamwise shape function is
% \begin{equation}
%   g(x) \;=\; \exp\!\left[-\,\frac{(x-x_f)^2}{2\sigma^2}\right],
% \end{equation}
% where $2\sigma^2 = 20\,(\delta_0^{*})^2$ and $x_f = 25\,\delta_0^{*}$. This placement and width ensure adequate grid resolution within the forcing envelope, keep the excitation spatially compact, and minimise interactions with the inflow and outflow boundaries.

% The non-dimensional frequency is defined as $\tilde f = f\,\delta_0/u_e$, and the discrete tones are selected so that the resulting wavepackets amplify within $x/\delta_0^{*}\in[0,\,750]$. 
% For Mach~6 we employ eight tones spanning $0.10 \le \tilde f \le 0.18$ (corresponding to $f^{*}\approx 80$–$145~\mathrm{kHz}$); for Mach~8, $0.09 \le \tilde f \le 0.17$ (again $f^{*}\approx 80$–$145~\mathrm{kHz}$). These bands target Mack-mode receptivity under the elevated post-shock pressure and temperature and are consistent with the accompanying \gls{LST} predictions. The construction provides a compact, well-resolved, reproducible, and inflow-independent disturbance source that excites the intended instability content while maintaining small perturbation levels.

% \subsection{Pressure response handling}
% The response to the initial perturbation was monitored through $400$ wall probes, equally spaced over $x_d \in [\,300\,\delta^{*}_{0},\,600\,\delta^{*}_{0}\,]$, recording wall-pressure time series at a sampling frequency $f_s$ (chosen to resolve the highest forced tone). The simulation time step provided $2000$ samples per period of the fundamental frequency. A sample of the wall-pressure evolution over a representative time interval is shown in \fref{figure:fftplot}(a), where pressure traces from $10$ monitor points are plotted against time.

% After discarding the initial transient portion, the wall-pressure fluctuation data were analysed using a discrete Fourier transform (DFT) to recover the imposed frequencies and their streamwise amplification in wall pressure, i.e.\ the linear response $\tilde{\boldsymbol{q}}$. To reduce spectral leakage, a Hann (Hanning) window was applied to each temporal record prior to the transform. For each probe location $x_d$ and forced frequency $f_i$, the peak of the DFT magnitude at $f_i$ was extracted to define the amplitude $a_{\mathrm{rms}}(x_d,f_i)$, which is used for amplitude comparisons and for the $N$-factor evaluation. The DNS-based $N$-factors are then taken as
% \begin{equation}
%     N_{\mathrm{DNS}}(x_d,f_i) \;=\; \ln\!\left(\frac{a_{\mathrm{rms}}(x_d,f_i)}{a_{\mathrm{rms},0}(f_i)}\right) \;+\; c,
% \end{equation}
% where $a_{\mathrm{rms},0}(f_i)$ is the neutral-point amplitude (taken from the early-growth baseline) and $c$ is a constant offset used to align the DNS-based $N$-factors with those from \gls{LST}, which does not include receptivity.

% To introduce controlled, small-amplitude perturbations, we employ a wall suction–blowing forcing strip in the spirit of \cite{AlamSandham2000}, adapted to our two-dimensional domain. A Gaussian-localised disturbance is centred at $x_f$ (see \cref{figure:schematicflow}) and imposed as a wall mass-flux perturbation,
% \begin{equation}
%   \rho v \;=\; a_0\, g(x)\, \sum_i \sin\!\big(2\pi f_i t\big),
% \end{equation}
% with $a_0 = 1\times10^{-5}\,u_e$ is the amplitude parameter of the forcing chosen to be small ensuring the disturbances are in the lienar region. The streamwise shape function is
% \begin{equation}
%   g(x) \;=\; \exp\!\left[-\,\frac{(x-x_f)^2}{2\sigma^2}\right],
% \end{equation}
% where $2\sigma^2 = 20\,(\delta_0^{*})^2$ and $x_f = 25\,\delta_0^{*}$. This placement and width ensure adequate grid resolution within the forcing envelope and prevent boundary interactions, while keeping the excitation spatially compact.

% The non-dimensional frequency is defined as $\tilde f = f\,\delta_0/u_e$; frequencies are selected so that the resulting wavepackets amplify within $x/\delta_0^{*}\in[0,\,750]$. For Mach~6 we use eight tones spanning $0.10 \le \tilde f \le 0.18$ (corresponding to $f^{*}\approx 80$–$145~\mathrm{kHz}$); for Mach~8, $0.09 \le \tilde f \le 0.17$ (again $f^{*}\approx 80$–$145~\mathrm{kHz}$). These bands target Mack-mode receptivity under the elevated post-shock pressure and temperature and are consistent with the accompanying \gls{LST} predictions. The construction therefore provides a compact, well-resolved, and reproducible disturbance source that excites the intended instability content while maintaining small perturbation levels.


% To induce small amplitude perturbations in the domain, a suction and blowing type forcing strip similar to \cite{alam_sandham_2000} is employed, updated to account for small amplitude perturbations and the two-dimensional domain. A Gaussian-shaped disturbance is placed at $x_f$ as shown in ~\cref{figure:schematicflow} and is simulated as a perturbation of mass flux on the wall \begin{equation}\label{equation:forcing}
%     \rho v = a_0 \  g(x) \ \sum_i \sin{ (2\pi f_i t)}
% \end{equation} where, $a_0$ = $1\times10^{-5} \ u_{e}$ is the amplitude parameter of the forcing chosen to be small ensuring the disturbances are in the linear region. The profile shape function, $g(x)$ is  \begin{equation}
%     g(x) = \exp{- \frac{(x- x_f)}{2 \sigma^2}}
% \end{equation} with the width defined by $2\sigma^2 = 20 (\delta_{0}^{*})^2$ and $x_f = 25 \delta_{0}^{*}$ from the origin of the simulation domain.  This was selected to ensure there were enough points within the shape function to accurately represent the disturbance. The chosen amplitude for the simulation is $a_0 = 1\times 10^{-5}$, which was taken from a linearity study where $a_0$ was varied between $a_0 = 1\times 10^{-3}$ to $a_0 = 1\times 10^{-9}$. While the highest amplitude case showed nonlinear behaviour, it was observed that amplitudes below $a_0 = 1\times 10^{-4}$ could be classed as having a linear response. The non-dimensional frequency $\tilde f$ is defined as $\tilde f= f \delta_{0}/u_e$, and the frequencies are selected to ensure that the response occurs within the streamwise domain of $x = 0 - 750\delta^{*}_{0}$. A set of 8 non-dimensional frequencies are included in the trip equation for each case, ranging over $0.1 \le \tilde f \le 0.18$ which corresponds to frequencies in the range of $f^{*} = 80-145 kHz$ for Mach 6 flows and $0.09 \le \tilde f \le 0.17$ which are in the range of $f^{*} = 80-145 kHz$ for the Mach 8 cases.  The chosen frequencies are within the region for Mack modes at the heightened free-stream pressure and temperature and are corroborated by the \gls{LST} studies. 




% \newpage


\section{Stability characteristics}
% A total of twelve two-dimensional cases were simulated, covering all combinations of \gls{CPG}, \gls{TEQ}, and \gls{TNE} at Mach~6 and Mach~8, with wall temperatures $T_w=1200\,\mathrm{K}$ and $T_w=1800\,\mathrm{K}$ (hereafter, \emph{cold wall} and \emph{hot wall}, respectively). A number of systematic surveys across Mach number and wall temperature exist, such as those by \citet{Mack1969,Mack1984,MasadEtAl1992,Malik1989}, but these were conducted at low enthalpy. By contrast, \citet{Bitter2015} examined high-enthalpy boundary layers and documented high-enthalpy effects on instability; however, the effect of thermally excited flows was not investigated in a high-fidelity setting. The present cases address this gap and are directly relevant to real flight conditions for a \gls{HGV} shown in \cref{figure:trajectory}. In this section, comparisons between \gls{LST} and \gls{DNS} are presented first; the base flows are then compared across the different thermal assumptions; subsequently, the growth of fluctuations and effects of thermal excitation across Mach numbers are examined.
A total of twelve two-dimensional cases were simulated, covering all combinations of \gls{CPG}, \gls{TEQ}, and \gls{TNE} at Mach~6 and Mach~8, with wall temperatures $T_w=1200\,\mathrm{K}$ and $T_w=1800\,\mathrm{K}$ (hereafter, \emph{cold wall} and \emph{hot wall}, respectively). Classical stability surveys across varying Mach numbers and wall conditions have been established by prior work \citep{Mack1969,Mack1984,Malik1989,MasadEtAl1992}; however, these were predominantly conducted at low enthalpy. More recent investigations at elevated enthalpies have identified novel high-temperature effects on modal instability \citep{Bitter2015}, yet the interplay between thermochemical nonequilibrium and instability dynamics remains underexplored in high-fidelity simulations. The present work addresses this gap by examining thermochemically consistent base flows under conditions representative of actual hypersonic vehicle trajectories, as depicted in \cref{figure:trajectory}. In this section, comparisons between \gls{LST} and \gls{DNS} are presented first; the base flows are then compared across the different thermal assumptions; subsequently, the growth of fluctuations and effects of thermal excitation across Mach numbers are examined.

\subsection{Base flows}

\begin{figure}[t]
\centering
\includegraphics[width=0.9\textwidth]{resources/figures/disturbance/result_profiles_m6tw12.pdf}
\caption{Mach 6, cold wall $Tw=1200 \ K$ - normalised streamwise velocity $u/u_e$, translational-rotational temperature $T/T_e$ and vibrational temperature $T_v/T_e$ profiles as a function of the similarity variable $\eta = y/\delta^{*}$ compared against respective locally self-similar solutions. Profiles were taken at 75\% of the domain length.}\label{figure:bcprofiles}
\end{figure}


A comparison of two-dimensional \gls{DNS} boundary-layer profiles with their locally self-similar counterparts is shown in \cref{figure:bcprofiles} for velocity (top row), temperature (middle row), and vibrational temperature (bottom row). The left column corresponds to \gls{CPG}, the middle column to \gls{TEQ}, and the right column to \gls{TNE}, for which the reference is a \gls{TFG} similarity solution. The profiles shown are for the Mach~6 cold-wall condition, extracted at approximately $75\%$ of the streamwise domain length to ensure fully developed upstream histories whilst remaining well upstream of the outflow. For both \gls{CPG} and \gls{TEQ}, the agreement with the corresponding similarity profiles is excellent across the wall-normal extent, including the near-wall gradients and outer-layer approach to the free stream, indicating that the numerical solution reproduces the expected self-similar structure. In \gls{TEQ}, the temperature peak is reduced by approximately $8\%$ relative to \gls{CPG} because a portion of the thermal energy is accommodated in vibrational modes, draining the total peak heating; here $T_v=T$ by assumption, whereas vibrational excitation is not present in \gls{CPG}. The \gls{TNE} simulations are initialised from the \gls{TFG} similarity state, and the subsequent relaxation $T_v\!\to\!T$ is clearly evident but incomplete at the station plotted: $T_v$ remains elevated, so additional energy resides in vibration relative to the \gls{TFG} baseline and the translation–rotation temperature is correspondingly depressed. Overall, the \gls{CPG} and \gls{TEQ} cases validate the local self-similarity framework, whilst the \gls{TNE} profiles capture the expected nonequilibrium adjustment from \gls{TFG} initialisation towards $T_v=T$, with the residual offset quantifying the remaining relaxation distance. The profiles of the other flight conditions show similar traits and are therefore omitted for brevity. 


% Figure~\ref{figure:bcprofiles} shows the evolution of displacement thickness $\delta^{*}$, momentum thickness $\theta^{*}$, energy thickness $\delta^{*}_{e}$, and enthalpy thickness $\delta^{*}_{h}$ for the Mach~6 cold–wall case.  The (compressible, density–weighted) definitions are
% \begin{equation}
% \delta^{*}=\int_{0}^{\infty}\left(1-\frac{\rho u}{\rho_{e} u_{e}}\right)\,dy,\qquad
% \theta^{*}=\int_{0}^{\infty}\frac{\rho u}{\rho_{e} u_{e}}\left(1-\frac{u}{u_{e}}\right)\,dy
% \end{equation}
% \begin{equation}
% \delta^{*}_{e}=\int_{0}^{\infty}\frac{\rho u}{\rho_{e} u_{e}}
% \left(1-\frac{u^{2}}{u_{e}^{2}}\right)\,dy,\qquad
% \delta^{*}_{h}=\int_{0}^{\infty}\frac{\rho u}{\rho_{e} u_{e}}
% \frac{T-T_{e}}{T_{w}-T_{e}}\,dy.
% \end{equation}
% \noindent where subscripts $e$ and $w$ denote edge and wall values, respectively. 

% Figure~\ref{figure:bcprofiles} shows the streamwise evolution of displacement thickness $\delta^{*}$ and momentum thickness $\theta$ for the Mach~6 and Mach~8 cold-wall cases. The compressible, density-weighted definitions are
% \begin{equation}
% \delta^{*}=\int_{0}^{\infty}\left(1-\frac{\rho u}{\rho_{e} u_{e}}\right)\,dy,\qquad
% \theta=\int_{0}^{\infty}\frac{\rho u}{\rho_{e} u_{e}}\left(1-\frac{u}{u_{e}}\right)\,dy,
% \end{equation}
% where subscripts $e$ and $w$ denote edge and wall values, respectively. 

% \begin{figure}[t]
% \centering
% \includegraphics[width=0.96\textwidth]{resources/figures/disturbance/result_boundarylayergrowth.pdf}
% \caption{Streamwise evolution of boundary-layer integral parameters for CPG, TEQ, TNE simulations for Mach 6 hot wall case.}\label{figure:tvprofiles}
% \end{figure}

% Both thicknesses increase monotonically as the boundary layer develops, but the models separate in a physically consistent way. For cold walls, \gls{TEQ} drives stronger near-wall cooling than \gls{CPG}, packing more mass flux into the low-speed region. This reduces the mass-flux deficit, yielding a smaller $\delta^{*}$ in \gls{TEQ}. Concurrently, temperature-dependent transport steepens the wall gradient, raising $c_f$ and producing a larger momentum thickness $\theta$ via the von Kármán relation $d\theta/dx\approx c_f/2$. \gls{TNE} sits between \gls{CPG} and \gls{TEQ} because vibrational relaxation damps temperature excursions and weakens near-wall cooling. This yields $\delta^{*}$ closer to \gls{CPG} and $\theta$ closer to \gls{TEQ}. These trends are consistent across both Mach numbers, though absolute magnitudes differ due to stronger compression heating at Mach~8.


% \begin{figure}[t]
% \centering
% \includegraphics[width=0.96\textwidth]{resources/figures/disturbance/result_boundarylayergrowth_mach.pdf}
% \caption{Streamwise evolution of boundary-layer integral parameters for TNE simulations at different Mach numbers and wall temperatures: (a) displacement thickness $\delta^*$, (b) momentum thickness $\theta^*$, (c) incompressible shape factor $H_i$, and (d) boundary-layer thickness $\delta_{99}$.}
% \label{figure:bcgrowth_mach}
% \end{figure}


% All thicknesses increase monotonically with $x/\delta^{*}_{0}$ as the layer develops, but the models separate in a physically consistent way. 
% For a cold wall, \gls{TEQ} drives stronger near–wall cooling than \gls{CPG}; the higher density packs more mass flux into the low–speed region, which increases $\rho u$ where $u/u_e\ll1$ and therefore reduces the integrand $1-\rho u/(\rho_e u_e)$. The net deficit shrinks, so $\delta^{*}$ grows slightly \emph{less} in \gls{TEQ} than in \gls{CPG}. 
% Concurrently, temperature–dependent transport steepens the wall gradient (\(\mu\) lower near the cold wall, higher aloft), raising $c_f$; with $dp/dx\simeq0$, the von Kármán relation gives $d\theta/dx\approx c_f/2$, hence a \emph{larger} momentum thickness $\theta^{*}$ in \gls{TEQ}. 
% The energy thickness $\delta^{*}_{e}$ follows the same logic its kernel $1-(u/u_e)^2$ weights the momentum deficit even more so the fuller inner profile and larger wall shear yield a slightly larger $\delta^{*}_{e}$ than in \gls{CPG}. 
% By contrast, the enthalpy thickness $\delta^{*}_{h}$ tracks the thermal field: stronger cooling in \gls{TEQ} compresses the thermal layer and reduces the integrated enthalpy deficit, giving a smaller $\delta^{*}_{h}$ than \gls{CPG}. \gls{TNE} sits between \gls{CPG} and \gls{TEQ} because vibrational relaxation alters only the \emph{thermal} response. With chemistry frozen, changes in $\rho$ and $u$ arise solely from how energy is partitioned between translational–rotational and vibrational modes. Where $T$ and $T_v$ decouple, the effective heat capacity increases; this damps $T$ excursions, weakens the near–wall cooling relative to \gls{TEQ}, and slightly reduces $c_f$. The result is $\delta^{*}$ and $\delta^{*}_{h}$ closer to \gls{CPG}, while $\theta_\delta$ and $\delta^{*}_{e}$ remain closer to (but typically below) \gls{TEQ}. The small ripples on some curves are numerical/forcing artefacts and do not affect these trends.


Figure~\ref{figure:bcprofiles} shows the streamwise evolution of displacement thickness $\delta^{*}$ and momentum thickness $\theta$ for the Mach~6 and Mach~8 cold-wall cases. The compressible, density-weighted definitions are
\begin{equation}
\delta^{*}=\int_{0}^{\infty}\left(1-\frac{\rho u}{\rho_{e} u_{e}}\right)\,dy,\qquad
\theta^{*}=\int_{0}^{\infty}\frac{\rho u}{\rho_{e} u_{e}}\left(1-\frac{u}{u_{e}}\right)\,dy,
\end{equation}
where subscripts $e$ and $w$ denote edge and wall values, respectively. 

% \begin{figure}[t]
% \centering
% \includegraphics[width=0.96\textwidth]{resources/figures/disturbance/result_boundarylayergrowth.pdf}
% \caption{Streamwise evolution of boundary-layer integral parameters for CPG, TEQ, and TNE simulations: Mach~6 and Mach~8 cold-wall cases.}
% \label{figure:bcprofiles}
% \end{figure}

% Both thicknesses increase monotonically as the boundary layer develops, but the models separate in a physically consistent way. For cold walls, \gls{TEQ} drives stronger near-wall cooling than \gls{CPG}, packing more mass flux into the low-speed region. This reduces the mass-flux deficit, yielding a smaller $\delta^{*}$ in \gls{TEQ}. Concurrently, temperature-dependent transport steepens the wall gradient, raising $c_f$ and producing a larger momentum thickness $\theta$ via the von Kármán relation $d\theta^{*}/dx\approx c_f/2$. \gls{TNE} sits between \gls{CPG} and \gls{TEQ} because vibrational relaxation damps temperature excursions and weakens near-wall cooling. This yields $\delta^{*}$ closer to \gls{CPG} and $\theta^{*}$ closer to \gls{TEQ}. These trends are consistent across both Mach numbers, though absolute magnitudes differ due to stronger compression heating at Mach~8.

\begin{figure}[t]
\centering
\includegraphics[width=0.96\textwidth]{resources/figures/disturbance/result_boundarylayergrowth.pdf}
\caption{Streamwise evolution of displacement thickness $\delta^*$ and momentum thickness $\theta$ for CPG, TEQ, and TNE simulations: (a) Mach~6 cold-wall and (b) Mach~8 cold-wall cases.}
\label{figure:blgrowth}
\end{figure}

Both thicknesses increase monotonically as the boundary layer develops, but the models separate in a physically consistent way. For cold walls, \gls{TEQ} drives stronger near-wall cooling than \gls{CPG}, packing more mass flux into the low-speed region. This reduces the mass-flux deficit, yielding a smaller $\delta^{*}$ in \gls{TEQ}. Concurrently, temperature-dependent transport steepens the wall gradient, raising $c_f$ and producing a larger momentum thickness $\theta^{*}$ via the von Kármán relation $d\theta^{*}/dx\approx c_f/2$. \gls{TNE} sits between \gls{CPG} and \gls{TEQ} because vibrational relaxation damps temperature excursions and weakens near-wall cooling, yielding $\delta^{*}$ closer to \gls{CPG} and $\theta^{*}$ closer to \gls{TEQ}. At Mach~6, the three models show distinct separation in both quantities, with the spread more pronounced in $\delta^*$ than in $\theta^{*}$. At Mach~8, however, the models nearly collapse, particularly for $\theta^{*}$, indicating that the stronger compression heating and thinner boundary layer reduce the relative influence of thermal nonequilibrium on momentum transport.

\begin{figure}[t]
\centering
\includegraphics[width=0.96\textwidth]{resources/figures/disturbance/result_boundarylayergrowth_mach.pdf}
\caption{Streamwise evolution of boundary-layer integral parameters for TNE simulations at different Mach numbers and wall temperatures: (a) displacement thickness $\delta^*$, (b) momentum thickness $\theta^*$, (c) energy thickness $\delta_e$, and (d) enthalpy thickness $\delta_h$.}
\label{figure:bcgrowth_mach}
\end{figure}

Isolating \gls{TNE} cases, Figure~\ref{figure:bcgrowth_mach} presents the streamwise evolution of all four integral thicknesses across different Mach numbers and wall temperatures. The additional quantities are defined as
\begin{equation}
\delta_{e}=\int_{0}^{\infty}\frac{\rho u}{\rho_{e} u_{e}}
\left(1-\frac{u^{2}}{u_{e}^{2}}\right)\,dy,\qquad
\delta_{h}=\int_{0}^{\infty}\frac{\rho u}{\rho_{e} u_{e}}
\frac{T-T_{e}}{T_{w}-T_{e}}\,dy.
\end{equation}
The displacement thickness $\delta^*$ exhibits higher growth rates at both higher wall temperatures and higher Mach numbers, with the hot-wall cases consistently producing thicker displacement layers. The momentum thickness $\theta^*$ shows markedly different sensitivity: at Mach~6, the cold- and hot-wall cases remain close throughout the domain, whilst at Mach~8, a substantial separation emerges, with the hot-wall case growing more rapidly. The energy thickness $\delta_e$ follows a similar trend to $\theta^*$, with its kernel $1-(u/u_e)^2$ weighting the momentum deficit more strongly, Mach~6 cases cluster together, whilst Mach~8 cases diverge significantly with wall temperature. The enthalpy thickness $\delta_h$ reveals distinct thermal behaviour: the Mach~8 cold-wall case starts with a markedly lower initial $\delta_h$ than the Mach~6 hot-wall case, reflecting stronger compression heating and a thinner thermal layer. Despite this lower initial value, the growth rate $d\delta_h/dx$ is substantially higher for the Mach~8 cold-wall condition, indicating more rapid thermal boundary-layer development driven by the steeper enthalpy gradient at higher edge enthalpies.


\subsubsection{Thermal relaxation}\label{section:disturbance:thermal_relaxation}
\begin{figure}[ht]
\centering
\includegraphics[width=0.96\textwidth]{resources/figures/disturbance/result_Tvcomp_ssprofiles_dns_update.pdf}
\caption{Comparison of $T_v$ profiles at 75\% of the domain length (a) $M_e=6$, $T_w=1200$, (b) $M_e=6$, $T_w=1200$, (c) $M_e=6$, $T_w=1200$ and (d) $M_e=8$, $T_w=1800$.}\label{figure:tvprofiles}
\end{figure}

% Figure~\ref{figure:tvprofiles} compares the $\theta_v$ and $\theta$ profiles for all four flow conditions, in each case showing the TNE case from DNS compared with the TFG similarity solution, which also constituted the inflow condition for the simulations. The cold wall cases in panels (a) and (c) show larger departures from the inflow similarity solution in both profiles compared to the hot wall cases (b) and (d) indicating faster relaxation towards equilibrium. Similarly, the lower Mach number cases (a) and (b) show larger departures than the high Mach number cases (c) and (d). As the flow goes through the domain, both $\theta_v$ and $theta$ profiles relax towards the TEQ solution. In each case, the DNS solution remains closer to the TFG than the TEQ similarity solution. Hence for the flight and wedge angle conditions chosen for this study, the TFG solution is preferred as the inflow condition for DNS, allowing the receptivity and growth of disturbances to occur in a region of space where the flow is relaxing slowly towards equilibrium.

Figure~\ref{figure:tvprofiles} compares the $\theta_v$ and $\theta$ profiles for all four flow conditions, in each case showing the \gls{TNE} result from \gls{DNS} against the \gls{TFG} similarity solution, which also provides the inflow condition for the simulations. The cold-wall cases in panels (a) and (c) show larger departures from the inflow similarity solution in both profiles than the hot-wall cases (b) and (d), indicating that lower wall temperatures promote equilibration. Similarly, the lower Mach number cases (a) and (b) depart more than the higher-Mach-number cases (c) and (d). As the flow progresses downstream, both $\theta_v$ and $\theta$ relax toward the \gls{TEQ} solution. The DNS \gls{TNE} curves sit between the two similarity predictions: $\theta_{\text{DNS}}$ lies between the self-similar \gls{TFG} and \gls{TEQ} $\theta$ curves, and likewise for $\theta_v$, indicating partial but incomplete relaxation. In each case, the \gls{DNS} solution remains closer to \gls{TFG} than to the \gls{TEQ} similarity solution. Hence, for the flight and wedge-angle conditions considered, the \gls{TFG} solution is preferred as the \gls{DNS} inflow: it supplies a self-similar, boundary-layer-consistent state at a fractional computational cost, whilst preserving slow vibrational relaxation so that receptivity and linear growth develop inside the domain rather than being set by the inflow; it also reduces sensitivity to inflow length and cleanly isolates the downstream instability physics.



\begin{figure}[t]
\centering
\includegraphics[width=0.96\textwidth]{resources/figures/disturbance/result_disturabnce_DNS_Dav.pdf}
\caption{Species-resolved vibrational Damköhler profiles against $y$, $Da_{v,s}$, for $s\in\{\mathrm{N_2},\mathrm{O_2},\mathrm{NO}\}$, sampled at multiple streamwise stations to illustrate vibrational relaxation. Panels: (a) $M_e=6$, $T_w=1200~\mathrm{K}$; (b) $M_e=6$, $T_w=1800~\mathrm{K}$; (c) $M_e=8$, $T_w=1200~\mathrm{K}$; (d) $M_e=8$, $T_w=1800~\mathrm{K}$. }
\label{figure:Davprofiles}
\end{figure}

% A complementary view is through a vibrational Damköhler number that compares
% advective residence time to vibrational relaxation:
% \begin{equation}
%   \mathrm{Da}_v \;=\; \frac{\tau_{\mathrm{flow}}}{\tau_v},
% \end{equation}
% with the flow time scale taken as
% \begin{equation}
%   \tau_{\mathrm{flow}} \sim \frac{L}{u_e},
% \end{equation}
% where \(L\) is a boundary-layer length scale chosen to be the displacement thickness and
% \(\tau_v\) is the vibrational relaxation time. The evolution of species-resolved vibrational Damköhler numbers, $Da_{v,s}$, across streamwise stations is shown in \cref{figure:Davprofiles}. Cold–wall cases show lower vibrational Damköhler numbers because cooling slows the physics and speeds the flow-through of the layer. Lower temperatures markedly lengthen vibrational relaxation, so vibrations equilibrate more slowly. At the same time, the colder wall compresses the thermal layer, reducing the residence time within the region where vibrational exchange is active. Together, slower relaxation but faster transit, drive $Da_{v,s}$ down across the cold–wall panels, with the exact drop set by each species’ vibrational characteristics. It is also evident that $Da_{v,\mathrm{N_2}}$ is consistently higher than $Da_{v,\mathrm{O_2}}$ and $Da_{v,\mathrm{NO}}$.
% Since $\mathrm{N_2}$ has a higher characteristic vibrational temperature $\Theta_v$ and therefore relaxes more slowly; $\mathrm{O_2}$ and $\mathrm{NO}$ relax faster than $\mathrm{N_2}$, yielding lower $Da_{v,s}$.  From $x/\delta_0^{*}=200$ to $800$, the $Da_{v,s}$ profiles amplify monotonically while retaining a similar shape. The peak magnitude increases at each downstream station, the peak location in $y/\delta_0^{*}$ remains approximately fixed, and the outer decay becomes more gradual as the layer thickens. The largest fractional rise of the peak $Da_{v,s}$ in panel (d) ($\approx70\%$) compared with panel (a) ($\approx40\%$) means that, downstream, the ratio of time scales is changing faster in (d): the vibrational relaxation time is dropping more rapidly with increasing temperature, and/or the advective residence time is growing more quickly as the layer thickens. In the hot–wall, higher–Mach case (d) both effects are stronger, so $Da_{v,s}$ becomes relatively larger with $x$; in the cold–wall case (a) these trends are weaker, giving a smaller increase. Therefore, higher Mach numbers and wall temperatures drive the flow toward vibrational equilibrium more rapidly: both increase $Da_{v,s}$ downstream, so relaxation outpaces advection and the vibration field approaches its local equilibrium sooner. It is important to characterise the flow and flight conditions to understand how disturbances grow within it.


A complementary characterisation employs the vibrational Damköhler number, which compares advective residence time with vibrational relaxation,
\begin{equation}
    \mathrm{Da}_v = \tau_{\mathrm{flow}}/\tau_v
\end{equation}
where the flow time scale is taken as \(\tau_{\mathrm{flow}} \sim L/u_e\), with \(L\) chosen as the displacement thickness, and \(\tau_v\) the vibrational relaxation time. The spanwise evolution of the species-resolved measures \(Da_{v,s}\) at different streamwise stations is presented in \cref{figure:Davprofiles}.

In configurations with a cooled wall, \(Da_{v,s}\) is systematically larger: reduced temperatures lengthen \(\tau_v\) (slower relaxation) but also compress the thermal layer, thereby shortening \(\tau_{\mathrm{flow}}\) (faster transit). The competing effects of slower relaxation and faster transit yield a net increase in \(Da_{v,s}\), indicating that the flow resides in the boundary layer for a sufficiently extended period relative to the vibrational time scale to achieve greater equilibration. The magnitude of \(Da_{v,s}\) varies with each species' vibrational properties: \(Da_{v,\mathrm{N_2}}\) exceeds \(Da_{v,\mathrm{O_2}}\) and \(Da_{v,\mathrm{NO}}\) because \(\mathrm{N_2}\) has a higher characteristic vibrational temperature \(\Theta_v\) and therefore relaxes more slowly.

Over \(x/\delta_0^{*}\in[200,\,800]\), the \(Da_{v,s}\) profiles increase monotonically whilst retaining a similar shape: the peak magnitude grows at successive stations, its wall-normal location remains approximately fixed within the boundary layer, and the outer-layer decay becomes more gradual as the boundary layer thickens. The larger fractional rise of the peak in panel~(d) (approximately \(70\%\)) relative to panel~(a) (approximately \(40\%\)) reflects the different baseline magnitudes and growth rates across configurations. As the boundary layer develops downstream, \(\tau_{\mathrm{flow}}\) increases due to layer thickening whilst temperatures rise and thereby reduce \(\tau_v\), causing \(Da_{v,s}\) to increase as relaxation progressively outpaces advection in the time-scale competition.

Comparing across all configurations, the approach to equilibrium follows a clear hierarchy based on the magnitude of \(Da_{v,s}\): case~(a) (\(M_e = 6\), \(T_w = 1200\)~K) achieves the highest degree of vibrational equilibration with the largest Damköhler numbers, followed by case~(c) (\(M_e = 8\), \(T_w = 1200\)~K), then case~(b) (\(M_e = 6\), \(T_w = 1800\)~K), with case~(d) (\(M_e = 8\), \(T_w = 1800\)~K) remaining furthest from equilibrium. Wall cooling promotes equilibration: cooler walls lengthen \(\tau_v\) through reduced temperatures but also shorten \(\tau_{\mathrm{flow}}\) through thinner boundary layers, with the residence-time effect dominating to yield higher \(Da_{v,s}\). Lower Mach numbers provide an additional effect in the same direction.



% Cold walls increase \(/\tau_v\) and thus reduce
% \(\mathrm{Da}_v\), while hot walls decrease \(\tau_v\) and raise \(\mathrm{Da}_v\); likewise, higher Mach numbers elevate post-shock temperatures, shortening \(\tau_v\) and increasing \(\mathrm{Da}_v\).




% In the chemically frozen cases considered here, \(\mathrm{Da}_v\) alone controls the proximity to
% vibrational equilibrium: low \(\mathrm{Da}_v\) corresponds to slow relaxation and larger departures
% from equilibrium, whereas high \(\mathrm{Da}_v\) indicates rapid relaxation toward \(T_v\!\to\!T\). Figure~\ref{figure:Davprofiles} shows the Damkohler profiles at various streamwise locations showcasing the trend to nonequilibrium within the cases. It can be clearly seen that N2 has a lower number than O2 which has a lower number than NO. The increase in the Damkohler number dictates how much the flow is in equilbrium, since that is the equiilibrium limit, Although the effect of this Damkohler number progression was noticed in citep Neveille Candler 2014 have studied the propagation of the the number at variou wall temperatures Kinsely and Zhong studied the effect on a cone, at low enthalpies and found the TNE is destabilizing compared with frozen flow, butstabilizing compared with TE. 




% Increasing $Da_{v,s}$ with $x$ signals transition from vibrationally frozen ($Da_{v,s}\!<\!1$) toward near-equilibrium ($Da_{v,s}\!>\!1$); the species-dependent ordering typically shows faster-relaxing NO (larger $Da_{v,\mathrm{NO}}$) relative to $\mathrm{O_2}$ and $\mathrm{N_2}$.

% A comparison of two-dimensional DNS boundary layer profiles to their respective locally self-similar profiles is shown in \cref{figure:flatplateprofiles} for velocity (top row), temperature (middle row) and vibrational temperature (bottom row). The left column is for the CPG assumption, the middle column is for TEQ and the right-hand column is for TNE, for which the comparison is with a TFG similarity solution. The profiles are of the Mach 6 cold wall condition taken at approximately 75\% of the simulation domain length. The profiles of the CPG and TEQ cases show an excellent agreement with the respective similarity profiles. It can be seen that the temperature peak in the TEQ case is reduced by approximately 8\% relative to the CPG case since some of the energy goes into the vibrational modes. For the TEQ case we have $T_v=T$, while vibrational excitation is not considered in the CPG model. The TNE simulations are initialised with the TFG similarity condition and the relaxation of $T_v$ toward $T$ can be observed. The relation is some way from being complete by the profile plotted. As a result, some additional energy appears in the vibration mode relative to the TFG similarity solution and the translation-rotation temperature decreases slightly.


% \subsubsection{Influence of thermal nonequilibrium}



% --- Growth rates of disturbances ---
% \subsection{Growth rates of disturbances}
% The base flow was perturbed via wall suction--blowing to excite the most unstable second--mode instability. These disturbances are acoustic in nature, arising from partial trapping of waves between the wall and the \emph{relative sonic line}. Within this region the disturbance behaves as a standing wave formed by repeated reflections, while outside the sonic line it becomes evanescent. This trapped, waveguide--like behaviour is illustrated in \cref{figure:mackmodes_mach8}, which shows the spatial structure of the pressure disturbance, \(\tilde{p}\), within a narrow streamwise region of the Mach~8 \gls{TNE} cold--wall case. Each panel corresponds to a successive window in \(x/\delta^{*}_{0}\) and is shown at a \(1{:}1\) aspect ratio to preserve the wave geometry. The disturbance field displays a clear acoustic signature: phase lines are inclined downstream, and oscillations are confined between the wall and the relative sonic line. The alternating regions of compression and rarefaction reveal repeated reflections within this trapped layer, consistent with a slow acoustic (second--mode) wave. The near--wall amplification visible toward the downstream end of each window indicates the persistence of instability growth despite local attenuation zones associated with broadband interference. The solid horizontal line denotes the boundary--layer edge, \(\delta_{99}\), beyond which the disturbance decays rapidly, confirming that the second--mode energy remains concentrated within the viscous sublayer. Overall, the field visualisation highlights the waveguide--like behaviour of the second--mode structure and its spatial modulation under realistic broadband excitation in thermal nonequilibrium conditions. The field is divided into five consecutive streamwise windows, each displayed at a \(1{:}1\) aspect ratio to highlight its evolution. The modulation of pressure amplitude and phase is characteristic of Mack modes: wall--peaked amplitudes, phase rotation across the sonic line, and regions of partial attenuation due to interference between tones of different frequencies and growth rates. Over \(x/\delta^{*}_{0} \in [400,\,600]\), the local reduction in amplitude reflects this broadband interference rather than a loss of instability. Within this region, the dominant energetic components correspond to frequencies \(\tilde f = 0.09,\,0.10,\text{ and }0.11\), which together produce the observed spatial beating pattern in the disturbance envelope. Their individual growth and decay rates differ slightly with \(x\), so that constructive and destructive superposition results in the modulated pressure field evident in \cref{figure:mackmodes_mach8}. The iteration indices of these modes, listed within the figure, identify their respective phase locations and relative amplitudes, providing a direct link between the broadband forcing spectrum and the spatial interference structure visible in the wall–normal pressure field. Similar patterns have been reported in broadband DNS of Mach~6 cones by \citet{SivasubramanianFasel2014} and \citet{UnnikrishnanGaitonde2021}, the latter has linked such amplitude modulation to non-linear second--mode interactions.




% \subsection{Growth rates of disturbances}
% The base flow was perturbed via wall suction--blowing to excite the most unstable second--mode instability. These disturbances are acoustic in nature, arising from partial trapping of waves between the wall and the \emph{relative sonic line}. Within this region the disturbance behaves as a standing wave formed by repeated reflections, while outside the sonic line it becomes evanescent. This trapped, waveguide--like behaviour is illustrated in \cref{figure:mackmodes_mach8}, which shows the spatial structure of the pressure disturbance, \(\tilde{p}\), within a narrow streamwise region of the Mach~8 \gls{TNE} cold--wall case. Each panel corresponds to a successive window in \(x/\delta^{*}_{0}\) and is shown at a \(1{:}1\) aspect ratio to preserve the wave geometry. The disturbance field displays a clear acoustic signature: phase lines are inclined downstream, and oscillations are confined between the wall and the relative sonic line. Alternating regions of compression and rarefaction reveal repeated reflections within this trapped layer, consistent with a slow acoustic (second--mode) wave. The near--wall amplification visible toward the downstream end of each window indicates the persistence of instability growth despite local attenuation zones associated with broadband interference. The solid horizontal line denotes the boundary--layer edge, \(\delta_{99}\), beyond which the disturbance decays rapidly, confirming that the second--mode energy remains concentrated within the viscous sublayer. The modulation of pressure amplitude and phase is characteristic of Mack modes: wall--peaked amplitudes, phase rotation across the sonic line, and regions of partial attenuation due to interference between tones of different frequencies and growth rates. Over \(x/\delta^{*}_{0} \in [400,\,600]\), the local reduction in amplitude reflects this broadband interference rather than a loss of instability. Within this region, the dominant energetic components correspond to frequencies \(\tilde{f} = 0.09,\,0.10,\text{ and }0.11\), which together produce the observed spatial beating pattern in the disturbance envelope. Their individual growth and decay rates differ slightly with \(x\), so that constructive and destructive superposition results in the modulated pressure field evident in \cref{figure:mackmodes_mach8}. The iteration indices of these modes, listed within the figure, identify their respective phase locations and relative amplitudes, providing a direct link between the broadband forcing spectrum and the spatial interference structure visible in the wall--normal pressure field. Similar patterns have been reported in broadband DNS of Mach~6 cones by \citet{SivasubramanianFasel2014} and \citet{UnnikrishnanGaitonde2021}, the latter linking such amplitude modulation to non--linear second--mode interactions.


% \begin{figure}[t]
%     \centering
%     \includegraphics[width=13cm]{resources/figures/disturbance/result_mackmodes_m812.pdf}
%     \caption{Contours of pressure disturbance $\hat{p}$ in a narrow streamwise window for the Mach~8 TNE cold–wall case. The solid horizontal line marks the boundary–layer thickness $\delta_{99}$}
%     \label{figure:mackmodes_mach8}
% \end{figure}


% which shows contours of the wall--normal pressure field obtained by \emph{superposing all imposed forcing frequencies}.
% The base flow was perturbed via wall suction–blowing to excite the most unstable second–mode instability. Second–mode disturbances are acoustic in nature and arise from partial trapping within the boundary layer, specifically between the wall and the \emph{relative sonic line} . Consequently, acoustic waves are oscillatory on the subsonic side of this line but become evanescent outside it, leading to partial reflection back into the boundary layer. This behaviour is illustrated in \cref{figure:mackmodes_mach8}, which shows contours of the pressure eigenfunction for an unstable second–mode, slow acoustic disturbance. The solid curve denotes the boundary–layer thickness, \(\delta_{99}\). The disturbance exhibits a trapped, waveguide–like structure: the wall–normal envelope peaks near the relative sonic line, with repeated reflections between this line and the wall. Outside the sonic line the field is evanescent, whereas within it the incident–reflected acoustic components form the standing–wave pattern characteristic of second–mode amplification.

% Figure~\ref{figure:mackmodes_mach8} shows the wall–normal pressure field obtained by \emph{superposing all imposed forcing frequencies}. For readability, the field is divided into five vertically stacked panels, each displayed at a \(1{:}1\) aspect ratio to track the evolution over successive streamwise windows. The streamwise modulation of the perturbation pressure is characteristic of Mack modes: amplitudes peak at the wall, and a clear phase rotation occurs across the relative sonic line, indicating partial reflection and trapping between the wall and the sonic line. The solid curve marks \(\delta_{99}\). A local attenuation is evident over \(x/\delta^{*}_{0}\in[400,\,600]\), arising from the superposition of components with broadband of frequencies realistic to the flight environment and growth rates some locally unstable and amplifying, others decaying, so that their interference yields a net reduction in amplitude. Whereas the growth and decay of a single–frequency Mack mode have been well documented since the works of \citet{GasterGrant1975}, the present \emph{broadband} forcing exposes the interference patterns and spatial intermittency that emerge when many tones are combined. The resulting behaviour is in qualitative agreement with broadband DNS studies on Mach~6 cones reported by \citet{SivasubramanianFasel2014} and \citet{UnnikrishnanGaitonde2021}. The latter further demonstrated that the observed decay followed by renewed growth is indicative of non-linear second–mode dynamics.



% Figure~\ref{figure:mackmodes_mach8} shows the wall–normal pressure field obtained by \emph{superposing all imposed forcing frequencies}. For readability, the field is divided into five vertically stacked panels, each displayed at a \(1{:}1\) aspect ratio to track the evolution over successive streamwise windows. The streamwise modulation of the perturbation pressure is characteristic of Mack modes: amplitudes peak at the wall, and a clear phase rotation occurs across the relative sonic line, indicating partial reflection and trapping between the wall and the sonic line. The solid curve marks \(\delta_{99}\). A local attenuation is evident over \(x/\delta^{*}_{0}\in[400,\,600]\), arising from the superposition of components with broadband of frequencies realistic to the flight environment and growth rates some locally unstable and amplifying, others decaying, so that their interference yields a net reduction in amplitude. Whereas the growth and decay of a single–frequency Mack mode have been well documented since the works of \citet{GasterGrant1975}, the present \emph{broadband} forcing exposes the interference patterns and spatial intermittency that emerge when many tones are combined. The resulting behaviour is in qualitative agreement with broadband DNS studies on Mach~6 cones reported by \citet{SivasubramanianFasel2014} and \citet{UnnikrishnanGaitonde2021}. The latter further demonstrated that the observed decay followed by renewed growth is indicative of non-linear second–mode dynamics.







 

% ; the dashed line indicates the relative sonic line $y_s$ (trapping layer). Panels show the trapped acoustic cell for three assumed phase speeds: (a) $c_r=u_e$, (b) $c_r=u_e+a_f$, and (c) $c_r=u_e-a_f$ (see inset labels). The disturbance is confined between the wall and $y_s$, consistent with the second–mode trapped–wave picture.

% The mode shapes of the Mack--mode instability are compared in \cref{figure:mach8eigenfuncs}, where the eigenfunctions of the pressure disturbance are normalised by their respective maxima. The profiles correspond to three streamwise locations and demonstrate excellent qualitative agreement between the \gls{LST} predictions and the \gls{DNS} data. In all cases, the pressure amplitude peaks at the wall and decays monotonically across the boundary layer, confirming the slow acoustic nature of the second--mode instability. The wall confinement strengthens with increasing frequency, as expected for higher--order acoustic modes whose wavelengths are shorter relative to the local boundary--layer thickness. Downstream, the envelope broadens slightly and the oscillatory structure weakens, consistent with the local reduction in growth rate as the boundary layer thickens. Notably, at \(x/\delta^{*}_{0}=500\), only the lowest frequency component (\(f=0.09\)) remains dynamically significant, whereas the higher--frequency modes (\(f=0.10\) and \(f=0.11\)) have decayed below the local instability threshold and thus no longer influence the disturbance field. This selective persistence illustrates how the broadband excitation evolves toward a single dominant tone, with the most amplified mode dictating the downstream behaviour. Overall, the eigenfunctions exhibit the characteristic wall--peaked, waveguide--like structure predicted by linear theory, confirming that the dominant disturbances observed in the \gls{DNS} correspond to the classical second--mode family.


\subsection{Growth rates of disturbances}
The base flow was perturbed via wall suction-blowing to excite the most unstable second-mode instability. These disturbances arise from partial trapping of waves between the wall and the relative sonic line. Within this region the disturbance behaves as a standing wave formed by repeated reflections, whilst outside the sonic line it becomes evanescent. This trapped, waveguide-like behaviour is illustrated in \cref{figure:mackmodes_mach8}, which shows the spatial structure of the pressure disturbance $\tilde{p}$ (Pa) within successive streamwise windows for the Mach~8 cold-wall \gls{TNE} case. Each panel is shown at a $1{:}1$ aspect ratio to preserve the wave geometry. 

The disturbance field displays a clear acoustic signature: phase lines are inclined downstream, and oscillations are confined between the wall and the relative sonic line (marked by the solid horizontal line at $\delta_{99}$). Alternating regions of compression and rarefaction reveal repeated reflections within this trapped layer, consistent with the second-mode acoustic wave. The near-wall amplification visible towards the downstream end of each window demonstrates the expected exponential growth of the instability.

\begin{figure}[t]
    \centering
    \includegraphics[width=13cm]{resources/figures/disturbance/result_mackmodes_m812.pdf}
    \caption{Contours of pressure disturbance $\tilde{p}$ in successive streamwise windows for the Mach~8 TNE cold-wall case. The solid horizontal line marks the boundary-layer thickness $\delta_{99}$.}
    \label{figure:mackmodes_mach8}
\end{figure}

The modulation of pressure amplitude and phase is characteristic of Mack modes: wall-peaked amplitudes, phase rotation across the sonic line, and regions of partial attenuation due to interference between tones. Over $x/\delta^{*}_{0} \in [400, 600]$, the local reduction in amplitude reflects broadband interference rather than loss of instability. Within this region, the dominant energetic components correspond to frequencies $\tilde{f} = 0.09$,\ $0.10$, and $0.11$, which together produce the observed spatial beating pattern in the disturbance envelope. Their individual growth rates differ slightly with $x$, leading to constructive and destructive superposition that results in the modulated pressure field evident in \cref{figure:mackmodes_mach8}. Similar patterns have been reported in broadband DNS of Mach~6 cones by \citet{SivasubramanianFasel2014} and \citet{UnnikrishnanGaitonde2021}, the latter linking such amplitude modulation to nonlinear second-mode interactions. However, the linearity of the current study shows the behaviour to be an effect of linear interactions among multiple frequency disturbances.




\begin{figure}[t]
    \centering
    \includegraphics[width=11cm]{resources/figures/disturbance/result_mach8_tne_eigenfuncs_turbo.pdf}
    \caption{Normalised pressure mode shapes at selected streamwise locations for the Mach~8 TNE cold-wall case, showing the wall-normal distribution of second-mode disturbances at different frequencies.}
    \label{figure:mach8eigenfuncs}
\end{figure}


The mode shapes of the Mack-mode instability are compared in \cref{figure:mach8eigenfuncs}, where the mode shapes of the pressure disturbance are normalised by their respective maxima. The profiles correspond to three streamwise locations and demonstrate excellent qualitative agreement between the \gls{LST} predictions and the \gls{DNS} data. In all cases, the pressure amplitude peaks at the wall and decays monotonically across the boundary layer, confirming the slow acoustic nature of the second-mode instability. The wall confinement strengthens with increasing frequency, as expected for higher-order acoustic modes whose wavelengths are shorter relative to the local boundary-layer thickness.

As the flow develops downstream, the boundary layer thickens and the local instability pocket shifts towards lower frequencies, leading to a gradual broadening of the wall-normal envelope. This trend is consistent with the reduction in local growth rate predicted by \gls{LST} and matches the spatial evolution observed in the \gls{DNS}. Notably, at $x/\delta^{*}_{0} = 500$, only the lowest frequency component ($f = 0.09$) remains dynamically significant, whereas the higher-frequency modes ($f = 0.10$ and $0.11$) have decayed below the local instability threshold. The higher-frequency waves amplify earlier and penetrate further downstream before saturating, while the lower-frequency component persists over a longer streamwise extent. This selective persistence reflects the natural frequency filtering of spatially developing hypersonic boundary layers under \gls{TNE} conditions, whereby the broadband excitation evolves towards a single dominant tone governed by the local growth-rate distribution.

Overall, the mode shapes exhibit the characteristic wall-peaked, waveguide-like structure predicted by linear theory, confirming that the dominant disturbances observed in the \gls{DNS} correspond to the classical second-mode family and that the simulation captures the correct progression of instability growth.

% The mode shapes of the Mack--mode instability are compared in \cref{figure:mach8eigenfuncs}, where the eigenfunctions of the pressure disturbance are normalised by their respective maxima. The profiles correspond to three streamwise locations and demonstrate excellent qualitative agreement between the \gls{LST} predictions and the \gls{DNS} data. In all cases, the pressure amplitude peaks at the wall and decays monotonically across the boundary layer, confirming the slow acoustic nature of the second--mode instability. The wall confinement strengthens with increasing frequency, as expected for higher--order acoustic modes whose wavelengths are shorter relative to the local boundary--layer thickness. As the flow develops downstream, the boundary layer thickens and the local instability pocket shifts toward lower frequencies, leading to a gradual broadening of the wall--normal envelope and a weakening of the oscillatory structure. This trend is consistent with the reduction in local growth rate predicted by \gls{LST} and matches the spatial evolution observed in the \gls{DNS}. Notably, at \(x/\delta^{*}_{0}=500\), only the lowest frequency component (\(f=0.09\)) remains dynamically significant, whereas the higher--frequency modes (\(f=0.10\) and \(f=0.11\)) have decayed below the local instability threshold and no longer contribute to the disturbance field. The higher--frequency waves amplify earlier and penetrate further downstream before saturating, while the lower--frequency component persists over a longer streamwise extent. This selective persistence reflects the natural frequency filtering of spatially developing hypersonic boundary layers under \gls{TNE} conditions, whereby the broadband excitation evolves toward a single dominant tone governed by the local growth--rate distribution. Overall, the eigenfunctions exhibit the characteristic wall--peaked, waveguide--like structure predicted by linear theory, confirming that the dominant disturbances observed in the \gls{DNS} correspond to the classical second--mode family and that the simulation captures the correct progression of instability growth.




% After a Fourier decomposition, the response at different frequencies was extracted, with results shown in \cref{figure:growthrates}. Each panel in the figure shows the disturbance amplitude (normalised by the forcing amplitude, $a_0$) as a function of streamwise distance (in m) from the leading edge of the wedge for three curves (blue for CPG, green for TEQ and red for TNE). The rows show the four combinations of inflow properties i.e.~$M_e=6$ cold, $M_e=6$ hot, $M_e=8$ cold and $M_e=8$ hot. The columns show three frequencies. At $M_e=6$ these are $f_1=105.08 \ \mathrm{kHz}$ $ f_2=113.17 \ \mathrm{kHz}$ and $f_3=121.25 \ \mathrm{kHz}$, whereas at $M_e=8$ they are $f_1=85.43 \ \mathrm{kHz}$ $f_2= 100.9 \ \mathrm{kHz}$ and $f_3= 108.73 \ \mathrm{kHz}$. It can be seen that the TEQ cases exhibit significantly higher amplitudes compared to the CPG and TNE simulations. The amplitudes of CPG and TNE cases are much closer, with TNE showing slightly higher values, consistent with the base flows tending slowly towards equilibrium, although we also need to note that the boundary layers grow differently in each case. Kline et al~\cite{Kline2018a} also explored varying relaxation times in TNE simulations and observed that amplification rates in TNE tend towards the TEQ case with appropriate time scales. These findings highlight the importance of considering thermal nonequilibrium effects in stability studies, as it can both stabilise the flow compared to thermal equilibrium and destabilise it compared to a calorically perfect gas flow. 




% --- Amplification independence ---
% \subsubsection{Amplification independence}
% \noindent\textit{What to include.}
% Test collapse of amplification curves under appropriate normalisations (e.g.\ plot $N$ vs.\ $Re_\theta$ or normalise $\int \alpha_i\,\mathrm{d}x$ by a local acoustic or semilocal viscous length). Demonstrate collapses across Mach and $T_w$ within a fixed thermal model; then show deviations when switching CPG$\leftrightarrow$TEQ$\leftrightarrow$TNE, attributing differences to changes in sound speed and bulk viscosity rather than vibrational excitation alone. Report a collapse metric (e.g.\ RMS spread of $N$ at fixed abscissa) and include a concise nondimensionalisation table.

% --- Receptivity to Numerical Schemes ---
% \subsubsection{Receptivity to Numerical Schemes}
% \noindent\textit{What to include.}
% Document how numerical choices (reconstruction order, WENO variant, filtering, sponge treatment, inflow handling) modify initial disturbance content and early growth. Show inlet spectra and near-inlet $\alpha_i$ from short-window FFTs; compare against \gls{LST} receptivity to separate physical from numerically seeded content. Provide grid/time refinement trends (amplitude error vs.\ $\Delta x,\Delta t$) and state the accepted production configuration.

% \newpage

% --- Scaling of instability parameters ---
% \subsection{Scaling of instability parameters}

\subsection{Comparative analysis of instability behaviour}
Having established suitable normalisations for frequency, phase speed, and wall–normal peak location, this section examines how these scaled instability parameters vary with thermal excitation model and Mach number. The objective is to assess whether a universal scaling can be achieved across the \gls{CPG}, \gls{TEQ}, and \gls{TNE} configurations, and to quantify the shifts in mode behaviour induced by thermal nonequilibrium. Emphasis is placed on identifying how excitation effects modify the most–amplified frequency and the location of the trapped acoustic layer, before extending the comparison to highlight systematic differences between the Mach~6 and Mach~8 cases.

\subsubsection{Validation of LST predictions against DNS}

\begin{figure}[ht]
    \centering
    \includegraphics[width=12cm]{resources/figures/disturbance/result_mach6_tne_dns_lst_nfactor.pdf}
    \caption{Streamwise evolution of $N$ factors for Mach~6 flow over the wedge configuration. Panel (a) presents the direct comparison between LST predictions and receptivity-shifted DNS data for multiple non-dimensional frequencies. (b) shows LST-predicted $N$ factors with the corresponding $N$-factor envelope. (c) displays DNS-derived $N$ factors after receptivity correction, with the envelope.}
    \label{figure:mackmodes_mach6_nfactor_comparison}
\end{figure}


While \gls{LST} provides rapid predictions of instability growth rates, its reliance on locally parallel flow assumptions necessitates validation against high-fidelity \gls{DNS} data. Direct comparison between \gls{LST}-derived $N$ factors and those extracted from \gls{DNS} is complicated by receptivity effects, the efficiency with which free-stream disturbances couple into boundary-layer instabilities depends on the excitation mechanism, frequency, and local flow conditions. To account for these receptivity variations, $N$ factors computed from \gls{DNS} are vertically shifted to align with corresponding \gls{LST} curves for each frequency and thermal model. These shifts quantify the receptivity coefficient, with smaller shifts indicating more efficient coupling between the imposed disturbance and the resulting instability mode. Once aligned, an $N$-factor envelope is constructed from the combined \gls{LST} and shifted \gls{DNS} data, providing a unified framework for assessing instability amplification across the streamwise domain.

Figure~\ref{figure:mackmodes_mach6_nfactor_comparison} illustrates the $N$-factor evolution for the Mach~6 configuration across multiple frequencies and thermal excitation models. Panel (a) demonstrates that after applying the receptivity shift, the \gls{DNS}-derived $N$ factors align well with \gls{LST} predictions, validating the theoretical framework for the range of frequencies examined. The $N$-factor envelopes constructed from both \gls{LST} (panel b) and \gls{DNS} (panel c) exhibit consistent maximum amplification trends, providing confidence in the stability analysis methodology. The close agreement between the two approaches confirms that, despite the inherent limitations of the parallel-flow assumption, \gls{LST} captures the essential physics of second-mode instability growth in hypersonic boundary layers when appropriately corrected for receptivity effects.

\begin{figure}[t]
    \centering
    \includegraphics[width=14cm]{resources/figures/disturbance/result_receptivity_coefficients.pdf}
    \caption{Receptivity coefficients quantifying the vertical shift required to align DNS-derived $N$ factors with LST predictions. (a) shows coefficients for Mach~6 cases across CPG, TEQ, and TNE thermal models at two wall temperatures ($T_w/T_\infty = 1200$~K and $1800$~K). (b) presents the corresponding Mach~8 data.}
    \label{figure:mackmodes_receptivity_coefficients}
\end{figure}

The receptivity coefficients required to align \gls{DNS} and \gls{LST} $N$ factors are quantified in Figure~\ref{figure:mackmodes_receptivity_coefficients}. Following \citet{DeTullio2013}, wall suction and blowing is employed as the disturbance source, deemed the most appropriate excitation mechanism for investigating second-mode instabilities. For the Mach~6 cases, both \gls{TNE} and \gls{CPG} models exhibit relatively low receptivity shifts across the frequency range, suggesting efficient disturbance coupling in these configurations. In contrast, the \gls{TEQ} model requires notably larger shifts, indicating reduced receptivity potentially linked to differences in mean temperature profiles and thermochemical damping mechanisms that alter the modal structure. At higher frequencies, where disturbance peaks occur closer to the wall suction/blowing region, transient effects become more pronounced, necessitating larger receptivity corrections across all thermal models. This trend reflects the increased sensitivity of near-wall modes to local flow perturbations and the departure from the parallel-flow assumption inherent in \gls{LST}. The Mach~8 cases exhibit qualitatively similar behaviour, with \gls{TNE} and \gls{CPG} models again demonstrating comparable and relatively modest receptivity shifts, whilst the \gls{TEQ} model shows more variable coupling efficiency. One notable exception appears at $\hat f = 0.09$, where an anomalous receptivity shift deviates from the otherwise consistent pattern this may warrant further investigation into local flow features or numerical artefacts at this specific frequency. Overall, the systematic differences in receptivity between thermal models underscore the importance of thermochemical effects not only in modifying instability growth rates, but also in governing the initial amplitude of disturbances entering the boundary layer. The systematic variation in receptivity with thermal model therefore validates the \gls{LST} framework whilst underscoring the central role of thermal processes in governing disturbance coupling to boundary-layer instabilities.


% \noindent\textit{What to include.}
% Establish scalings for most-amplified frequency, phase speed, and wall-normal peak location using local means (e.g.\ $f\,\delta^\ast/u_e$, $f\,\mu/(\rho a^2)$, $y_{\max}/\delta$). Identify which scaling best collapses first- and second-mode trends across Mach and $T_w$. Track the synchronisation point and show how TEQ/TNE shift the acoustic trapping layer through altered sound speed and bulk viscosity. Provide compact regression fits for downstream use in maps/prediction.



% \subsubsection{Influence of Thermal excitations}

% \subsubsection{Comparison across Mach numbers}

% Focusing on a short window of the Mach~6 cold–wall simulation, \cref{figure:mackmodes_mach6_temp} presents the translational–rotational temperature disturbance over \(x/\delta^{*}_{0}\in[351,\,369]\) for \gls{CPG}, \gls{TEQ}, and \gls{TNE}; this window is representative of the behaviour observed across the remaining configurations. The field exhibits near–wall maxima and downstream–tilted phase lines, consistent with a trapped second–mode wave between the wall and the relative sonic line (dash–dotted). The solid curve marks \(\delta_{99}\), and thin overlays (where shown) indicate the phase–speed loci \(c_r=u_e\) and \(c_r=u_e+a_f\), which align with the observed phase progression and relate the disturbance pattern to the local base–flow acoustics. Within this region, the dominant nondimensional frequencies are \(\tilde f\in[0.14,\,0.16]\). A clear phase rotation is evident across the relative sonic line, with evanescent behaviour outside it, consistent with partial trapping and reflection. This frequency band also accords with linear–stability predictions of the most–amplified second–mode range for the present conditions.




% Across the three thermally excited models the qualitative structure within the boundary layer is similar, although \gls{TEQ} yields the largest disturbance amplitudes and \gls{TNE} displays a slight downstream phase (or onset) delay. In \gls{CPG} and \gls{TEQ} the waves decay more rapidly inside the layer, whereas in \gls{TNE} a small fraction of energy leaks into the free stream via partial transmission at the second sonic line, producing weak dispersive signatures above the layer. This behaviour is compatible with a weakly radiating supersonic component; nevertheless, interactions within the broadband slow–acoustic spectrum must also be considered, and the observed leakage is too small to imply sustained radiation into the free stream~\citep{KniselyZhong2019}. 

% \subsubsection{Influence of thermal excitations}
% To illustrate the influence of thermal excitations on the instability structure, a representative window of the Mach~6 cold-wall simulation is first examined. \Cref{figure:mackmodes_mach6_temp} shows the translational-rotational temperature disturbance over $x/\delta^{*}_{0}\in[351,\,369]$ for \gls{CPG}, \gls{TEQ}, and \gls{TNE}, capturing the characteristic second-mode behaviour observed across all configurations. The disturbance field exhibits near-wall maxima and downstream-tilted phase lines, consistent with a trapped acoustic wave between the wall and the relative sonic line (dash-dotted). The solid curve denotes $\delta_{99}$, and thin overlays (where shown) indicate the phase-speed loci $c_r=u_e$ and $c_r=u_e+a_f$, which align with the observed phase progression and link the disturbance to the local base-flow acoustics. Within this region, the dominant non dimensional frequencies are $\tilde{f}\in[0.13,\,0.15]$, with clear phase rotation across the sonic line and evanescent behaviour outside it, consistent with partial trapping and reflection. This frequency range agrees with linear-stability predictions for the most-amplified second-mode band under the present conditions.

% Across the three thermally excited models, the qualitative structure within the boundary layer is similar, although \gls{TEQ} yields the largest disturbance amplitudes and \gls{TNE} displays a slight downstream phase delay. In \gls{CPG} and \gls{TEQ} the waves decay more rapidly inside the layer, whereas in \gls{TNE} a small fraction of energy leaks into the freestream via partial transmission at the relative sonic line, producing weak dispersive signatures above the layer. This behaviour is compatible with a weakly radiating supersonic component; nevertheless, interactions within the broadband slow-acoustic spectrum must also be considered, and the observed leakage is too small to imply sustained radiation into the freestream~\citep{KniselyZhong2019}.

% \begin{figure}[t]
%     \centering
%     \includegraphics[width=13cm]{resources/figures/disturbance/result_mackmodes_m612_temperature.pdf}
%     \caption{Contours of pressure disturbance $\hat{T}$ in a narrow streamwise window for the Mach~6 hot–wall case.}
%     \label{figure:mackmodes_mach6_temp}
% \end{figure}

% The small phase delay and reduced confinement observed in the \gls{TNE} field indicate that vibrational relaxation weakens the thermoacoustic feedback responsible for second-mode amplification. This behaviour corresponds to a downstream shift in the local growth-rate envelope, consistent with the delayed onset of oscillations visible in \cref{figure:mackmodes_mach6_temp}. The faint signatures above the boundary layer edge in the \gls{TNE} case imply a modest reduction in acoustic impedance, allowing limited energy leakage into the outer flow. Such partial transmission is compatible with radiative damping effects previously reported for high-enthalpy conditions~\citep{Bitter2015}. Overall, whilst the second-mode structure remains qualitatively similar across the three models, the relative amplitudes and onset locations demonstrate how thermal excitation influences both the efficiency and timing of instability growth, trends that will be quantified in the subsequent growth-rate analysis.

% \begin{figure}[t]
%     \centering
%     \includegraphics[width=13cm]{resources/figures/disturbance/result_mach612_eigenfuncs.pdf}
%     \caption{Contours of pressure disturbance $\hat{p}$ in a narrow streamwise window for the Mach~8 TNE cold–wall case. The solid horizontal line marks the boundary–layer thickness $\delta_{99}$}
%     \label{figure:mach6_tempeigenfuncs}
% \end{figure}

% In contrast to the similar phase organisation observed in \cref{figure:mackmodes_mach6_temp}, the temperature eigenfunctions in \cref{figure:mach6_tempeigenfuncs} reveal clear quantitative differences between the models when normalized consistently. The \gls{TNE} case exhibits noticeably larger temperature eigenfunction amplitudes compared to \gls{CPG} and \gls{TEQ}, particularly in the near-wall region and within the acoustic cavity. This behaviour arises from the delayed energy exchange between translational and vibrational reservoirs under thermal nonequilibrium. In \gls{TNE}, the vibrational temperature $T_v$ does not equilibrate rapidly enough to absorb energy from the disturbance, leaving a larger fraction concentrated in the translational-rotational mode and enhancing the local compressibility. As a result, the modal temperature response is amplified. The \gls{CPG} and \gls{TEQ} cases show comparable eigenfunction amplitudes, both smaller than \gls{TNE}. In \gls{TEQ}, part of the internal energy is absorbed by vibrational modes in equilibrium, effectively buffering the translational temperature oscillations. Here, both the translational-rotational and vibrational temperatures represent the same equilibrium energy packet, so the modal disturbance energy is distributed across internal modes rather than concentrated in the translational component.

% The apparent discrepancy between these eigenfunction amplitudes and the spatial disturbance field observed in \cref{figure:mackmodes_mach6_temp} where \gls{TEQ} exhibits the largest spatial amplitudes highlights the distinction between modal structure and wave packet evolution. The spatial field $\tilde{T}$ reflects the cumulative effect of acoustic resonance efficiency, growth-rate integration, and multi-frequency interactions within the boundary layer. The larger spatial amplitudes in \gls{TEQ} arise from enhanced wave trapping and more efficient acoustic feedback within the cavity, leading to greater cumulative amplification despite the individual modal response being buffered. Conversely, the \gls{TNE} case exhibits reduced spatial amplitudes due to partial energy leakage and weakened acoustic confinement, even though the single-frequency modal sensitivity is enhanced. This demonstrates that thermal equilibrium improves the effectiveness of the second-mode resonance mechanism increasing spatial disturbance build-up through better acoustic trapping whilst thermal nonequilibrium amplifies individual modal responses but reduces overall wave packet growth through radiative damping.


% \begin{figure}[t]
%     \centering
%     \includegraphics[width=13cm]{resources/figures/disturbance/result_mackmodes_m812_temperature.pdf}
%     \caption{Contours of pressure disturbance $\hat{T}$ in a narrow streamwise window for the Mach~6 hot–wall case.}
%     \label{figure:mackmodes_mach8_temp}
% \end{figure}

% At Mach~8, the second-mode disturbance field exhibits lower dominant frequencies and stronger penetration towards the boundary layer edge compared to the Mach~6 case, consistent with the increased boundary layer thickness and modified acoustic cavity dimensions at higher Mach numbers. The temperature disturbance fields reveal characteristic second-mode wave packets with alternating thermal structures and elongated wavelengths. Across the three thermochemical models, \gls{TEQ} exhibits substantially larger spatial disturbance amplitudes than \gls{CPG} and \gls{TNE}, mirroring the behaviour observed at Mach~6. This consistent amplitude hierarchy reflects enhanced acoustic resonance efficiency under thermal equilibrium conditions. At Mach~8, the higher freestream enthalpy leads to elevated boundary layer temperatures, thereby increasing the degree of vibrational excitation across all models. In \gls{TEQ}, the equilibrium assumption ensures that vibrational modes are fully populated according to the local temperature, effectively reducing the frozen sound speed and modifying the acoustic impedance distribution within the boundary layer. This alteration strengthens wave trapping within the acoustic cavity and enhances the constructive interference of reflected acoustic waves, resulting in larger cumulative disturbance amplitudes. The \gls{CPG} case, which neglects vibrational energy entirely, displays well-defined wave structures with moderate amplitudes and sharp thermal gradients. The \gls{TNE} case shows comparable overall amplitude levels to \gls{CPG} but with slightly diffused thermal signatures, reflecting the delayed vibrational relaxation that prevents full energy absorption. The finite relaxation time in \gls{TNE} allows partial vibrational excitation but insufficient equilibration, leading to reduced acoustic confinement and modest energy leakage into the freestream. Although the wavelength and spatial organisation of the Mack modes remain qualitatively similar across all three models, the pronounced amplitude difference in \gls{TEQ} observed consistently at both Mach~6 and Mach~8 demonstrates that thermal equilibrium significantly enhances the spatial build-up of second-mode disturbances through improved acoustic feedback mechanisms, an effect that becomes more pronounced with increasing Mach number due to the greater thermal excitation of vibrational modes.

% To quantify the amplitude variations across the cases, \cref{figure:mach_amplitudecomp} presents the streamwise evolution of pressure disturbance amplitudes extracted from FFT analysis for a range of nondimensional frequencies at both Mach~6 (cold wall) and Mach~8 (hot wall) conditions. The qualitative behaviour observed in the spatial disturbance fields is clearly reflected in these plots: all models exhibit a distinct amplification region associated with second-mode instability, followed by decay downstream of the saturation point. At Mach~6, the \gls{CPG} case shows disturbance growth that peaks sharply near $x/\delta^{*}_{0}\approx 400$, with a well-defined dominant frequency near $\tilde{f}=0.14$. The \gls{TEQ} case exhibits stronger overall amplification with a peak near $x/\delta^{*}_{0}\approx 600$, and displays a broader frequency response, consistent with the enhanced acoustic resonance efficiency discussed previously. The \gls{TNE} case shows a delayed onset and downstream shift of the amplitude peak to approximately $x/\delta^{*}_{0}\approx 650$, with slightly reduced peak amplitudes compared to \gls{TEQ}, reflecting the reduced acoustic confinement and partial energy leakage. At Mach~8, all three models exhibit more complex multi-peaked behaviour with broader frequency distributions, indicative of increased nonlinear interactions and multiple competing instability modes. The \gls{CPG} case displays multiple amplitude peaks across the frequency range, whilst \gls{TEQ} shows the strongest amplification with distinct peaks around $\tilde{f}=0.14$ and lower frequencies. The \gls{TNE} case again exhibits reduced amplitudes and more distributed energy across frequencies compared to \gls{TEQ}. The progressive shift of the amplitude envelope towards lower frequencies downstream in both Mach number cases further illustrates the spatial filtering characteristics of the boundary layer, whereby higher frequencies saturate and decay earlier whilst lower-frequency disturbances persist over longer streamwise extents.

% \begin{figure}[t]
%     \centering
%     \includegraphics[width=11cm]{resources/figures/disturbance/result_mach_amplitude_comp.pdf}
%     \caption{Streamwise evolution of pressure disturbance amplitude spectra for (a) Mach 6 cold wall and (b) Mach 8 hot wall cases at selected frequencies, obtained via Fourier transformation. Legend: Legend: CPG (\protect\cpgline), TEQ (\protect\teqline), TNE (\protect\tneline).}
%     \label{figure:mach_amplitudecomp}
% \end{figure}


% The observed differences between \gls{TEQ} and \gls{TNE} can be understood through the vibrational relaxation timescale relative to the acoustic disturbance period. The vibrational relaxation time $\tau_v$ characterises the rate at which energy transfers between translational and vibrational modes, typically modelled via the Landau-Teller formulation as
% \begin{equation}
% \frac{dT_v}{dt} = \frac{T_{tr} - T_v}{\tau_v},
% \end{equation}
% where $\tau_v$ depends strongly on local temperature and density. The relevant nondimensional parameter governing thermal nonequilibrium effects is the vibrational Damköhler number, $\text{Da}_v = \tau_{\text{flow}}/\tau_v$, which compares the characteristic flow timescale to the relaxation timescale. For second-mode instabilities, the appropriate flow timescale is the acoustic period $\tau_{\text{acoustic}} = 1/f$. When $\text{Da}_v \gg 1$ (fast relaxation), the vibrational modes remain in local equilibrium with the translational temperature, corresponding to the \gls{TEQ} limit. Conversely, when $\text{Da}_v \ll 1$ (slow relaxation), vibrational modes cannot respond to rapid fluctuations, approaching the frozen limit. In the \gls{TNE} case, $\text{Da}_v \sim O(1)$, indicating partial relaxation: vibrational modes are excited but lag behind translational fluctuations, creating a phase difference that reduces the effective compressibility and alters the acoustic wave propagation characteristics. This relaxation lag manifests in two key ways. First, the effective sound speed in \gls{TNE} lies between the frozen (\gls{CPG}) and equilibrium (\gls{TEQ}) values, modifying the acoustic cavity geometry and wave trapping efficiency. Second, the out-of-phase vibrational response introduces a dissipative mechanism that damps acoustic oscillations, reducing disturbance amplitudes and promoting energy leakage into the freestream. The effective heat capacity ratio can be expressed as
% \begin{equation}
% \gamma_{\text{eff}} = \gamma_{\text{frozen}} + (\gamma_{\text{eq}} - \gamma_{\text{frozen}}) \frac{\text{Da}_v}{1 + \text{Da}_v},
% \end{equation}
% illustrating the smooth transition from frozen to equilibrium behaviour as the relaxation rate increases. In \gls{TEQ}, the lower effective $\gamma$ reduces the sound speed and increases acoustic impedance contrast at the boundary layer edge, strengthening wave trapping. In \gls{TNE}, the intermediate $\gamma_{\text{eff}}$ weakens this contrast, allowing partial transmission of acoustic energy. 



\subsubsection{Influence of thermal excitations}
To illustrate how thermal nonequilibrium alters the instability structure, a representative window of the Mach~6 cold-wall simulation is examined. \Cref{figure:mackmodes_mach6_temp} shows the translational-rotational temperature disturbance over $x/\delta^{*}_{0}\in[351,\,369]$ for \gls{CPG}, \gls{TEQ}, and \gls{TNE}—spanning frozen, fully equilibrated, and partially relaxed states—capturing the characteristic second-mode behaviour observed across all configurations. The disturbance field exhibits near-wall maxima and downstream-tilted phase lines, consistent with a trapped acoustic wave between the wall and the relative sonic line (dash-dotted), qualitatively similar to \gls{LST} predictions~\citep{Hudson1997,ZanusEtAl2019}. The solid curve denotes $\delta_{99}$, and thin overlays (where shown) indicate the phase-speed loci $c_r=u_e$ and $c_r=u_e+a_f$, which align the disturbance with the local base-flow acoustics. Within this region, the dominant non dimensional frequencies are $\tilde{f}\in[0.13,\,0.15]$, with clear phase rotation across the sonic line and evanescent behaviour outside it, consistent with partial trapping and reflection. This frequency range agrees with linear-stability predictions for the most-amplified second-mode band under the present conditions~\citep{BitterShepherd2015}.

Across the three models, the qualitative wave structure within the boundary layer remains similar: near-wall acoustic trapping and downstream-tilted phase lines are observed in all cases. However, quantitative differences emerge that reflect the competing timescales of vibrational relaxation and acoustic oscillation. \gls{TEQ} yields the largest spatial disturbance amplitudes, whilst \gls{TNE} displays a slight downstream phase delay and reduced spatial confinement. In \gls{CPG} and \gls{TEQ}, acoustic energy remains confined within the boundary layer with rapid decay inside the layer. Conversely, in \gls{TNE}, a small fraction of energy leaks into the freestream via partial transmission at the relative sonic line, producing weak dispersive signatures above the layer behaviour compatible with weakly radiating acoustic components, though the leakage remains too modest to imply sustained radiation into the freestream~\citep{KniselyZhong2019}.

\begin{figure}[t]
    \centering
    \includegraphics[width=13cm]{resources/figures/disturbance/result_mackmodes_m612_temperature.pdf}
    \caption{Contours of temperature disturbance $\hat{T}$ in a narrow streamwise window for the Mach~6 cold-wall case.}
    \label{figure:mackmodes_mach6_temp}
\end{figure}

The phase delay and reduced confinement in \gls{TNE} provide immediate signatures of weakened thermoacoustic feedback. This behaviour corresponds to a downstream shift in the local growth-rate envelope, consistent with the delayed oscillation onset visible in \cref{figure:mackmodes_mach6_temp}. The faint signatures above the boundary layer edge in \gls{TNE} imply a modest reduction in acoustic impedance, allowing limited energy leakage into the outer flow behaviour compatible with radiative damping previously reported for high-enthalpy conditions~\citep{BitterShepherd2015}. The fact that \gls{TEQ} maintains stronger confinement whilst \gls{TNE} shows leakage suggests that the degree of equilibration directly modulates the acoustic cavity efficiency.

\begin{figure}[t]
    \centering
    \includegraphics[width=13cm]{resources/figures/disturbance/result_mach612_eigenfuncs.pdf}
    \caption{Contours of pressure disturbance $\hat{p}$ in a narrow streamwise window for the Mach~8 TNE cold–wall case. The solid horizontal line marks the boundary–layer thickness $\delta_{99}$}
    \label{figure:mach6_tempeigenfuncs}
\end{figure}

In contrast to the similar phase organisation observed in spatial fields, the temperature mode shape in \cref{figure:mach6_tempeigenfuncs} reveal quantitative differences in modal sensitivity between the models. The \gls{TNE} case exhibits noticeably larger temperature mode shapes amplitudes compared to \gls{CPG} and \gls{TEQ}, with peak values distributed throughout the boundary layer depth. This amplification arises from the delayed energy exchange between translational and vibrational reservoirs: in \gls{TNE}, the vibrational temperature $T_v$ does not equilibrate rapidly enough to absorb disturbance energy, leaving a larger fraction concentrated in the translational-rotational mode and enhancing local compressibility. The \gls{CPG} and \gls{TEQ} cases show comparable and substantially smaller mode shapes amplitudes; \gls{TEQ} benefits from equilibrium energy buffering across vibrational modes, which distributes disturbance energy and reduces the instantaneous translational temperature response. Whilst qualitatively similar modal structures have been observed in high-Mach low-Reynolds number flows~\citep{MiroMiro2019}, the pronounced amplitude hierarchy observed here is distinctive, likely reflecting the specific equilibration regimes of the present configurations. This distinction reveals a fundamental decoupling: \gls{TNE} exhibits greater local modal sensitivity to pressure fluctuations, yet produces smaller cumulative spatial disturbances compared to \gls{TEQ}. This apparent paradox reflects the competition between instantaneous modal response and cumulative wave packet evolution, where enhanced local sensitivity in \gls{TNE} is offset by weakened acoustic confinement.

The spatial field $\tilde{T}$ reflects the cumulative effect of acoustic resonance efficiency, growth-rate integration, and multi-frequency interactions within the boundary layer. The larger spatial amplitudes in \gls{TEQ} arise from enhanced wave trapping and more efficient acoustic feedback, leading to greater cumulative amplification despite the individual modal response being buffered by energy distribution across modes. Conversely, \gls{TNE} exhibits reduced spatial amplitudes due to partial energy leakage and weakened acoustic confinement, even though single-frequency modal sensitivity is enhanced. This demonstrates that thermal equilibrium improves second-mode resonance efficiency through better wave trapping, whilst thermal nonequilibrium amplifies individual modal responses but reduces overall wave packet growth through energy leakage and dissipation.

\begin{figure}[t]
    \centering
    \includegraphics[width=13cm]{resources/figures/disturbance/result_mackmodes_m812_temperature.pdf}
    \caption{Contours of temperature disturbance $\hat{T}$ in a narrow streamwise window for the Mach~8 hot-wall case.}
    \label{figure:mackmodes_mach8_temp}
\end{figure}

At Mach~8, the second-mode disturbance field exhibits lower dominant frequencies and stronger penetration towards the boundary layer edge compared to Mach~6, consistent with the increased boundary layer thickness and modified acoustic cavity dimensions. The characteristic second-mode wave packets display alternating thermal structures and elongated wavelengths. Across the three models, \gls{TEQ} exhibits substantially larger spatial disturbance amplitudes than \gls{CPG} and \gls{TNE}, reinforcing the Mach~6 hierarchy. At higher Mach, the elevated freestream enthalpy increases boundary layer temperatures and vibrational excitation across all models. In \gls{TEQ}, the equilibrium assumption ensures vibrational modes are fully populated, effectively reducing the frozen sound speed and modifying acoustic impedance distribution to strengthen wave trapping. The \gls{CPG} case displays well-defined wave structures with moderate amplitudes and sharp thermal gradients, whilst \gls{TNE} shows diffused thermal signatures reflecting the delayed vibrational relaxation that prevents full energy absorption. Although wavelength and spatial organisation remain qualitatively similar across all models, the persistent amplitude hierarchy—with \gls{TEQ} consistently dominant at both Mach numbers—demonstrates that thermal equilibration significantly enhances spatial disturbance build-up through improved acoustic feedback, an effect amplified with increasing Mach number.


To quantify the amplitude variations, \cref{figure:mach_amplitudecomp} presents the streamwise evolution of pressure disturbance amplitudes extracted from FFT analysis for a range of nondimensional frequencies at both Mach~6 (cold wall) and Mach~8 (hot wall) conditions. All models exhibit a distinct amplification region associated with second-mode instability, followed by downstream decay. At Mach~6, \gls{CPG} shows sharp growth peaking near $x/\delta^{*}_{0}\approx 400$ at $\tilde{f}=0.14$. \gls{TEQ} exhibits stronger amplification with a peak near $x/\delta^{*}_{0}\approx 600$ and broader frequency response, consistent with enhanced acoustic resonance. \gls{TNE} shows delayed onset and a downstream-shifted amplitude peak near $x/\delta^{*}_{0}\approx 650$, with reduced amplitudes reflecting weakened acoustic confinement. These trends confirm prior observations of amplitude suppression in thermochemically nonequilibrium flows~\citep{KinsleyZhong2019}; however, the present results demonstrate that thermal equilibrium represents the amplification limit, with \gls{TEQ} substantially exceeding both \gls{CPG} and \gls{TNE}. At Mach~8, all three models display multi-peaked behaviour with broader frequency distributions, indicative of increased nonlinear interactions, yet the same hierarchy persists: \gls{TEQ} shows the strongest amplification, whilst \gls{TNE} exhibits reduced amplitudes and more distributed energy. This amplitude hierarchy across both Mach numbers is consistent with findings that thermochemical nonequilibrium produces a stabilising effect relative to chemical equilibrium~\citep{Kline2019}; the present thermal nonequilibrium results follow a similar pattern, with thermal equilibrium emerging as the dominant amplification regime. In both cases, the progressive shift of the amplitude envelope towards lower frequencies downstream illustrates spatial filtering: higher frequencies saturate and decay earlier, whilst lower-frequency disturbances persist over longer streamwise extents~\citep{Fedorov2011}.

\begin{figure}[t]
    \centering
    \includegraphics[width=11cm]{resources/figures/disturbance/result_mach_amplitude_comp.pdf}
    \caption{Streamwise evolution of pressure disturbance amplitude spectra for (a) Mach 6 cold wall and (b) Mach 8 hot wall cases at selected frequencies, obtained via Fourier transformation. Legend: CPG (\protect\cpgline), TEQ (\protect\teqline), TNE (\protect\tneline).}
    \label{figure:mach_amplitudecomp}
\end{figure}

% The physical basis for these differences resides in the competition between acoustic timescales and vibrational relaxation. Analogous to the baseline flow equilibration characterised by the global Damköhler number $\text{Da}_v = \tau_{\text{flow}}/\tau_v$ (Section~\ref{section:disturbance:thermal_relaxation}), second-mode disturbances are governed by a localised acoustic Damköhler number where the relevant timescale is the acoustic period $\tau_{\text{acoustic}} = 1/f$ rather than the mean residence time. In \gls{TEQ}, where $\text{Da}_v \gg 1$ globally, vibrational modes remain phase-locked with translational fluctuations, maintaining an effective heat capacity ratio close to the equilibrium value $\gamma_{\text{eq}}$. The reduced sound speed and enhanced acoustic impedance contrast at the boundary layer edge strengthen wave trapping and produce the observed larger spatial amplitudes~\citep{ZhangZhang2022,Varma2025}. In \gls{TNE}, the intermediate $\text{Da}_v \sim O(1)$ characteristic of the present configurations means vibrational modes cannot fully respond to rapid acoustic oscillations: energy excitation lags behind translational fluctuations, creating a phase offset that reduces local compressibility. This manifests through an intermediate effective heat capacity ratio
% \begin{equation}
% \gamma_{\text{eff}} = \gamma_{\text{frozen}} + (\gamma_{\text{eq}} - \gamma_{\text{frozen}}) \frac{\text{Da}_v}{1 + \text{Da}_v},
% \end{equation}
% which weakens acoustic impedance contrast compared to \gls{TEQ}, reducing wave trapping efficiency. Additionally, the out-of-phase vibrational response introduces dissipation that damps acoustic oscillations and promotes energy leakage into the freestream, explaining the observed downstream phase delay, reduced spatial confinement, and amplitude suppression in \gls{TNE} relative to \gls{TEQ}.

The physical basis for these differences resides in the competition between acoustic timescales and vibrational relaxation. Analogous to the baseline flow equilibration characterised by the global Damköhler number $\text{Da}_v = \tau_{\text{flow}}/\tau_v$ (Section~\ref{section}), second-mode disturbances are governed by a localised acoustic Damköhler number where the relevant timescale is the acoustic period $\tau_{\text{acoustic}} = 1/f$ rather than the mean residence time. In \gls{TEQ}, where $\text{Da}_v \gg 1$ globally, vibrational modes remain phase-locked with translational fluctuations, maintaining an effective heat capacity ratio close to the equilibrium value $\gamma_{\text{eq}}$. The reduced sound speed and enhanced acoustic impedance contrast at the boundary layer edge strengthen wave trapping and produce the observed larger spatial amplitudes~\citep{ZhangZhang2022,Varma2025}. In \gls{TNE}, the intermediate $\text{Da}_v \sim O(1)$ characteristic of the present configurations means vibrational modes cannot fully respond to rapid acoustic oscillations: energy excitation lags behind translational fluctuations, creating a phase offset that reduces local compressibility. A critical distinction not previously identified in thermochemical nonequilibrium studies is that the thermal conductivity partition between translational-rotational ($\kappa_{tr}$) and vibrational ($\kappa_v$) modes fundamentally alters the acoustic impedance structure. In \gls{TEQ}, the vibrational thermal conductivity $\kappa_v$ actively extracts energy from boundary layer disturbances, redistributing it into the vibrational reservoir and suppressing peak temperature oscillations. This energy redistribution reduces the local temperature sensitivity and acoustic impedance contrast. In \gls{TNE}, insufficient vibrational equilibration means $\kappa_v$ cannot effectively withdraw this energy on the acoustic timescale; consequently, the translational-rotational temperature response remains amplified, but the asymmetric energy partitioning introduces a dissipative mechanism that damps acoustic oscillations and promotes energy leakage into the freestream. This interplay between thermal conductivity allocation and relaxation timescale explains the observed downstream phase delay, reduced spatial confinement, and amplitude suppression in \gls{TNE} relative to \gls{TEQ}, whilst the effective heat capacity ratio
\begin{equation}
\gamma_{\text{eff}} = \gamma_{\text{frozen}} + (\gamma_{\text{eq}} - \gamma_{\text{frozen}}) \frac{\text{Da}_v}{1 + \text{Da}_v}
\end{equation}
provides only a partial description of the underlying physical mechanisms.





% In the Mach 8 case, the lower frequency response is evident within the domain, and the disturbance is closer to the boundary layer edge, the structures  \citep{KniselyZhong2018} show the excitaiton of the second mode. 

\paragraph{Phase speed and growth-rate structure} The phase speed $c_{\phi}$ and spatial growth rate $\alpha_i$ for the $\tilde{f}=0.16$ disturbance component across Mach~6 cold-wall and hot-wall configurations reveal several key features, as shown in \cref{figure:mach6_cph}. A striking feature emerges in the growth-rate profiles: \gls{TNE} and \gls{CPG} exhibit virtually identical downstream phase delays relative to \gls{TEQ}, despite operating under distinctly different thermal conditions. This phase delay manifests as a delayed onset of amplification and a downstream-shifted growth-rate peak, independent of wall temperature. The phase speeds further reveal quantitative differences: both \gls{TNE} and \gls{CPG} sustain higher phase speeds than \gls{TEQ} throughout the domain, reflecting the altered acoustic impedance structure arising from reduced vibrational excitation. This trend is less pronounced than the acoustic behaviour observed in comparable low-Reynolds number high-enthalpy conditions studied via linear stability theory~\citep{MortensenZhong2013}, suggesting that the nonequilibrium effect on phase speed may saturate or become modulated at higher Reynolds numbers.

\begin{figure}[t]
    \centering
    \includegraphics[width=13cm]{resources/figures/disturbance/result_phase_speed_growth_612_618_five.pdf}
    \caption{Spatial growth rate and phase speed for Mach~6 at $\tilde{f}=0.16$, extracted from DNS. Left: cold-wall conditions. Right: hot-wall conditions. Dashed lines denote acoustic reference speeds $c_{\phi}=1\pm 1/M_e$..}
    \label{figure:mach6_cph}
\end{figure}


A notable feature in the phase-speed profiles is the appearance of regimes where $c_{\phi}$ dips below the $1-1/M_e$ line, indicating local excursion into the acoustic supersonic mode band. Such crossings have been associated with second-mode excitation mechanisms~\citep{KniselyZhong2018}. In the present cold-wall case, these supersonic mode signatures are most pronounced in \gls{TEQ}, whilst remaining subdued in \gls{CPG} and \gls{TNE}. Although a detailed investigation of supersonic mode dynamics is beyond the scope of this study, the observation suggests that thermal equilibration enhances the excitation of these higher-order acoustic components, implying that thermal equilibrium influences not only the primary second-mode growth but also the structure of competing instability branches.

\begin{figure}[t]
    \centering
    \includegraphics[width=13cm]{resources/figures/disturbance/result_Nfactor_612.pdf}
    \caption{$N$ factors acquired from DNS and LST for Mach 6 $T_w=1200 \ K$ cases.Legend: CPG DNS (\protect\cpgline), CPG LST (\protect\cpgdottedline), TEQ DNS   (\protect\teqline), TEQ LST(\protect\teqdottedline), TNE DNS (\protect\tneline), TFG LST (\protect\tnedottedline)}
    \label{figure:m6tw12_Nfactor}
\end{figure}

Finally \gls{DNS} $N$-factors are extracted via spatial integration of the growth rate and compared against \gls{LST} predictions. A vertical shift is applied to account for absolute amplitude differences, with no frequency adjustment. \Cref{figure:m6tw12_Nfactor} presents results for four frequencies ($\tilde{f}=0.13$, $0.14$, $0.15$, $0.16$) across \gls{CPG}, \gls{TEQ}, and \gls{TNE} models for the $M_e=6$, $T_w=1200$ configuration. DNS upstream noise reflects receptivity behaviour; in the growth phase, DNS and \gls{LST} agreement is excellent regarding peak location and magnitude. Notably, \gls{CPG} and \gls{TEQ} agreement quality is comparable despite \gls{LST} using a frozen \gls{CPG} base flow in both cases, confirming that base-flow changes dominate over thermochemical modelling in \gls{LST}. In \gls{TNE}, the frozen base-flow assumption in \gls{LST} introduces an additional approximation (vibrational relaxation not modelled), yet agreement remains reasonable, indicating that nonequilibrium effects on $N$-factor prediction are modest. \gls{TEQ} tends to overestimate $N$-factors slightly, whilst \gls{CPG} remains closest to \gls{TNE} results. 




% At Mach~6, the \gls{CPG} case shows disturbance growth that peaks sharply near $x/\delta^{*}_{0}\approx 400$, with a well-defined dominant frequency near $\tilde{f}=0.14$. The \gls{TEQ} case exhibits stronger overall amplification with a peak near $x/\delta^{*}_{0}\approx 600$, and displays a broader frequency response, consistent with the enhanced acoustic resonance efficiency arising from equilibrium vibrational excitation. The \gls{TNE} case shows a delayed onset and downstream shift of the amplitude peak to approximately $x/\delta^{*}_{0}\approx 650$, with slightly reduced peak amplitudes compared to \gls{TEQ}, reflecting the reduced acoustic confinement due to partial relaxation effects. At Mach~8, all three models exhibit more complex multi-peaked behaviour with broader frequency distributions, indicative of increased nonlinear interactions and the presence of multiple competing instability modes at higher enthalpy conditions. The \gls{CPG} case displays multiple amplitude peaks across the frequency range, whilst \gls{TEQ} shows the strongest amplification with distinct peaks around $\tilde{f}=0.14$ and lower frequencies. The \gls{TNE} case again exhibits reduced amplitudes and more distributed energy across frequencies compared to \gls{TEQ}, consistent with the dissipative effects of finite-rate vibrational relaxation. The progressive shift of the amplitude envelope towards lower frequencies downstream in both Mach number cases further illustrates the spatial filtering characteristics of the boundary layer, whereby higher frequencies saturate and decay earlier whilst lower-frequency disturbances persist over longer streamwise extents.



% After a Fourier decomposition, the response at different frequencies was extracted, with results shown in \cref{figure:growthrates}. Each panel in the figure shows the disturbance amplitude (normalised by the forcing amplitude, $a_0$) as a function of streamwise distance (in m) from the leading edge of the wedge for three curves (blue for CPG, green for TEQ and red for TNE). The rows show the four combinations of inflow properties i.e.~$M_e=6$ cold, $M_e=6$ hot, $M_e=8$ cold and $M_e=8$ hot. The columns show three frequencies. At $M_e=6$ these are $f_1=105.08 \ \mathrm{kHz}$ $ f_2=113.17 \ \mathrm{kHz}$ and $f_3=121.25 \ \mathrm{kHz}$, whereas at $M_e=8$ they are $f_1=85.43 \ \mathrm{kHz}$ $f_2= 100.9 \ \mathrm{kHz}$ and $f_3= 108.73 \ \mathrm{kHz}$. It can be seen that the TEQ cases exhibit significantly higher amplitudes compared to the CPG and TNE simulations. The amplitudes of CPG and TNE cases are much closer, with TNE showing slightly higher values, consistent with the base flows tending slowly towards equilibrium, although we also need to note that the boundary layers grow differently in each case. Kline et al~\cite{Kline2018} also explored varying relaxation times in TNE simulations and observed that amplification rates in TNE tend towards the TEQ case with appropriate time scales. These findings highlight the importance of considering thermal nonequilibrium effects in stability studies, as it can both stabilise the flow compared to thermal equilibrium and destabilise it compared to a calorically perfect gas flow. 

% Finally, the general trend of this is a generalisation. 

% \newpage

% --- Comparison across Mach numbers ---
\subsubsection{Comparison across Mach numbers}

To assess how thermal nonequilibrium effects scale with flow conditions, the disturbance amplitude response and RMS pressure evolution are compared across Mach~6 and Mach~8 at both cold-wall ($T_w=1200$~K) and hot-wall ($T_w=1800$~K) configurations. \Cref{figure:lstdns_amplitude} shows the disturbance amplitude profiles for the most-amplified frequency at each condition, arranged as: top row Mach~6 cold-wall, second row Mach~6 hot-wall, third row Mach~8 cold-wall, and bottom row Mach~8 hot-wall, with \gls{CPG}, \gls{TEQ}, and \gls{TNE} in successive columns. \gls{TEQ} consistently exhibits the largest disturbance amplitudes across all conditions, with \gls{CPG} and \gls{TNE} showing marginal differences in magnitude. Wall temperature elevation produces a notable effect in the \gls{TEQ} case, progressively shifting the amplitude peak closer to the leading edge, whereas the response profiles of \gls{CPG} and \gls{TNE} remain relatively insensitive to wall temperature variation. At higher Mach numbers, the growth characteristics of \gls{CPG} become increasingly similar to those of \gls{TNE}, a trend that is particularly evident in the Mach~8 hot-wall configuration where the frozen-gas model exhibits comparable amplification to the nonequilibrium response.

\begin{figure}[t]
    \centering
    \includegraphics[width=13cm]{resources/figures/disturbance/result_lst_dns_amplitude_comp.pdf}
    \caption{Amplitude of disturbances versus distance from wedge leading edge. Each row represents the configuration (\textit{rows}: (a) $M_e=6$, $T_w=1200$, (b) $M_e=6$, $T_w=1200$, (c) $M_e=6$, $T_w=1200$, (d) $M_e=8$, $T_w=1800$). Each column represents a different frequency, with frequency increasing from left to right, $\tilde f = \ 0.13,0.14,0.15$ for Mach 6 cases and $\tilde f = \ 0.11,0.12,0.13$ for Mach 8 cases.}
    \label{figure:lstdns_amplitude}
\end{figure}

The pressure root-mean-square evolution, shown in \cref{figure:prms_allcases}, provides complementary insight into the cumulative disturbance development across the flow domain. Peak $p_{\text{rms}}$ values decrease systematically from the cold-wall Mach~6 case towards the hot-wall Mach~8 case, which is consistent with the reduced instability growth observed at higher Mach numbers and elevated wall temperatures. \gls{TEQ} shows the largest $p_{\text{rms}}$ magnitudes across all conditions, whilst \gls{CPG} and \gls{TNE} remain reduced but with increasingly similar values at higher Mach numbers. These observations are broadly consistent with prior findings that thermal nonequilibrium has limited effect relative to perfect-gas assumptions~\citep{BitterShepherd2015}; however, the present results reveal that thermal equilibrium produces substantial amplification, an effect not previously emphasised in the literature. The hierarchy across conditions reflects the underlying Damköhler number dependence: Mach~6 configurations achieve rapid vibrational equilibration and thus approach \gls{TEQ} behaviour, whilst the Mach~8 hot-wall case remains closer to frozen-gas conditions due to insufficient relaxation time~\citep{KlineEtAl2018}. Cooled walls further promote equilibration by shortening the boundary layer residence time relative to the vibrational relaxation timescale, enhancing the transition towards \gls{TEQ} dynamics and explaining the higher peak $p_{\text{rms}}$ values observed in cold-wall configurations. These observations confirm that thermal equilibration enhances disturbance growth, although this advantage diminishes with increasing Mach number and wall temperature as the flow moves away from equilibrium conditions.

\begin{figure}[t]
    \centering
    \includegraphics[width=13cm]{resources/figures/disturbance/result_mach_pressure_rms.pdf}
    \caption{Streamwise evolution of pressure root-mean-square ($p_{\text{rms}}$) across the cases. Panels show: (a) Mach~6 cold-wall ($T_w=1200$~K), (b) Mach~6 hot-wall ($T_w=1800$~K), (c) Mach~8 cold-wall ($T_w=1200$~K), (d) Mach~8 hot-wall ($T_w=1800$~K).}
    \label{figure:prms_allcases}
\end{figure}


The N-factor comparison for all cases is showing in~\cref{figure:Nfactors_allcases}. Here one frequency has been selected in each case. The top row shows the $M_e=6$, $T_w=1200$ case, the second row $M_e=6$, $T_w=1200$, the third row $M_e=6$, $T_w=1200$ and the bottom row $M_e=8$, $T_w=1800$. The three columns show the CPG, TEQ and TNE cases respectively. It can be seen that the peak $N$-factors are higher for the cooled wall cases and lower for the higher Mach number cases and range from 6.5 in the $M_e=6$, $T_w=1200$ TEQ case down to 2.7 in the same case at $M_e=8$, $T_w=1800$. The $y$-shift correction applied for $M_e=8$ cases was considerably higher. The agreement between LST and DNS is still good, for the $M_e=8$ cases, despite the lower growth rates and greater non-parallelism of the base flows for these cases.

\begin{figure}[t]
    \centering
    \includegraphics[width=13cm]{resources/figures/disturbance/result_lst_dns_mach_Nfactor.pdf}
    \caption{Comparison of $N$-factors versus distance from the wedge leading edge across all configurations. Each row represents: (a) $M_e=6$, $T_w=1200$~K, (b) $M_e=6$, $T_w=1800$~K, (c) $M_e=8$, $T_w=1200$~K, (d) $M_e=8$, $T_w=1800$~K.}
    \label{figure:Nfactors_allcases}
\end{figure}

% Spatial integration of the growth rate yields N-factors that are compared against \gls{LST} predictions using a vertical amplitude shift to account for receptivity effects. The rate of N-factor increase in \gls{CPG} cases aligns well with \gls{LST} predictions, particularly when compared to lower Mach number cases, suggesting good agreement in the growth trend between DNS and \gls{LST} for the frozen-gas model. However, peak N-factors in \gls{CPG} cases deviate from \gls{LST} predictions, highlighting inconsistencies in their final amplitudes. Furthermore, the \gls{TEQ} model yields DNS N-factors that are consistently lower than \gls{LST} predictions, with this discrepancy most pronounced at lower frequencies in the hot-wall run, reaching differences of $\Delta N \approx 0.5$. This underprediction at larger streamwise distances is attributed to growing boundary layer non-parallelism, which violates the parallel-flow assumption inherent to \gls{LST}. The equilibrium model remains the most unstable of the simulations, and despite the underprediction at lower frequencies, still yields higher amplitude responses across most cases. Nonetheless, \gls{CPG} cases exhibit lower N-factors compared to \gls{TEQ} cases at both wall temperatures, consistent with prior studies. Notably, the amplitude differences between \gls{CPG}, \gls{TEQ}, and \gls{TNE} reveal that \gls{TNE} and \gls{CPG} responses are remarkably similar, particularly in hot-wall conditions where they yield virtually identical results.


Spatial integration of the growth rate yields N-factors that are compared against \gls{LST} predictions using a vertical amplitude shift to account for receptivity effects. The rate of N-factor increase in \gls{CPG} cases aligns well with \gls{LST} predictions, particularly when compared to lower Mach number cases, suggesting good agreement in the growth trend between DNS and \gls{LST} for the frozen-gas model. However, peak N-factors in \gls{CPG} cases deviate from \gls{LST} predictions, highlighting inconsistencies in their final amplitudes. Furthermore, the \gls{TEQ} model yields DNS N-factors that are consistently lower than \gls{LST} predictions, with this discrepancy most pronounced at lower frequencies in the hot-wall run, reaching differences of $\Delta N \approx 0.5$. This underprediction at larger streamwise distances is attributed to growing boundary layer non-parallelism, which violates the parallel-flow assumption inherent to \gls{LST}. 

\begin{figure}[t]
    \centering
    \includegraphics[width=13cm]{resources/figures/disturbance/result_all_lst_dns_envelope_comparison.pdf}
    \caption{N-factor envelopes for all Mach and wall temperature configurations. Rows: (a) $M_e=6$, $T_w=1200$~K, (b) $M_e=6$, $T_w=1800$~K, (c) $M_e=8$, $T_w=1200$~K, (d) $M_e=8$, $T_w=1800$~K. Columns show frequencies $\tilde{f}=0.13, 0.14, 0.15$ for Mach~6 cases and $\tilde{f}=0.11, 0.12, 0.13$ for Mach~8 cases.}
    \label{figure:lstdns_nfactor}
\end{figure}

The equilibrium model remains the most unstable of the simulations, and despite the underprediction at lower frequencies, still yields higher amplitude responses across most cases. Nonetheless, \gls{CPG} cases exhibit lower N-factors compared to \gls{TEQ} cases at both wall temperatures, consistent with prior studies. Notably, the amplitude differences between \gls{CPG}, \gls{TEQ}, and \gls{TNE} reveal that \gls{TNE} and \gls{CPG} responses are remarkably similar, particularly in hot-wall conditions where they yield virtually identical results. This convergence between frozen-gas and nonequilibrium models is significant: for the conditions considered in this study, the thermally frozen-gas model provides a robust \gls{LST} approximation for predicting nonequilibrium DNS growth behaviour, suggesting that detailed thermochemical modelling within \gls{LST} may be unnecessary when \gls{CPG} assumptions sufficiently capture the base-flow effects on instability dynamics. The comparison across all configurations demonstrates that whilst \gls{TEQ} sustains an amplification advantage over both \gls{CPG} and \gls{TNE}, the distinction between frozen and nonequilibrium responses diminishes significantly at higher Mach numbers and hot-wall conditions, reinforcing the Damköhler-based hierarchy established in the preceding analysis.

% \paragraph{Effect of vibrational nonequilibrium}

% For flows in thermal nonequilibrium (\gls{TNE}), it is important to establish the vibrational nonequilibrium signature across the flow conditions investigated to understand its influence on second-mode growth. 

% For acoustic disturbance analysis in thermochemically nonequilibrium flows, the vibrational and translational energy conservation equations are linearised about a base flow state. Decomposing quantities as $\phi = \bar{\phi} + \tilde{\phi}$ (base flow plus perturbation), the vibrational energy disturbance equation becomes

% \begin{equation}\label{equation:vibrational_disturbance}
% \frac{D\tilde{e}_v'}{Dt} = -\bar{u}_i \frac{\partial \tilde{e}_{v,0}}{\partial x_i} + (e_{v,0}^* - \bar{e}_v) \tau_{s,0}^{-1} + (e_{v,0}^* - \bar{e}_v) \tau_{s,0}^{-2} \left( \frac{\partial \tau_s}{\partial T} \tilde{T} + \frac{\partial \tau_s}{\partial p} \tilde{p} \right),
% \end{equation}

% and the translational energy disturbance equation is

% \begin{equation}\label{equation:translational_disturbance}
% \begin{split}
% \frac{D\tilde{e}_t'}{Dt} = -\left(e_{v,0}^* - \bar{e}_v\right) \tau_{s,0}^{-1} - (e_{v,0}^* - \bar{e}_v) \tau_{s,0}^{-2} \left( \frac{\rho'}{\rho_0} + \frac{\partial \tau_s}{\partial T} \tilde{T} + \frac{\partial \tau_s}{\partial p} \tilde{p} \right) \\
% - \left( \frac{\rho'}{\rho_0} \bar{U}_{i,0} + \bar{u}_i' \right) \frac{\partial \bar{e}_0}{\partial x_i} - \frac{P_0 \partial \bar{u}_i'}{\rho_0 \partial x_i} - \frac{p'}{
% \rho_0} \frac{\partial \bar{U}_{i,0}}{\partial x_i}.
% \end{split}
% \end{equation}

% In these equations, the three dominant production mechanisms are: (i) the mean-flow gradient production term $-\bar{u}_i (\partial \tilde{e}_{v,0}/\partial x_i)$, which arises from disturbance advection and couples to the streamwise variation of the base-flow state; (ii) the relaxation production term $(e_{v,0}^* - \bar{e}_v) \tau_{s,0}^{-1}$, which represents the energy exchange between translational and vibrational modes via the Landau-Teller formulation and introduces a phase lag between the two disturbance fields; and (iii) the thermodynamic production term proportional to the mean-flow nonequilibrium $(e_{v,0}^* - \bar{e}_v)$ and the relaxation timescale derivatives $(\partial \tau_s / \partial T)$ and $(\partial \tau_s / \partial p)$, which governs how disturbances interact with the background thermal nonequilibrium state.
\paragraph{Effect of vibrational nonequilibrium}

For flows in thermal nonequilibrium (\gls{TNE}), it is important to establish the vibrational nonequilibrium signature across the flow conditions investigated to understand its influence on second-mode growth. For acoustic disturbance analysis in thermochemically nonequilibrium flows, the vibrational and translational energy conservation equations are linearised about a base flow state. Following the framework of \citet{WagnildEtAl2012}, we decompose quantities as  \begin{equation}\label{equation:vibrational_disturbance}
\frac{d\tilde{e}_v'}{dt} = -\bar{u}_i \frac{\partial \tilde{e}_{v}}{\partial x_i} + (e_{v,0}^* - \bar{e}_v) \cfrac{1}{\bar{\tau}_{s}} + (e_{v,0}^* - \bar{e}_v) \cfrac{1}{\bar{\tau}_{s}^2} \left( \frac{\partial \bar{\tau}_s}{\partial T} \tilde{T} + \frac{\partial \bar{\tau}_s}{\partial p} \tilde{p} \right),
\end{equation} and the linearised total energy disturbance equation: \begin{equation}\label{equation:translational_disturbance}
\begin{split}
\frac{d\tilde{e}_t'}{dt} = -\left(e_{v,0}^* - \bar{e}_v\right) \cfrac{1}{\bar{\tau}_{s}} - (e_{v,0}^* - \bar{e}_v) \cfrac{1}{\bar{\tau}_{s}^2} \left( \frac{\tilde{\rho}}{\bar{\rho}} + \frac{\partial \bar{\tau}_s}{\partial T} \tilde{T} + \frac{\partial \bar{\tau}_s}{\partial p} \tilde{p} \right) \\
- \left( \frac{\tilde{\rho}}{\bar{\rho}} \bar{u}_{i} + \tilde{u}_i \right) \frac{\partial \bar{e}_0}{\partial x_i} - \cfrac{\bar{p}}{\bar{\rho}} \frac{\partial \tilde{u}_i}{\partial x_i} - \cfrac{\tilde{p}}{\bar{\rho}} \frac{\partial \bar{u}_{i}}{\partial x_i}.
\end{split}
\end{equation}
The three dominant production mechanisms in these equations are: (i) the mean-flow gradient production term $-\bar{u}_i (\partial \tilde{e}_{v}/\partial x_i)$, which represents disturbance advection by the base-flow velocity; (ii) the relaxation production term $(e_{v,0}^* - \bar{e}_v)/\bar{\tau}_s$, which models energy exchange between vibrational and translational modes via the Landau-Teller mechanism and introduces phase lag between the two disturbance fields; and (iii) the thermodynamic production term proportional to $(\partial \bar{\tau}_s/\partial T) \tilde{T}$ and $(\partial \bar{\tau}_s/\partial p) \tilde{p}$, which governs how the relaxation timescale responds to temperature and pressure perturbations, thereby modulating energy exchange efficiency during acoustic oscillations.

The linearised energy equations reveal a critical distinction between thermal equilibrium and nonequilibrium responses: the phase relationship between vibrational temperature disturbance $\tilde{T}_v$ and translational temperature disturbance $\tilde{T}$. In thermal equilibrium (\gls{TEQ}), where the relaxation timescale is very short, the vibrational mode remains phase-locked with the translational oscillations. The coupling term $(e_{v,0}^* - \bar{e}_v)/\bar{\tau}_s$ in Equation~\eqref{equation:vibrational_disturbance} responds synchronously to pressure fluctuations, amplifying the acoustic response and strengthening acoustic impedance contrast within the boundary layer acoustic cavity.

\begin{figure}[t]
    \centering
    \includegraphics[width=13cm]{resources/figures/disturbance/result_mackmodes_mach_Tv_perturbation.pdf}
    \caption{Contours of vibrational temperature disturbance $\hat{T}_v$ in a narrow streamwise window  for (a) $M_e=6$, $T_w=1200$, (b)$M_e=6$, $T_w=1200$, (c) $M_e=6$, $T_w=1200$, (d) $M_e=8$, $T_w=1800$).}
    \label{figure:mackmodes_tv}
\end{figure}


% In thermal nonequilibrium (\gls{TNE}), the finite relaxation timescale creates a phase lag: the vibrational energy cannot fully respond to rapid acoustic oscillations. For slow acoustic waves (low-frequency disturbances), where $\partial \bar{u}_i / \partial p < 0$, the mean-flow gradient production term generates positive energy production in the vibrational equation. However, the relaxation lag term $(e_{v,0}^* - \bar{e}_v)/\bar{\tau}_s^2 (\partial \bar{\tau}_s / \partial T) \tilde{T}$ introduces an out-of-phase component: when translational temperature rises, the vibrational response lags, creating a temporal offset between $\tilde{T}_v$ and $\tilde{T}$. This phase offset acts as an effective dissipative mechanism, decoupling the two energy reservoirs and reducing the cumulative acoustic amplification. The consequence, observed in the DNS results, is reduced spatial disturbance growth and downstream-shifted amplitude peaks in \gls{TNE} compared to \gls{TEQ}.

In thermal nonequilibrium (\gls{TNE}), the finite relaxation timescale creates a phase lag: the vibrational energy cannot fully respond to rapid acoustic oscillations. For slow acoustic waves (low-frequency disturbances), where $\partial \bar{u}_i / \partial p < 0$, the mean-flow gradient production term generates positive energy production in the vibrational equation. However, the relaxation lag term $(e_{v,0}^* - \bar{e}_v)/\bar{\tau}_s^2 (\partial \bar{\tau}_s / \partial T) \tilde{T}$ introduces an out-of-phase component: when translational temperature rises, the vibrational response lags, creating a temporal offset between $\tilde{T}_v$ and $\tilde{T}$. This phase offset acts as an effective dissipative mechanism, decoupling the two energy reservoirs and reducing the cumulative acoustic amplification. The consequence, observed in the DNS results, is reduced spatial disturbance growth and downstream-shifted amplitude peaks in \gls{TNE} compared to \gls{TEQ}.

\begin{figure}[t]
    \centering
    \includegraphics[width=13cm]{resources/figures/disturbance/result_T_Tv_perturbation.pdf}
    \caption{Comparison of $N$ factors versus the distance from the wedge leading edge. Each row represents the configuration (first row $M_e=6$, $T_w=1200$, second row $M_e=6$, $T_w=1200$, third row $M_e=6$, $T_w=1200$, bottom row $M_e=8$, $T_w=1800$). The frequency is $f_1 = 113.17 \ \mathrm{kHz}$ for the $M_e=6$ cases and $f_2 = 100.96 \mathrm{kHz}$ for the $M_e=8$ cases.}
    \label{figure:disturbance_tv}
\end{figure}


\Cref{figure:mackmodes_tv} presents an instantaneous snapshot of the vibrational temperature disturbance $\tilde{T}_v = T_v^{\text{forced}} - T_v^{\text{base}}$ extracted at the boundary layer edge ($y/\delta_{99} \approx 1$) across all four configurations. The contours reveal a striking out-of-phase behaviour: whilst disturbance quantities such as velocity $\tilde{u}_i$, $\tilde{v}_i$ and pressure $\tilde{p}$ oscillate in phase with the translational temperature (exhibiting the characteristic phase-tilted wave structures expected from acoustic cavity resonance), the vibrational temperature disturbance manifests a substantially delayed oscillation. The lag is most pronounced in the Mach~8 hot-wall case (panel d), where $\tilde{T}_v$ appears nearly 90 degrees out of phase with $\tilde{T}$. In the colder Mach~6 conditions (panels a and c), the lag is less extreme but clearly visible as a downstream shift in the phase progression. This spatial manifestation of the temporal lag confirms the theoretical prediction: the vibrational mode cannot maintain synchrony with the fast translational oscillations driven by acoustic pressure perturbations. The energy available in the acoustic wave is continuously being redirected into vibrational degrees of freedom at a rate determined by the relaxation timescale, and this partitioning occurs with systematic time delay, effectively removing energy from the acoustic growth cycle and suppressing second-mode amplification in thermal nonequilibrium flows.

The similarity between \gls{TNE} and \gls{CPG} N-factors reflects a fundamental decoupling mechanism: phase-lagged vibrational response renders the vibrational mode acoustically inert. In \gls{CPG}, vibrational excitation is frozen entirely. In \gls{TNE}, vibrational modes respond with substantial phase lag, extracting acoustic energy out-of-sync with the disturbance cycle. Both configurations therefore experience identical effective acoustic impedance controlled by translational-rotational modes alone. The vibrational mode in \gls{TNE}, despite theoretical excitation, contributes negligibly to acoustic dynamics due to desynchronisation, explaining why growth rates remain close to the frozen-gas limit across all flow conditions.

\section{Chapter summary}

Direct numerical simulations and linear stability theory have been carried out for a combination of Mach numbers (Mach 6 and Mach 8), wall temperatures (cold-wall $T_w = 1200$~K and hot-wall $T_w = 1800$~K), and thermal models (calorically perfect gas, thermal equilibrium, and thermal nonequilibrium) relevant to hypersonic flow over wedge-shaped geometries. Boundary layer similarity solutions provide inflow conditions for the simulations and base-flow profiles for linear stability analysis. A key finding is that thermally frozen base flow provides more physically consistent inflow conditions, since the relaxation toward thermal equilibrium coincides temporally with the excitation of second-mode instabilities, rendering the choice of thermal model critical to disturbance dynamics.

The simulations successfully capture Mack's second-mode acoustic instability across all configurations, with growth rates and spatial amplification profiles showing excellent agreement with linear stability theory predictions across the most-amplified frequency bands. When normalised for receptivity effects, N-factor comparisons between DNS and \gls{LST} remain robust across all thermochemical regimes, validating the parallel-flow assumption up to saturation. Critically, the results demonstrate a consistent hierarchy: thermal equilibrium cases are substantially more unstable than thermal nonequilibrium configurations, which in turn exhibit growth rates comparable to calorically perfect gas predictions. This hierarchy reflects underlying changes in acoustic impedance structure and wave trapping efficiency within the boundary layer acoustic cavity.

Analysis across Mach numbers and wall temperatures reveals competing timescale effects governed by the vibrational Damköhler number. At Mach 6 with cooled walls, vibrational equilibration occurs rapidly, and thermal equilibrium conditions dominate instability characteristics. At Mach 8 with heated walls, the boundary layer becomes thicker and translational temperatures rise, extending the vibrational relaxation timescale significantly. The Mach 8 cold-wall case exhibits faster equilibration than the Mach 6 hot-wall case, demonstrating that wall cooling promotes equilibration by shortening boundary layer residence time relative to relaxation timescale. These findings establish the physical basis for the observed amplitude hierarchy: configurations approaching equilibrium exhibit larger spatial disturbances and higher N-factors, whilst configurations remaining farther from equilibrium show suppressed growth.

The phase-speed and growth-rate analysis reveals that thermal nonequilibrium introduces a systematic downstream phase delay relative to thermal equilibrium, with this delay independent of wall temperature. Both frozen-gas (\gls{CPG}) and nonequilibrium (\gls{TNE}) models sustain higher phase speeds than thermal equilibrium, reflecting altered acoustic impedance arising from reduced or phase-lagged vibrational excitation. Critically, supersonic mode signatures (where phase speed dips below the $1 - 1/M_e$ line) are most pronounced in thermal equilibrium cases, indicating that thermal equilibration enhances the excitation of higher-order acoustic components and extends the effective instability mechanism beyond the primary second-mode band.

The thermal conductivity partition between translational-rotational ($\kappa_{tr}$) and vibrational ($\kappa_v$) modes emerges as a previously unidentified mechanism governing acoustic impedance structure. In thermal equilibrium, the vibrational thermal conductivity $\kappa_v$ actively extracts energy from boundary layer disturbances, suppressing peak temperature oscillations and reducing local temperature sensitivity. This energy redistribution weakens acoustic impedance contrast compared to frozen-gas assumptions. In thermal nonequilibrium, insufficient vibrational equilibration prevents $\kappa_v$ from fully withdrawing this energy on acoustic timescales, leaving translational-rotational temperature responses amplified. However, this enhancement is offset by a dissipative mechanism arising from the phase lag between vibrational and translational energy responses.

% Linearisation of the vibrational and translational energy conservation equations reveals three dominant production mechanisms governing disturbance evolution: (i) mean-flow gradient advection, (ii) the Landau-Teller relaxation term $(e_{v,0}^* - \bar{e}_v)/\bar{\tau}_s$, and (iii) thermodynamic production proportional to relaxation timescale derivatives $(\partial \bar{\tau}_s/\partial T)$ and $(\partial \bar{\tau}_s/\partial p)$. The relaxation term introduces a critical phase lag between vibrational and translational temperature disturbances, arising from the finite timescale required for energy transfer between molecular translational and vibrational degrees of freedom. During rapid acoustic oscillations, the vibrational temperature $\tilde{T}_v$ lags behind translational temperature $\tilde{T}$ by a phase angle $\phi = \arctan(\omega \bar{\tau}_s)$, where $\omega$ is the acoustic frequency. This phase lag is most pronounced at higher Mach numbers and elevated wall temperatures, where the relaxation timescale extends and the Damköhler number moves away from equilibrium.

% DNS contours of vibrational temperature disturbance at the boundary layer edge provide direct evidence of this phase lag: whilst pressure, velocity, and translational temperature disturbances oscillate in phase with characteristic acoustic phase-tilt patterns, the vibrational temperature manifests a substantially delayed oscillation. In the Mach 8 hot-wall case, $\tilde{T}_v$ appears nearly 90 degrees out of phase with $\tilde{T}$, indicating severe decoupling between energy reservoirs during acoustic cycles. This out-of-phase response acts as an effective energy sink: work performed by acoustic compression is partitioned into vibrational degrees of freedom asynchronously, preventing energy return to the acoustic field in phase to reinforce wave growth. The result is effective acoustic damping that suppresses cumulative second-mode amplification.

The convergence of thermal nonequilibrium and frozen-gas N-factors across all conditions is now explained through this decoupling mechanism. Phase-lagged vibrational response renders the vibrational mode acoustically inert despite formal inclusion in the model, functionally equivalent to complete vibrational freezing. In \gls{CPG}, no vibrational energy is excited. In \gls{TNE}, vibrational energy is excited but transferred asynchronously with acoustic oscillations, producing identical acoustic impedance structure. Only in thermal equilibrium does phase-locked vibrational coupling enhance acoustic impedance contrast and amplify second-mode growth. This distinction underscores that thermal nonequilibrium effects are fundamentally determined not merely by the degree of vibrational excitation, but by the synchronisation of energy exchange with acoustic timescales—a finding with significant implications for predicting transition in hypersonic nonequilibrium flows.




% Direct numerical simulation reveals compelling evidence of the vibrational temperature lag predicted by linear theory. \Cref{figure:Tv_disturbance_contours} presents contours of vibrational temperature disturbance $\tilde{T}_v$ across a representative streamwise window for the Mach~6 cases at four distinct configurations. The contour plots display characteristic acoustic wave structures: alternating regions of elevated (blue) and depressed (red) vibrational temperature oscillations, with wavelengths consistent with second-mode frequency content. Critically, the spatial organisation and phase progression of these contours differ subtly but systematically from the translational temperature fields shown previously. In the cold-wall cases (panels a and c), the vibrational disturbance penetration into the boundary layer is noticeably reduced compared to translational temperature, and the phase lines appear shifted downstream relative to the translational response. In the hot-wall cases (panels b and d), the vibrational disturbance amplitude diminishes markedly, and the spatial structures become increasingly diffuse.

% To isolate the phase relationship quantitatively, time series of vibrational and translational temperature disturbances are extracted at fixed spatial locations across the boundary layer depth. \Cref{figure:Tv_timeseries} displays these time traces for representative streamwise stations: the upper panels show translational-rotational temperature oscillations $\tilde{T}$, whilst the lower panels show the corresponding vibrational disturbances $\tilde{T}_v$. A striking observation emerges: at every spatial location and in every configuration examined, $\tilde{T}_v$ oscillates with systematically reduced amplitude compared to $\tilde{T}$, and crucially, the peaks and troughs of $\tilde{T}_v$ lag temporally behind those of $\tilde{T}$. This phase lag is most pronounced in the Mach~8 hot-wall case (right panels), where the vibrational oscillations appear nearly 90 degrees out of phase with the translational signal, indicating that when translational temperature reaches its maximum, vibrational temperature is still ascending. This temporal offset is the direct manifestation of the finite relaxation timescale: energy is being transferred from translational to vibrational reservoirs, but the transfer rate cannot keep pace with the rapid acoustic oscillations.

% The physical consequence of this lag is that vibrational and translational energies never achieve synchronised oscillation during a complete acoustic cycle. When translational temperature rises, vibrational energy begins absorbing energy, but by the time vibrational excitation is complete, translational temperature has already begun to fall. Similarly, when translational temperature drops, the vibrational mode continues to decay, arriving at its minimum after the translational field has already begun recovering. This asynchronous exchange acts as an effective energy sink: work done by acoustic compression is partitioned into vibrational degrees of freedom, but that energy is not returned to the acoustic field in phase to reinforce the oscillation. Instead, it dissipates through this out-of-phase coupling, directly suppressing the acoustic amplification observed in the growth-rate profiles and N-factor comparisons. The magnitude of this phase lag, quantified implicitly through the relaxation timescale derivatives $(\partial \bar{\tau}_s / \partial T)$ in the energy equations, increases with Mach number and wall temperature elevation, explaining the systematic reduction in disturbance amplitudes across the flow conditions investigated.

% \paragraph{Connection to second-mode (Mack mode) instability growth}

% The linearised energy equations reveal the physical pathways through which vibrational nonequilibrium modulates second-mode instability. In an idealised acoustic disturbance route, a planar acoustic wave enters the hypersonic boundary layer and, through thermoacoustic feedback, generates unstable second-mode oscillations that grow exponentially downstream. The three production mechanisms in Equations~\eqref{equation:vibrational_disturbance} and \eqref{equation:translational_disturbance} directly control this growth:

% \textit{Mean-flow advection} (term i) transports the energy disturbances downstream at the base-flow velocity $\bar{u}_i$. This mechanism determines the spatial location of the growth region and couples the disturbance evolution to the local base-flow structure, including its thermal gradients and acoustic cavity geometry.

% \textit{Relaxation coupling} (term ii), proportional to $(e_{v,0}^* - \bar{e}_v)/\bar{\tau}_s$, represents the rate at which vibrational and translational energies exchange. In thermal equilibrium (\gls{TEQ}), where $e_{v,0}^* \approx \bar{e}_v$ (small nonequilibrium departure), this term is weak. However, the exchange is \textit{phase-locked} with the acoustic oscillations, meaning vibrational energy absorbs and releases energy in synchrony with pressure–temperature fluctuations, enhancing acoustic impedance contrast and strengthening wave trapping within the acoustic cavity. This mechanism directly amplifies second-mode growth through improved acoustic resonance efficiency.

% In thermal nonequilibrium (\gls{TNE}), the finite relaxation timescale $\bar{\tau}_s$ prevents complete phase-locking. The vibrational response lags behind the translational oscillations, creating a phase offset. This decoupling introduces an effective dissipative mechanism: energy oscillates between the two modes without maintaining constructive interference, reducing cumulative amplification. Crucially, the sign convention of Equations~\eqref{equation:vibrational_disturbance} and \eqref{equation:translational_disturbance} shows that energy extracted by the vibrational term (negative sign in translational equation) is not fully recovered in the growth cycle, simulating radiative damping.

% \textit{Thermodynamic sensitivity} (term iii), through the derivatives $(\partial \bar{\tau}_s/\partial T)$ and $(\partial \bar{\tau}_s/\partial p)$, captures how the relaxation timescale itself varies with acoustic perturbations. When a pressure wave compresses the flow, both $T$ and $p$ increase, and the derivatives quantify whether the relaxation time shortens or lengthens. If $\partial \bar{\tau}_s / \partial T < 0$ (typical for high-temperature air), then a temperature rise accelerates relaxation, causing the vibrational mode to respond more quickly to the disturbance. This modifies the effective acoustic impedance in a nonlinear fashion: strong compression (large $\tilde{T}$ and $\tilde{p}$) forces rapid vibrational equilibration, temporarily altering wave trapping efficiency. This feedback loop either enhances or suppresses second-mode growth depending on the sign and magnitude of these derivatives and the degree of base-flow nonequilibrium.

% In the idealised disturbance route, second-mode growth emerges when these three mechanisms combine to produce a net positive production of acoustic energy. In \gls{TEQ}, all three terms cooperate: advection sustains propagation, phase-locked relaxation enhances impedance contrast, and thermodynamic feedback reinforces equilibration. In \gls{TNE}, the relaxation lag and thermodynamic sensitivity introduce competing effects: whilst advection and thermodynamic terms may still support growth, the phase-lagged relaxation damps the response, resulting in reduced spatial amplification and downstream-shifted growth peaks observed in the DNS results.


% \paragraph{Effect of vibrational nonequilibrium}

% For flows in \gls{TNE}, it is important to establish the vibrational nonequilibrium signature across the flow conditions investigated to understand its influence on second-mode growth. For acoustic disturbance analysis in vibrationally nonequilibrium flows, the vibrational and translational energy conservation equations are linearised about a base flow state. Following the framework of \citet{Wagnild}, we decompose quantities as $\phi = \bar{\phi} + \tilde{\phi}$ (base flow plus perturbation), yielding the vibrational energy disturbance equation

% \begin{equation}\label{equation:vibrational_disturbance}
% \frac{d\tilde{e}_v'}{dt} = -\bar{u}_i \frac{\partial \tilde{e}_{v,0}}{\partial x_i} + (e_{v,0}^* - \bar{e}_v) \cfrac{1}{\tau_{s,0}} + (e_{v,0}^* - \bar{e}_v) \cfrac{1}{\tau_{s,0}^2} \left( \frac{\partial \tau_s}{\partial T} \tilde{T} + \frac{\partial \tau_s}{\partial p} \tilde{p} \right),
% \end{equation}

% and the translational energy disturbance equation

% \begin{equation}\label{equation:translational_disturbance}
% \begin{split}
% \frac{d\tilde{e}_t'}{dt} = -\left(e_{v,0}^* - \bar{e}_v\right) \cfrac{1}{\tau_{s,0}} - (e_{v,0}^* - \bar{e}_v) \cfrac{1}{\tau_{s,0}^2} \left( \frac{\rho'}{\rho_0} + \frac{\partial \tau_s}{\partial T} \tilde{T} + \frac{\partial \tau_s}{\partial p} \tilde{p} \right) \\
% - \left( \frac{\rho'}{\rho_0} \bar{U}_{i,0} + \bar{u}_i' \right) \frac{\partial \bar{e}_0}{\partial x_i} - \frac{P_0}{\rho_0} \frac{\partial \bar{u}_i'}{\partial x_i} - \cfrac{p'}{\rho_0} \frac{\partial \bar{U}_{i,0}}{\partial x_i}.
% \end{split}
% \end{equation}




% For acoustic disturbance analysis, the vibrational and translational energy equations are linearized around a base flow state. Decomposing quantities as $\phi = \phi^0 + \phi'$ (base flow plus disturbance), the vibrational energy disturbance equation becomes

% \begin{equation}\label{equation:vibrational_disturbance}
% \begin{split}
% \frac{D\rho e_v'}{Dt} = -u_j^0 \frac{\partial \rho e_v'}{\partial x_j} - \rho_j^0 e_v^0 \frac{\partial u_j'}{\partial x_j} + \frac{\partial}{\partial x_j}\left(\kappa_v^0 \frac{\partial T_v'}{\partial x_j}\right) \\
% + \frac{\partial}{\partial x_j}\left(\rho^0 \sum_s h_{v,s}^0 D_s^0 \frac{\partial X_s'}{\partial x_j}\right) + \sum_s \rho_s^0 \frac{e_{v,\text{eq},s}' - e_{v,s}'}{\tau_s^0} + \sum_s \rho_s^0 e_{v,\text{eq},s}^0 \frac{\partial}{\partial \tau_s}\frac{\tau_s'}{\tau_s^0}.
% \end{split}
% \end{equation}

% Similarly, the translational energy disturbance equation is obtained by subtracting the base-flow total energy equation from the perturbed form. Since $E = e_t + e_v + \frac{1}{2}u_i u_i$, and separating translational and vibrational components:

% \begin{equation}\label{equation:translational_disturbance}
% \begin{split}
% \frac{D\rho e_t'}{Dt} = -u_j^0 \frac{\partial \rho e_t'}{\partial x_j} - \rho^0 e_t^0 \frac{\partial u_j'}{\partial x_j} + \frac{\partial}{\partial x_j}\left(\kappa_{tr}^0 \frac{\partial T'}{\partial x_j}\right) \\
% + \frac{\partial}{\partial x_j}\left(\rho^0 \sum_s h_{t,s}^0 D_s^0 \frac{\partial X_s'}{\partial x_j}\right) - \sum_s \rho_s^0 \frac{e_{v,\text{eq},s}' - e_{v,s}'}{\tau_s^0} + \frac{\partial u_i^0 \tau_{ij}'}{\partial x_j} + u_i' \frac{\partial \tau_{ij}^0}{\partial x_j}.
% \end{split}
% \end{equation}

% The three dominant production mechanisms in these equations are: (i) the mean-gradient production arising from velocity-pressure coupling in the base flow, (ii) the relaxation production term $\mathcal{Q}_{tv}' = \sum_s \rho_s^0 (e_{v,\text{eq},s}' - e_{v,s}')/\tau_s^0$, which couples the two energy modes through the Landau-Teller formulation and introduces phase lag between vibrational and translational disturbances, and (iii) the thermodynamic production proportional to the mean-flow nonequilibrium state $(T_v^0 - T^0)$, which modulates how disturbances interact with the background thermal nonequilibrium.





% In thermally nonequilibrium boundary layers, the vibrational mode does not necessarily track the translational--rotational temperature, introducing a distinct modification to the thermodynamic state. The TEQ model assumes $T_v = T$ throughout, while the TNE formulation accommodates finite-rate relaxation, allowing $T_v$ to lag behind $T$ depending on local flow conditions. This mismatch alters the effective heat-capacity ratio and the associated acoustic properties that underpin linear instability mechanisms. Although the resulting differences are generally moderate, they becom











% \begin{figure}[t]
%     \centering
%     \includegraphics[width=13cm]{resources/figures/disturbance/result_lst_dns_mach_Nfactor.pdf}
%     \caption{Comparison of $N$ factors versus the distance from the wedge leading edge. Each row represents the configuration (first row $M_e=6$, $T_w=1200$, second row $M_e=6$, $T_w=1200$, third row $M_e=6$, $T_w=1200$, bottom row $M_e=8$, $T_w=1800$). The frequency is $f_1 = 113.17 \ \mathrm{kHz}$ for the $M_e=6$ cases and $f_2 = 100.96 \mathrm{kHz}$ for the $M_e=8$ cases.}
%     \label{figure:Nfactors_allcases}
% \end{figure}


% The responses to the disturbances had an identical treatment to the Mach 6 cases with multiple eigenmodes chosen from an LST sweep of the $2.0 \ m $ long domain compared against the vertically adjusted DNS results. Figures~\ref{LST:m6tw18CPG} and~\ref{LST:m6tw18TEQ} show the growth in amplitude (here plotted as an N factor) of individual forced frequencies in DNS for the CPG and TEQ cases, compared with the respective LST prediction (dashed line) and y-shift added due to the receptivity coefficient. The rate of N-factor increase in CPG cases aligns well with LST predictions, especially compared to lower Mach number cases. This suggests good agreement in the growth trend between DNS and LST for CPG. However, the peak N-factors in CPG cases deviate from LST predictions, highlighting inconsistencies in their final amplitudes. Furthermore, unlike the low Mach number case, N-factors from DNS for TEQ cases are consistently lower than LST predictions. This discrepancy is most pronounced at lower frequencies in the hot wall run, reaching a difference of $N=0.5$. This suggests a more significant underprediction by DNS for TEQ, particularly at lower frequencies and hot wall conditions. The equilibrium proved to be the most unstable of the simulations, the underprediction at lower frequencies still gives a higher response in amplitudes except for $f=0.09$. Nonetheless, CPG cases exhibit lower N-factors for both hot and cold wall conditions compared to TEQ cases which is consistent in the previous studies. The amplitude differences in CPG, TEQ and TNEQ follow a similar trend where responses from TEQ are higher than CPG and TEQ with the latter two cases giving virtually identical results for hot wall conditions.

% \begin{figure}[t]
%     \centering
%     \includegraphics[width=13cm]{resources/figures/disturbance/result_all_lst_dns_envelope_comparison.pdf}
%     \caption{Amplitude of disturbances versus distance from wedge leading edge. Each row represents the configuration (\textit{rows}: (a) $M_e=6$, $T_w=1200$, (b) $M_e=6$, $T_w=1200$, (c) $M_e=6$, $T_w=1200$, (d) $M_e=8$, $T_w=1800$). Each column represents a different frequency, with frequency increasing from left to right, $\tilde f = \ 0.13,0.14,0.15$ for Mach 6 cases and $\tilde f = \ 0.11,0.12,0.13$ for Mach 8 cases.}
%     \label{figure:lstdns_nfactor}
% \end{figure}

% \subsubsection{Vibrational nonequilibrium}


% --- Stability maps ---
% \subsection{Stability maps}
% \noindent\textit{What to include.}
% Assemble $N_{\max}$ (or $N$ at a reference $x$) over the $3\times2\times2$ grid to produce contour maps vs.\ (Mach, $T_w$) for each thermal model, with iso-$N$ lines indicating likely transition risk. Add markers for synchronisation location and dominant frequency bands. Provide a brief lookup recipe (given Mach and $T_w$, choose model $\rightarrow$ read $N$, $f^\ast$, growth window) and flag regions where CPG/TEQ/TNE diverge most.

% --- Effects of vibrational nonequilibrium ---
% \subsection{Influence of Thermal exciations}
% \noindent\textit{What to include.}
% Isolate TNE effects by comparing TEQ vs.\ TNE at matched macroscopic base profiles (or by documenting any unavoidable differences). Show shifts in $\alpha_i$, $N$, phase speed, and modal energy distribution with/without vibrational relaxation. Discuss the roles of bulk viscosity and relaxation time scales in moving the acoustic trapping layer and altering damping, supported by frequency-dependent comparisons.


% Building on Kline’s cone study, we narrow the focus to \emph{vibrational nonequilibrium} alone and compare three clean limits: \emph{CPG} (frozen vibration), \emph{TEQ} (instantaneous vibration), and \emph{TNE} (finite-rate vibration). We deliberately decouple composition chemistry and wall catalysis so that any change in second- and “supersonic”-mode behaviour can be attributed to vibration rather than exothermic reactions. All cases are set up at matched, flight-relevant conditions (same reduced frequency band, comparable boundary-layer Reynolds measures, and common wall-to-edge temperature ratio) to make trends comparable across Mach number and wall temperature. Using DNS, we also document \emph{non-parallel} effects directly: the location and rate of growth shift with the streamwise development of the mean flow, explaining the systematic differences with local–parallel LST without introducing additional modelling. Taken together, this framing provides practical upper and lower bounds (TEQ vs.\ CPG), with TNE as the physically relevant bracket, and offers a transferable way to think about vibrational effects in high–Reynolds-number, flight-like conditions.





% --- Influence of thermal excitation ---
% \subsubsection{Influence of thermal excitation}
% \noindent\textit{What to include.}
% Quantify how thermal excitation in the \emph{base flow} (TEQ) vs.\ in the \emph{perturbations} (if applicable) modifies amplification. Report cases where holding the base fixed while toggling perturbation thermodynamics leaves growth nearly unchanged, and cases where it matters; connect to synchronisation shifts and wall-layer temperature gradients. Conclude with guidance on when a CPG stability solver suffices and when TEQ/TNE treatment is necessary.








% \section{Chapter summary}

% Direct numerical simulations and linear stability theory have been carried out for a combination of Mach numbers, wall temperatures and thermal models relevant to Mach 10 flow over wedge-shaped geometries. Boundary layer similarity solutions provide inflow conditions for the simulations and base flow for the linear theory. It is shown that a thermally-frozen base flow is more appropriate as an inflow condition for simulations since the relaxation toward thermal equilibrium coincides with the excitation of linear instabilities. The simulations capture Mack's second mode in thermally excited 2D simulations. The growth in amplitude obtained from the Fourier transform of DNS simulations compares well to the anticipated eigenmode growth calculated using LST for varying frequencies. The results show thermal equilibrium cases to be more unstable than cases with thermal nonequilibrium. Both of the thermally excited cases show higher eigenmode growth when compared to a calorically perfect gas simulation. Need to talk more about how thermal excitation varies across in one case, then go and talk about wall temperature variation. Following that, 

% Mach 6 cold wall case reached equilibrium quickly while Mach 8 hot wall took longer to reach equilibrium, however Mach 8 cold wall equilibrated faster than Mach 6 hot wall. Comparisons of \gls{LST} and \gls{DNS} were made for the flows, and the assumption of \gls{CPG} disturbance has yielded reasonable comparisons confirming the results of \cite{Wartemann2019} but has unpredicted results for \gls{TEQ}. The receptivity coefficients from \gls{DNS} are accounted for as \citep{Zhong2012} showed that DNS under predicted \gls{LST} results which has been true for \gls{TNE} and \gls{CPG} conditions but not for \gls{TEQ} showing that accurate modelling of disturbances is necessary for \gls{TEQ} conditions. 

% Across the three thermally excited models, the qualitative structure within the boundary layer is similar, although \gls{TEQ} yields the largest disturbance amplitudes and \gls{TNE} displays a slight downstream phase delay. In \gls{CPG} and \gls{TEQ} the waves decay more rapidly inside the layer, whereas in \gls{TNE} a small fraction of energy leaks into the freestream via partial transmission at the relative sonic line, producing weak dispersive signatures above the layer. This behaviour is compatible with a weakly radiating supersonic component; nevertheless, interactions within the broadband slow-acoustic spectrum must also be considered, and the observed leakage is too small to imply sustained radiation into the freestream~\citep{KniselyZhong2019}.


% \paragraph{Linearisation of the relaxation term}

% The relaxation term in the base-flow vibrational energy equation is:

% \begin{equation}
% \mathcal{Q}_{tv} = \sum_s \rho_s \frac{e_{v,\text{eq},s} - e_{v,s}}{\tau_s(T,p)}
% \end{equation}

% To linearise this term, we decompose around the base flow:

% \begin{equation}
% \mathcal{Q}_{tv} = \mathcal{Q}_{tv}^0 + \mathcal{Q}_{tv}' + O(\tilde{\phi}^2).
% \end{equation}

% The base-flow component is:

% \begin{equation}
% \mathcal{Q}_{tv}^0 = \sum_s \bar{\rho}_s \frac{\bar{e}_{v,\text{eq},s} - \bar{e}_{v,s}}{\bar{\tau}_s}.
% \end{equation}

% For the perturbation, we expand each term. The density perturbation is:

% \begin{equation}
% \rho_s + \tilde{\rho}_s = \bar{\rho}_s + \tilde{\rho}_s.
% \end{equation}

% The equilibrium vibrational energy varies with translational temperature. Expanding around $\bar{T}$:

% \begin{equation}
% e_{v,\text{eq},s}(T) = e_{v,\text{eq},s}(\bar{T}) + \frac{\partial e_{v,\text{eq},s}}{\partial T}\bigg|_{\bar{T}} \tilde{T} = \bar{e}_{v,\text{eq},s} + \frac{\partial \bar{e}_{v,\text{eq},s}}{\partial T} \tilde{T}.
% \end{equation}

% The actual vibrational energy perturbation is simply $\tilde{e}_{v,s}$.

% The relaxation time $\tau_s(T,p)$ depends on both temperature and pressure. Expanding around the base flow:

% \begin{equation}
% \tau_s(T,p) = \tau_s(\bar{T}, \bar{p}) + \frac{\partial \tau_s}{\partial T}\bigg|_{\bar{T},\bar{p}} \tilde{T} + \frac{\partial \tau_s}{\partial p}\bigg|_{\bar{T},\bar{p}} \tilde{p} = \bar{\tau}_s + \frac{\partial \bar{\tau}_s}{\partial T} \tilde{T} + \frac{\partial \bar{\tau}_s}{\partial p} \tilde{p}.
% \end{equation}

% Now, the reciprocal of the relaxation time is linearised using:

% \begin{equation}
% \frac{1}{\tau_s} = \frac{1}{\bar{\tau}_s + \frac{\partial \bar{\tau}_s}{\partial T} \tilde{T} + \frac{\partial \bar{\tau}_s}{\partial p} \tilde{p}} \approx \frac{1}{\bar{\tau}_s} \left(1 - \frac{1}{\bar{\tau}_s}\left(\frac{\partial \bar{\tau}_s}{\partial T} \tilde{T} + \frac{\partial \bar{\tau}_s}{\partial p} \tilde{p}\right)\right),
% \end{equation}

% which simplifies to:

% \begin{equation}
% \frac{1}{\tau_s} \approx \frac{1}{\bar{\tau}_s} - \frac{1}{\bar{\tau}_s^2}\left(\frac{\partial \bar{\tau}_s}{\partial T} \tilde{T} + \frac{\partial \bar{\tau}_s}{\partial p} \tilde{p}\right).
% \end{equation}

% Thus, the linear perturbation to the relaxation term becomes:

% \begin{equation}
% \mathcal{Q}_{tv}' = \sum_s \tilde{\rho}_s \frac{\bar{e}_{v,\text{eq},s} - \bar{e}_{v,s}}{\bar{\tau}_s} + \sum_s \bar{\rho}_s \frac{\partial \bar{e}_{v,\text{eq},s}}{\partial T} \tilde{T} \cfrac{1}{\bar{\tau}_s} - \sum_s \bar{\rho}_s (\bar{e}_{v,\text{eq},s} - \bar{e}_{v,s}) \cfrac{1}{\bar{\tau}_s^2}\left(\frac{\partial \bar{\tau}_s}{\partial T} \tilde{T} + \frac{\partial \bar{\tau}_s}{\partial p} \tilde{p}\right) \\
% - \sum_s \bar{\rho}_s \tilde{e}_{v,s} \cfrac{1}{\bar{\tau}_s}.
% \end{equation}

% Grouping by the nonequilibrium measure $(e_{v,\text{eq}}^* - e_v)$, and noting that this represents the departure from equilibrium in the base flow:

% \begin{equation}
% \mathcal{Q}_{tv}' = (e_{v,0}^* - \bar{e}_v) \cfrac{1}{\bar{\tau}_s} + (e_{v,0}^* - \bar{e}_v) \cfrac{1}{\bar{\tau}_s^2}\left(\frac{\partial \bar{\tau}_s}{\partial T} \tilde{T} + \frac{\partial \bar{\tau}_s}{\partial p} \tilde{p}\right).
% \end{equation}

% The **$1/\bar{\tau}_s^2$ term arises from the nonlinear expansion of the reciprocal relaxation time**, and the **$(\partial \bar{\tau}_s / \partial T)$ and $(\partial \bar{\tau}_s / \partial p)$ derivatives quantify how the relaxation timescale varies with thermodynamic variables**, thus modulating the sensitivity of the energy exchange to temperature and pressure perturbations.

% \subsection{Growth rates of disturbances}
% \subsubsection{Amplification independence}
% \subsubsection{Receptivity to Numerical Schemes} 

% \subsection{Scaling  of instability parameters}
% \subsection{Comparison across Mach numbers}
% \subsection{Stability maps}
% \subsection{Effects of vibrational nonequilibrium}
% \subsubsection{Influence of thermal excitation}
% \subsubsection{Comparison with thermal equilibrium cases}

% \subsection{Effect of vibrational nonequilibrium}
% \subsubsection{Comparison between Mach 6 and Mach 8}

% \subsection{Influence of wall temperature}


% \subsection{Summary of disturbance behaviour}

% \subsection{Linear Stability Theory}
% \subsection{Domain setup and configuration}
% \subsection{Results}
% \subsubsection{Base flows}
% \subsubsection{Growth rates}
% \subsubsection{N factors}





% The relaxation term in the base-flow vibrational energy equation is:
% \begin{equation}
% \mathcal{Q}_{tv} = \sum_s \rho_s \frac{e_{v,\text{eq},s} - e_{v,s}}{\tau_s(T,p)}.
% \end{equation}

% To linearise this term, the equilibrium vibrational energy is expanded around $\bar{T}$:
% \begin{equation}
% e_{v,\text{eq},s}(T) = \bar{e}_{v,\text{eq},s} + \frac{\partial \bar{e}_{v,\text{eq},s}}{\partial T} \tilde{T},
% \end{equation}

% and the relaxation time is expanded around the base flow:
% \begin{equation}
% \tau_s(T,p) = \bar{\tau}_s + \frac{\partial \bar{\tau}_s}{\partial T} \tilde{T} + \frac{\partial \bar{\tau}_s}{\partial p} \tilde{p}.
% \end{equation}

% The reciprocal of the relaxation time is linearised using $\frac{1}{\tau_s} \approx \frac{1}{\bar{\tau}_s} - \frac{1}{\bar{\tau}_s^2}\left(\frac{\partial \bar{\tau}_s}{\partial T} \tilde{T} + \frac{\partial \bar{\tau}_s}{\partial p} \tilde{p}\right)$, yielding the linearised vibrational energy disturbance equation:

